%==============================================================
%  CAPITOLO 8
%==============================================================
\chapter{Esempi di circuiti SRP/C}
\label{cap:08}
Questo report inizia trattando la progettazione di sistemi di controllo sicuri in termini generali. I paragrafi~\ref{sec:05-7}, 6.5 e 7.6 illustrano quindi, con riferimento all'esempio di una ghigliottina per il taglio della carta, come possono essere implementati i metodi per la progettazione di sistemi di controllo sicuri. I metodi per la determinazione del PL sono descritti passo-passo qui e nella norma EN ISO 13849-1; tuttavia, alcuni di questi passaggi, come la derivazione del diagramma a blocchi relativo alla sicurezza dallo schema elettrico, richiedono una certa pratica. Il manuale di SISTEMA~\cite{cookbook1} fornisce indicazioni su come derivare il diagramma a blocchi relativo alla sicurezza e il file SISTEMA dallo schema elettrico.

Tuttavia, a causa della varietà delle possibili funzioni di sicurezza e della loro implementazione, i singoli passaggi non si prestano a una descrizione generica. Per questo motivo, questo capitolo presenterà ora la valutazione di numerosi esempi di circuiti che implementano le funzioni di sicurezza in varie categorie e livelli di prestazione e mediante diverse tecnologie. Negli esempi di circuiti, il concetto di sistema di controllo generalmente comprende solo le parti dei sistemi di controllo relative alla sicurezza. Gli esempi sono limitati agli aspetti essenziali e servono quindi principalmente a illustrare la metodologia. Nella loro selezione è stata data importanza a un ampio spettro di tecnologie e possibili applicazioni. I lettori che hanno familiarità con il report del 1997~\cite{en954-1} sulle Categorie per i sistemi di controllo relativi alla sicurezza secondo la norma EN 954-1 riconosceranno alcuni degli esempi, ai quali è stato aggiunto, ad esempio, il calcolo della probabilità di guasto. Rispetto al Rapporto BGIA 2/2008e~\cite{bgiareport2-2008}, alcuni esempi non più aggiornati sono stati eliminati; è stato tuttavia aggiunto un nuovo esempio. Gli esempi rappresentano un'interpretazione delle Categorie e sono stati compilati dagli autori sulla base di molti anni di esperienza con i sistemi di controllo delle macchine relativi alla sicurezza e del lavoro svolto in comitati di normazione nazionali ed europei. Gli esempi servono a fornire ai progettisti una guida efficace per i loro sviluppi.

Poiché gli esempi sono stati creati da autori diversi, esiste inevitabilmente una certa variabilità, ad esempio nella presentazione dei dettagli o nel ragionamento alla base di alcuni dati numerici. Tutti i calcoli per gli esempi di circuito sono stati eseguiti con l'ausilio della versione 2.0 dell'applicazione software SISTEMA (vedere Allegato H), la versione disponibile al momento della produzione di questo rapporto. Ulteriori esempi di circuito, inclusi i file SISTEMA, sono descritti anche nel rapporto IFA 4/2018e, "Controlli di azionamento sicuri con inverter di frequenza"~\cite{werner2019}.

La descrizione di ciascun esempio è strutturata come segue:
\begin{itemize}
    \item funzione di sicurezza
    \item descrizione funzionale
    \item caratteristiche progettuali
    \item osservazioni
    \item calcolo della probabilità di guasto
    \item riferimenti più dettagliati
\end{itemize}

Alla voce "funzione di sicurezza" viene indicato il nome della funzione di sicurezza, unitamente agli eventi che la attivano e alle risposte di sicurezza richieste.

La "descrizione funzionale" descrive le funzioni essenziali legate alla sicurezza, sulla base di uno schema concettuale. Viene spiegato il comportamento in caso di guasto, e vengono indicate le misure per il rilevamento dei guasti.
Le caratteristiche particolari della progettazione dell'esempio in questione, quali l'applicazione di principi di sicurezza ben collaudati e l'impiego di componenti ben collaudati, sono elencate alla voce "caratteristiche progettuali".

Gli schemi circuitali sono schemi concettuali, limitati esclusivamente alla rappresentazione della/e funzioni di sicurezza con i componenti pertinenti necessari a questo particolare scopo. Nell'interesse della chiarezza, sono stati omessi alcuni circuiti aggiuntivi normalmente necessari, ad esempio quelli per garantire la protezione contro la scossa elettrica, per il controllo della sovratensione/sottotensione e della sovrapressione o bassa pressione, per il rilevamento di guasti di isolamento, cortocircuiti e guasti verso terra, ad esempio su linee instradate esternamente, o per garantire la richiesta immunità ai disturbi elettromagnetici. Dettagli circuitali non essenziali per la determinazione dello schema a blocchi legato alla sicurezza sono stati quindi deliberatamente omessi. Tali dettagli comprendono i circuiti di protezione nel sistema elettrico, quali fusibili e diodi, ad esempio sotto forma di diodi di ricircolo (free-wheeling diode). Gli schemi omettono inoltre i diodi di disaccoppiamento nei circuiti in cui, ad esempio, i segnali dei sensori vengono letti in modo ridondante da più unità logiche. Questa disposizione è finalizzata a impedire che, in caso di guasto, un ingresso diventi un'uscita nei sistemi ridondanti, influenzando così il secondo canale. Tutti questi componenti sono essenziali affinché un sistema di comando venga realizzato in conformità a una categoria e a un Performance Level. Conformemente agli elenchi dei guasti della EN ISO 13849-2, questioni quali l'influenza dei cortocircuiti dei conduttori devono naturalmente essere considerate anche in relazione alla funzione di sicurezza interessata e alle condizioni d'uso. Tutti i componenti impiegati devono pertanto essere selezionati tenendo conto della loro idoneità in base alla relativa specifica. Il sovradimensionamento è uno dei principi di sicurezza ben collaudati.

Ulteriori esempi sono elencati nelle osservazioni specifiche per tecnologia relative alla tecnica dei fluidi.

Le caratteristiche progettuali sono indicate soltanto ove siano rilevanti per le funzioni di sicurezza descritte. Si tratta generalmente di una "funzione di arresto legata alla sicurezza, avviata da un mezzo di protezione". Altre funzioni di sicurezza, quali la "prevenzione dell'avviamento inatteso" o una "funzione di ripristino manuale" e la "funzione di avviamento/riavviamento", non sono considerate in tutti gli esempi. Qualora vengano utilizzate apparecchiature ad azionamento manuale (pulsanti) per la realizzazione di tali funzioni di sicurezza, occorre garantire che, laddove la funzione di sicurezza sia realizzata in combinazione con l'elettronica in particolare, essa debba essere avviata dal rilascio (azione di apertura) di un pulsante già premuto.

Ove pertinente per l'esempio in questione, un riferimento particolare ad aspetti specifici di una possibile applicazione è riportato alla voce "Osservazioni".

Alla voce "Calcolo della probabilità di guasto" viene fornita una descrizione del calcolo del PL a partire dai parametri categoria, MTTF\textsubscript{D}, DC\textsubscript{avg} e CCF, sulla base dello schema a blocchi legato alla sicurezza derivato dallo schema concettuale. La categoria viene determinata a partire dalla descrizione funzionale e dalle caratteristiche progettuali.

I valori di MTTF\textsubscript{D} impiegati nei calcoli sono contrassegnati come valori del fabbricante ("[M]" per manufacturer), valori tipici da banche dati ("[D]" per database), oppure valori tratti dalla EN ISO 13849-1 ("[S]" per standard). Conformemente alla norma, la priorità dovrebbe essere data ai dati dei fabbricanti. Per alcuni componenti, al momento della redazione del rapporto non erano disponibili né dati affidabili dei fabbricanti né valori da banche dati. In questo caso è stato impiegato il metodo del conteggio dei componenti per la stima di valori esemplificativi tipici (contrassegnati "[E]" per estimated, stimato). I valori di MTTF\textsubscript{D} riportati in questo capitolo devono pertanto essere considerati, in alcuni casi, più come stime.

La presentazione delle misure ipotizzate per la diagnostica (DC) e contro i guasti per causa comune (CCF) è limitata a informazioni di carattere generale. I valori specifici per questi due criteri dipendono dalla realizzazione, dall'applicazione e dal fabbricante. È pertanto possibile che, per componenti simili, in esempi diversi vengano ipotizzati valori di DC differenti. Anche in questo caso, tutte le ipotesi relative a DC e CCF devono essere riesaminate nelle realizzazioni concrete; i valori ipotizzati non sono vincolanti e sono destinati esclusivamente a scopo illustrativo.

Nella presentazione, l'attenzione è rivolta principalmente alle categorie intese come "resistenza ai guasti", allo schema a blocchi e ai metodi "matematici" per la determinazione del PL. Al contrario, alcune fasi parziali, come l'esclusione di guasto, i principi di sicurezza di base e ben collaudati o le misure contro i guasti sistematici (compreso il software), sono menzionate solo brevemente. Durante la realizzazione, occorre prestare a questo aspetto un'attenzione adeguata, poiché valutazioni errate o una realizzazione inadeguata di tali misure potrebbero comportare un peggioramento della tolleranza ai guasti o della probabilità di guasto. Come ausilio per la comprensione degli esempi circuitali e per la loro realizzazione pratica, l'attenzione del lettore è pertanto richiamata sul Capitolo~\ref{cap:07} e sull'Allegato~\ref{app:C}, nei quali, ad esempio, i principi di sicurezza di base e ben collaudati sono descritti in dettaglio.

Infine, si rimanda ai "riferimenti più dettagliati", ove disponibili.

Per ciascun tipo di tecnologia, nei successivi punti specifici per tecnologia vengono formulati alcuni commenti di carattere generale, al fine di fornire una migliore comprensione degli esempi e per la realizzazione delle categorie.
Alcuni degli esempi circuitali rappresentano "sistemi di comando che coinvolgono tecnologie multiple". Questi esempi circuitali "misti" si basano sul concetto, sancito dalla norma, secondo cui una funzione di sicurezza viene sempre realizzata mediante "ricezione", "elaborazione" e "commutazione", indipendentemente dalla tecnologia impiegata."

%--------------------------------------------------------------
\section{Osservazioni generali di natura tecnologica sui sistemi di controllo di esempio}
\label{sec:08-1}

\subsection{Comandi elettromeccanici}
\label{subsec:08-1-1}
I comandi elettromeccanici impiegano principalmente componenti elettromeccanici sotto forma di dispositivi di controllo (ad esempio interruttori di posizione, selettori, pulsanti) e apparecchiature di commutazione (relè contattori, relè, contattori). Questi dispositivi hanno posizioni di commutazione definite. Generalmente non cambiano il loro stato di commutazione se non azionati esternamente o elettricamente. Se selezionati correttamente e utilizzati come previsto, sono in gran parte immuni a disturbi, come interferenze elettriche o elettromagnetiche. Sotto questo aspetto differiscono, in alcuni casi considerevolmente, dalle apparecchiature elettroniche. La loro durata e le modalità di guasto possono essere influenzate da un'adeguata selezione, dimensionamento e disposizione. Lo stesso vale per i conduttori impiegati, se opportunamente instradati all'interno e all'esterno dei vani elettrici.

Per le ragioni sopra indicate, i componenti elettromeccanici generalmente soddisfano i “principi di sicurezza di base” e in molti casi devono essere considerati anche “componenti collaudati” per applicazioni di sicurezza. Ciò vale, tuttavia, solo quando vengono rispettati i requisiti della norma IEC 60204-1~\cite{iec60204-1} per le apparecchiature elettriche della macchina/impianto. In alcuni casi, sono possibili esclusioni di guasto, ad esempio su un contattore di comando per quanto riguarda l'attivazione in assenza di tensione di comando, o la mancata apertura di un contatto di interruzione con azione di apertura diretta su un interruttore secondo la norma IEC 60947-5-1~\cite{iec60947-5-1}, Allegato K.

Informazioni dettagliate sulla modellazione dei componenti elettromeccanici sono reperibili nell'Allegato~\ref{app:D}.


%--------------------------------------------------------------
\subsection{Controlli idraulici}
\label{sec:08-1-2}
Negli impianti oleodinamici, la zona delle valvole, ovvero le valvole che controllano movimenti o stati pericolosi, deve essere considerata in particolare una "parte del sistema di controllo rilevante per la sicurezza". I circuiti oleodinamici elencati di seguito costituiscono solo esempi. Di norma, le funzioni di sicurezza richieste possono essere implementate anche mediante logiche di controllo alternative che impiegano tipologie di valvole appropriate, o in alcuni casi anche mediante soluzioni meccaniche aggiuntive come dispositivi di blocco o freni.

Nei sistemi idraulici (vedere Figura 8.1), in questo contesto vanno considerate anche le misure per la limitazione della pressione nel sistema (1V2) e per la filtrazione del fluido idraulico (1Z2).
I componenti 1Z1, 1S1 e 1S2 mostrati in Figura 8.1 sono presenti nella maggior parte dei sistemi idraulici e rivestono grande importanza, in particolare per le condizioni del fluido idraulico e di conseguenza per il funzionamento delle valvole.
Il filtro di sfiato del serbatoio 1Z1, disposto sul serbatoio del fluido, impedisce l'ingresso di impurità esterne. L'indicatore di livello del fluido 1S2 garantisce che il livello del fluido rimanga entro i limiti specificati. L'indicatore di temperatura 1S1 costituisce una misura idonea per la limitazione dell'intervallo di temperatura di esercizio e quindi dell'intervallo di viscosità di esercizio del fluido idraulico. Se necessario, devono essere previsti dispositivi di riscaldamento e/o raffreddamento in combinazione con il controllo della temperatura a circuito chiuso (vedere anche l'Allegato~\ref{app:C} a questo proposito).
Gli elementi di azionamento e i componenti per la conversione la trasmissione di energia nei sistemi oleodinamici generalmente non rientrano nell'ambito di applicazione della norma.

Nei sistemi pneumatici (vedere Figura 8.2), i componenti per la prevenzione dei rischi associati alla conversione di energia e l'unità di manutenzione per il condizionamento dell'aria compressa devono essere considerati, dal punto di vista della sicurezza, in combinazione con l'area delle valvole. Per controllare le possibili conversioni di energia tenendo conto degli aspetti di sicurezza, si utilizza spesso una valvola di scarico in combinazione con un pressostato. Negli esempi di circuito in questo capitolo, questi componenti sono indicati con 0V1 (valvola di scarico) e 0S1 (pressostato).
L'unità di manutenzione 0Z (vedere Figura 8.2) è generalmente composta da una valvola di intercettazione manuale 0V10, un filtro con separatore d'acqua 0Z10 con monitoraggio della contaminazione del filtro e una valvola di controllo della pressione 0V11 (con sfiato secondario adeguatamente dimensionato). L'indicatore di pressione 0Z11 soddisfa il requisito per il monitoraggio dei parametri del sistema.

% La minipage tiene insieme disegno e didascalia (figura non fluttuante)
\begin{center}
    \begin{minipage}{\textwidth}
        \centering
        \captionsetup{hypcap=false}
        \resizebox{0.92\textwidth}{!}{%==============================================================
%  Figura 8.1 - Campo di applicazione della EN ISO 13849 per i
%  sistemi idraulici
%  Adattamento in italiano della Figura 8.1 dell'IFA Report 2/2017
%
%  I simboli dei componenti riproducono quelli dell'originale
%  (ISO 1219): le coordinate sono ricavate dal disegno di partenza
%  (1 cm = 70 px, origine nell'angolo in alto a sinistra della
%  zona tratteggiata).
%
%  NOTA: nel disegno compare il simbolo matematico $\sim$: come
%  nelle Figure 3.2 e 3.3 non vanno usati cambi di corpo del font
%  (font=\large & c.), che farebbero fallire la composizione.
%==============================================================
\begin{tikzpicture}[
  font=\normalsize,
  linea/.style={draw=black, line width=0.9pt},
  simbolo/.style={draw=black, line width=0.9pt},
  freccia/.style={simbolo, -{Latex[length=5pt,width=4pt]}},
  zona/.style={draw=black, line width=0.9pt, dash pattern=on 7pt off 5pt},
  pilota/.style={draw=black, line width=0.7pt, dash pattern=on 3pt off 3pt},
  testo/.style={align=center, inner sep=1pt},
  nodo/.style={circle, fill=black, inner sep=0pt, minimum size=4pt}
]

%--------------------------------------------------------------
% Zone tratteggiate: in alto i componenti che realizzano la
% funzione di sicurezza, in basso conversione e trasmissione
% dell'energia
%--------------------------------------------------------------
\draw[zona] (0,0) rectangle (13.21,-12.50);
\draw[zona] (0,-4.64) -- (13.21,-4.64);
\draw[zona] (0,-7.71) -- (13.21,-7.71);

\node[testo, text width=5.6cm] at (6.71,-2.20)
  {Componenti che realizzano\\ la funzione di sicurezza,\\ ad es. valvole};
\node[testo, text width=4.4cm] at (10.60,-8.35)
  {Conversione di energia\\ Trasmissione di energia};

%--------------------------------------------------------------
% Cilindro 1A (elemento di azionamento): camicia, pistone e stelo
% a doppia linea, blocco di attacco con le due bocche
%--------------------------------------------------------------
\draw[simbolo] (5.73,1.36) rectangle (7.43,2.29);
\draw[simbolo] (5.90,1.36) -- (5.90,2.29);
\draw[simbolo] (6.26,1.36) -- (6.26,2.29);
\draw[simbolo, fill=white] (6.26,1.74) rectangle (7.97,1.90);
\draw[linea] (5.76,1.36) -- (5.76,0.67) -- (6.39,0.67) -- (6.39,0);
\draw[linea] (7.39,1.36) -- (7.39,0.67) -- (6.74,0.67) -- (6.74,0);

\node[anchor=east] at (5.55,2.11) {1A};
\node[anchor=east] at (5.55,1.36) {Elementi di azionamento};

%--------------------------------------------------------------
% Linee di collegamento
%--------------------------------------------------------------
% aspirazione dal serbatoio e mandata della pompa verso le valvole
\draw[linea] (4.03,-12.50) -- (4.03,-11.31);
\draw[linea] (4.03,-10.21) -- (4.03,-9.47);
\draw[linea] (4.03,-8.77) -- (4.03,-4.64);
% diramazione verso la valvola limitatrice 1V2 e scarico al serbatoio
\draw[linea] (4.03,-6.25) -- (6.48,-6.25) -- (6.48,-12.50);
\node[nodo] at (4.03,-6.25) {};
% ritorno attraverso il filtro 1Z2
\draw[linea] (8.23,-4.64) -- (8.23,-5.69);
\draw[linea] (8.23,-6.74) -- (8.23,-12.50);

%--------------------------------------------------------------
% 1V2: valvola limitatrice di pressione con molla regolabile e
% pilotaggio
%--------------------------------------------------------------
\draw[simbolo] (5.16,-5.87) rectangle (5.91,-6.61);
\draw[freccia] (5.16,-6.43) -- (5.88,-6.43);
% molla verticale sopra la valvola
\draw[simbolo] (5.38,-5.87) -- (5.56,-5.80) -- (5.16,-5.71) -- (5.56,-5.62)
              -- (5.16,-5.53) -- (5.32,-5.43);
% freccia di regolazione
\draw[freccia] (4.98,-5.73) -- (5.80,-5.53);
% pilotaggio
\draw[pilota] (5.16,-6.25) -- (4.79,-6.63) -- (4.79,-6.96) -- (5.34,-6.96)
             -- (5.34,-6.61);
\node[anchor=east] at (4.98,-5.96) {1V2};

%--------------------------------------------------------------
% 1Z2: filtro sulla linea di ritorno con valvola di by-pass e
% indicatore di intasamento
%--------------------------------------------------------------
\draw[simbolo] (6.90,-5.31) rectangle (8.97,-7.14);
\node[nodo] at (8.23,-5.46) {};
\node[nodo] at (8.23,-6.97) {};
% simbolo del filtro
\draw[simbolo] (8.23,-5.69) -- (8.76,-6.21) -- (8.23,-6.74) -- (7.71,-6.21) -- cycle;
\draw[pilota] (7.71,-6.21) -- (8.76,-6.21);
% by-pass con valvola di non ritorno e molla
\draw[linea] (8.23,-5.46) -- (7.29,-5.46) -- (7.29,-5.89);
\draw[simbolo] (7.18,-6.10) -- (7.29,-5.90) -- (7.40,-6.10);
\draw[simbolo] (7.29,-6.00) circle (0.10);
\draw[simbolo] (7.29,-6.11) -- (7.47,-6.17) -- (7.10,-6.24) -- (7.47,-6.31)
              -- (7.10,-6.38) -- (7.47,-6.45) -- (7.29,-6.52);
\draw[linea] (7.29,-6.52) -- (7.29,-6.97) -- (8.23,-6.97);
% indicatore di intasamento
\draw[linea] (8.23,-5.46) -- (9.36,-5.46);
\draw[simbolo] (9.36,-5.17) rectangle (10.29,-5.74);
\draw[freccia] (9.76,-5.62) -- (9.47,-5.33);
\draw[simbolo] (10.05,-5.46) circle (0.14);
\draw[simbolo] (9.95,-5.36) -- (10.15,-5.56);
\draw[simbolo] (9.95,-5.56) -- (10.15,-5.36);
\node[anchor=south west] at (6.93,-5.24) {1Z2};

%--------------------------------------------------------------
% 1V1: valvola di non ritorno con molla
%--------------------------------------------------------------
\draw[simbolo] (4.03,-8.77) -- (4.21,-8.83) -- (3.83,-8.91) -- (4.21,-8.99)
              -- (3.83,-9.07) -- (4.21,-9.14) -- (4.03,-9.20);
\draw[simbolo] (4.03,-9.33) circle (0.09);
\draw[simbolo] (3.90,-9.30) -- (4.03,-9.47) -- (4.16,-9.30);
\node[anchor=east] at (3.72,-8.95) {1V1};

%--------------------------------------------------------------
% 1M e 1P: motore elettrico, giunto e pompa
%--------------------------------------------------------------
\draw[simbolo] (1.60,-10.76) circle (0.55);
\node[testo] at (1.62,-10.55) {M};
\node[testo] at (1.57,-10.92) {3\,$\sim$};
\node[anchor=east] at (0.97,-10.76) {1M};

% albero a doppia linea con giunto
\draw[simbolo] (2.14,-10.67) -- (2.74,-10.67);
\draw[simbolo] (2.14,-10.86) -- (2.74,-10.86);
\draw[simbolo] (2.74,-10.57) -- (2.74,-10.97);
\draw[simbolo] (2.90,-10.57) -- (2.90,-10.97);
\draw[simbolo] (2.90,-10.67) -- (3.47,-10.67);
\draw[simbolo] (2.90,-10.86) -- (3.47,-10.86);

\draw[simbolo] (4.03,-10.76) circle (0.55);
\fill[black] (4.03,-10.23) -- (3.83,-10.53) -- (4.21,-10.53) -- cycle;
\draw[freccia] ([shift={(-32:0.86)}]4.03,-10.76) arc (-32:28:0.86);
\node[anchor=west] at (5.03,-10.76) {1P};

%--------------------------------------------------------------
% 1S1, 1S2 e 1Z1: accessori del serbatoio
%--------------------------------------------------------------
% termometro
\draw[simbolo] (9.96,-11.19) circle (0.40);
\draw[simbolo] (9.96,-10.90) -- (9.96,-11.40);
\fill[black] (9.96,-11.43) circle (0.06);
\draw[linea] (9.96,-11.59) -- (9.96,-12.07);
\node[anchor=south] at (9.96,-10.75) {1S1};
% indicatore di livello
\draw[simbolo] (11.07,-11.19) circle (0.40);
\draw[simbolo] (10.71,-11.21) -- (11.43,-11.21);
\fill[black] (10.96,-11.04) -- (11.19,-11.04) -- (11.07,-11.19) -- cycle;
\draw[linea] (11.07,-11.59) -- (11.07,-12.07);
\node[anchor=south] at (11.07,-10.75) {1S2};
% filtro di sfiato del serbatoio
\draw[simbolo] (12.40,-10.58) -- (12.92,-11.10) -- (12.40,-11.62)
              -- (11.88,-11.10) -- cycle;
\draw[pilota] (11.88,-11.10) -- (12.92,-11.10);
\draw[linea] (12.40,-11.62) -- (12.40,-12.07);
\draw[simbolo] (12.31,-10.40) -- (12.49,-10.40) -- (12.40,-10.57) -- cycle;
\draw[linea] (12.40,-10.40) -- (12.40,-10.19);
\draw[simbolo] (12.31,-10.19) -- (12.49,-10.19) -- (12.40,-10.03) -- cycle;
\node[anchor=east] at (12.24,-10.10) {1Z1};

%--------------------------------------------------------------
% Graffe e campi di applicazione
%--------------------------------------------------------------
\draw[linea, decorate, decoration={brace, amplitude=8pt}]
  (14.29,-4.50) -- (14.29,-0.36);
\node[testo, text width=5.2cm, anchor=west] at (14.95,-2.43)
  {Campo di applicazione della\\ parte del sistema di comando\\ legata alla sicurezza};

\draw[linea, decorate, decoration={brace, amplitude=8pt}]
  (14.29,-12.50) -- (14.29,-5.07);
\node[testo, text width=5.2cm, anchor=west] at (14.95,-8.79)
  {Possibile rilevanza per il\\ rispetto dei principi di\\ sicurezza di base e collaudati};

%--------------------------------------------------------------
% Cornice
%--------------------------------------------------------------
\draw[draw=black, line width=1pt]
  ([shift={(-0.45,0.45)}]current bounding box.north west) rectangle
  ([shift={(0.45,-0.45)}]current bounding box.south east);

\end{tikzpicture}
}
        \captionof{figure}{Campo di applicazione della EN ISO 13849 per i sistemi idraulici.}
        \label{fig:08-01}
    \end{minipage}
\end{center}

% La minipage tiene insieme disegno e didascalia (figura non fluttuante)
\begin{center}
    \begin{minipage}{\textwidth}
        \centering
        \captionsetup{hypcap=false}
        \resizebox{!}{0.62\textheight}{%==============================================================
%  Figura 8.2 - Campo di applicazione della EN ISO 13849 per i
%  sistemi pneumatici
%  Adattamento in italiano della Figura 8.2 dell'IFA Report 2/2017
%  (stesso stile della Figura 8.1)
%
%  I simboli dei componenti riproducono quelli dell'originale
%  (ISO 1219): le coordinate sono ricavate dal disegno di partenza
%  (1 cm = 70 px, origine nell'angolo in alto a sinistra della
%  zona tratteggiata). Nessun cambio di corpo del font: vedere la
%  nota nella Figura 8.1.
%==============================================================
\begin{tikzpicture}[
  font=\normalsize,
  linea/.style={draw=black, line width=0.9pt},
  simbolo/.style={draw=black, line width=0.9pt},
  freccia/.style={simbolo, -{Latex[length=5pt,width=4pt]}},
  doppia/.style={simbolo, {Latex[length=5pt,width=4pt]}-{Latex[length=5pt,width=4pt]}},
  zona/.style={draw=black, line width=0.9pt, dash pattern=on 7pt off 5pt},
  pilota/.style={draw=black, line width=0.7pt, dash pattern=on 3pt off 3pt},
  testo/.style={align=center, inner sep=1pt},
  nodo/.style={circle, fill=black, inner sep=0pt, minimum size=4pt}
]

%--------------------------------------------------------------
% Zone tratteggiate: componenti della funzione di sicurezza,
% componenti contro le fluttuazioni dell'energia, unita' di
% manutenzione
%--------------------------------------------------------------
\draw[zona] (0,0) rectangle (11.09,-16.71);
\draw[zona] (0,-5.89) -- (11.09,-5.89);
\draw[zona] (0,-9.60) -- (11.09,-9.60);
\draw[zona] (1.87,-10.86) rectangle (10.54,-14.83);

\node[testo, text width=5.6cm] at (5.51,-2.86)
  {Componenti che realizzano\\ la funzione di sicurezza,\\ ad es. valvole};
\node[testo, text width=5.6cm] at (5.51,-7.83)
  {Componenti per la prevenzione\\ dei rischi in caso di\\ fluttuazioni dell'energia};
\node[anchor=east] at (1.72,-10.98) {0Z};
\node[anchor=east] at (10.90,-16.00) {\textquotedblleft Unit\`a di manutenzione\textquotedblright};

%--------------------------------------------------------------
% Cilindro 1A (elemento di azionamento): camicia, pistone e stelo
% a doppia linea, blocco di attacco con le due bocche
%--------------------------------------------------------------
\draw[simbolo] (4.44,2.17) rectangle (6.57,3.33);
\draw[simbolo] (4.67,2.17) -- (4.67,3.33);
\draw[simbolo] (5.14,2.17) -- (5.14,3.33);
\draw[simbolo, fill=white] (5.14,2.61) rectangle (7.29,2.86);
\draw[linea] (4.50,2.17) -- (4.50,1.31) -- (5.29,1.31) -- (5.29,0);
\draw[linea] (6.50,2.17) -- (6.50,1.31) -- (5.71,1.31) -- (5.71,0);

\node[anchor=east] at (4.30,3.31) {1A};
\node[anchor=east] at (4.30,2.14) {Elementi di azionamento};

%--------------------------------------------------------------
% Linea di alimentazione pneumatica (interrotta dal filtro 0Z10 e
% dalla valvola 0V10)
%--------------------------------------------------------------
\draw[linea] (0.50,-9.60) -- (0.50,-13.43) -- (6.29,-13.43);
\draw[linea] (7.60,-13.43) -- (8.57,-13.43);
\draw[linea] (9.54,-13.43) -- (11.29,-13.43);
% sorgente di aria compressa
\draw[simbolo] (11.29,-13.43) -- (11.69,-13.20) -- (11.69,-13.66) -- cycle;

%--------------------------------------------------------------
% 0Z11: manometro
%--------------------------------------------------------------
\draw[simbolo] (2.96,-12.49) circle (0.46);
\draw[freccia] (3.15,-12.70) -- (2.74,-12.30);
\draw[linea] (2.96,-12.95) -- (2.96,-13.43);
\node[nodo] at (2.96,-13.43) {};
\node[anchor=south] at (2.96,-11.86) {0Z11};

%--------------------------------------------------------------
% 0V11: regolatore di pressione con molla regolabile, sfiato
% secondario e pilotaggio
%--------------------------------------------------------------
\draw[simbolo] (4.33,-12.73) rectangle (5.27,-13.66);
\draw[doppia] (4.37,-13.43) -- (5.23,-13.43);
% molla verticale sopra la valvola
\draw[simbolo] (5.07,-12.73) -- (5.50,-12.64) -- (4.93,-12.54) -- (5.50,-12.44)
              -- (4.93,-12.34) -- (5.50,-12.24) -- (5.11,-12.16);
% freccia di regolazione
\draw[freccia] (5.53,-12.56) -- (4.46,-12.30);
% sfiato secondario
\draw[simbolo] (5.27,-12.86) -- (5.50,-12.98) -- (5.27,-13.10);
% pilotaggio
\draw[pilota] (4.33,-13.49) -- (4.11,-13.66) -- (4.11,-14.13) -- (5.04,-14.13)
             -- (5.04,-13.66);
\node[anchor=south] at (4.83,-11.86) {0V11};

%--------------------------------------------------------------
% 0Z10: filtro con separatore d'acqua e indicatore di intasamento
%--------------------------------------------------------------
\draw[simbolo] (6.93,-12.79) -- (7.60,-13.43) -- (6.93,-14.09) -- (6.29,-13.43) -- cycle;
\draw[pilota] (6.93,-12.79) -- (6.93,-13.66);
\draw[simbolo] (6.54,-13.66) -- (7.33,-13.66) -- (6.93,-13.96) -- cycle;
\draw[linea] (6.93,-14.09) -- (6.93,-15.29);
% indicatore di intasamento
\draw[simbolo] (6.34,-12.06) rectangle (7.53,-12.77);
\draw[freccia] (6.88,-12.60) -- (6.49,-12.21);
\draw[simbolo] (7.19,-12.39) circle (0.16);
\draw[simbolo] (7.08,-12.28) -- (7.30,-12.50);
\draw[simbolo] (7.08,-12.50) -- (7.30,-12.28);
\node[anchor=south] at (6.93,-11.86) {0Z10};

%--------------------------------------------------------------
% 0V10: valvola di intercettazione manuale 3/2 con scarico
%--------------------------------------------------------------
\draw[simbolo] (8.57,-12.24) rectangle (9.54,-14.14);
\draw[simbolo] (8.57,-13.20) -- (9.54,-13.20);
% posizione di passaggio (porta chiusa a destra)
\draw[doppia] (8.61,-12.50) -- (9.49,-12.50);
\draw[simbolo] (9.24,-12.89) -- (9.24,-13.03);
\draw[simbolo] (9.24,-12.96) -- (9.54,-12.96);
% posizione di scarico (porta chiusa sulla linea)
\draw[doppia] (8.60,-13.43) -- (9.54,-13.89);
\draw[simbolo] (9.24,-13.36) -- (9.24,-13.51);
\draw[simbolo] (9.24,-13.43) -- (9.54,-13.43);
\draw[linea] (9.54,-13.89) -- (9.97,-13.89);
\draw[simbolo] (9.97,-13.79) -- (9.97,-14.00) -- (10.19,-13.89) -- cycle;
% azionamento manuale
\draw[linea] (9.24,-12.24) -- (9.24,-11.31);
\draw[linea] (9.00,-12.24) -- (9.00,-11.91) -- (9.11,-11.79) -- (9.00,-11.67)
            -- (9.00,-11.31);
\draw[simbolo] (8.89,-11.79) -- (8.97,-11.79);
\draw[simbolo] (8.89,-11.31) -- (9.34,-11.31);
\node[anchor=north] at (9.06,-14.26) {0V10};

%--------------------------------------------------------------
% Graffe e campi di applicazione
%--------------------------------------------------------------
\draw[linea, decorate, decoration={brace, amplitude=8pt}]
  (13.11,-5.64) -- (13.11,-0.29);
\node[testo, text width=5.2cm, anchor=west] at (13.77,-2.96)
  {Campo di applicazione della\\ parte del sistema di comando\\ legata alla sicurezza};

\draw[linea, decorate, decoration={brace, amplitude=8pt}]
  (13.11,-16.43) -- (13.11,-6.14);
\node[testo, text width=5.2cm, anchor=west] at (13.77,-11.46)
  {Possibile rilevanza per il\\ rispetto dei principi di\\ sicurezza di base e collaudati};

%--------------------------------------------------------------
% Cornice
%--------------------------------------------------------------
\draw[draw=black, line width=1pt]
  ([shift={(-0.45,0.45)}]current bounding box.north west) rectangle
  ([shift={(0.45,-0.45)}]current bounding box.south east);

\end{tikzpicture}
}
        \captionof{figure}{Campo di applicazione della EN ISO 13849 per i sistemi pneumatici.}
        \label{fig:08-02}
    \end{minipage}
\end{center}

Oltre alla parte del sistema di comando legata alla sicurezza, i circuiti a fluido presentati come esempi in questo capitolo contengono soltanto i componenti aggiuntivi necessari per la comprensione del sistema a fluido, oppure direttamente correlati alla tecnica di comando. I requisiti che devono essere soddisfatti dai sistemi a fluido sono descritti integralmente in [58; 59]. [60–63] costituiscono ulteriori norme pertinenti.
La maggior parte degli esempi di sistemi di comando sono comandi elettroidraulici o elettropneumatici. Una serie di requisiti di sicurezza su questi sistemi di comando è soddisfatta dalla parte elettrica del sistema di comando, ad esempio il requisito secondo cui le variazioni di energia sui sistemi di comando elettroidraulici devono essere controllate.
Negli esempi di comando qui descritti, la funzione di sicurezza richiesta è l'arresto di un movimento pericoloso o l'inversione di una direzione di movimento. La prevenzione dell'avviamento inatteso è implicitamente inclusa.

La funzione di sicurezza richiesta può tuttavia essere anche, ad esempio, un determinato livello di pressione o uno scarico di pressione.

Le strutture della maggior parte dei sistemi di comando a fluido sono realizzate nelle categorie 1, 3 o 4. Poiché la categoria B richiede già l'osservanza delle norme pertinenti e dei principi di sicurezza di base, i sistemi di comando a fluido di categoria B e 1 non differiscono sostanzialmente nella loro struttura di comando, ma soltanto per la maggiore affidabilità legata alla sicurezza delle relative valvole. Per questo motivo, il presente rapporto non presenta alcun sistema di comando a fluido di categoria B.

Ulteriori informazioni su idraulica e pneumatica si trovano sul sito web dell'IFA (www.dguv.de/ifa, Webcode: d1029520).

\subsection{Sistemi di comando elettronici ed elettronici programmabili}
\label{sec:08-1-3}
I componenti elettronici sono generalmente più sensibili alle influenze ambientali esterne rispetto ai componenti elettromeccanici. Se non vengono adottate misure particolari, l'impiego di componenti elettronici a temperature inferiori a 0°C è soggetto a vincoli sostanzialmente maggiori rispetto ai componenti elettromeccanici. Esistono inoltre influenze ambientali praticamente irrilevanti per gli elementi circuitali elettromeccanici, ma che rappresentano problemi cruciali per i sistemi elettronici, ossia qualsiasi disturbo elettromagnetico che si accoppia nei sistemi elettronici sotto forma di disturbo condotto o di campi elettromagnetici. In alcuni casi è richiesto uno sforzo maggiore affinché venga raggiunta un'adeguata immunità ai disturbi per l'uso industriale. Sui componenti elettronici l'esclusione di guasto è praticamente impossibile. Di conseguenza, la sicurezza non può in linea di principio essere garantita dalla progettazione di un particolare componente, ma soltanto da determinati concetti circuitali e dall'applicazione di misure appropriate per il controllo dei guasti.

Secondo gli elenchi dei guasti per i componenti elettrici/elettronici della EN ISO 13849-2, vengono essenzialmente ipotizzati i guasti di cortocircuito, circuito aperto, variazione di un parametro o di un valore, e guasti di tipo stuck-at (blocco su un valore fisso). Si tratta, senza eccezioni, di effetti di guasto ipotizzati come permanenti. I guasti transitori (che si verificano sporadicamente), come i soft error causati dall'inversione di carica di un condensatore in un chip a seguito di particelle ad alta energia quali le particelle alfa, possono generalmente essere rilevati solo con difficoltà e sono per lo più controllati mediante misure strutturali.

Il modo di guasto dei componenti elettronici è spesso difficile da valutare; in generale non è possibile definire un modo di guasto predominante. Ciò può essere illustrato con un esempio: se un relè o un contattore non viene azionato elettricamente, ossia se la corrente non attraversa la sua bobina, non vi è motivo per cui i contatti si chiudano, quando il componente viene utilizzato entro i vincoli della propria specifica. In altre parole, un relè o un contattore diseccitato non si inserisce spontaneamente a seguito di un guasto interno. La situazione è diversa per la maggior parte dei componenti elettronici, come i transistor. Anche se un transistor è interdetto, ossia in assenza di una corrente di base sufficientemente elevata, non può comunque essere esclusa la possibilità che esso diventi improvvisamente conduttivo senza influenza esterna, a seguito di un guasto interno, e di conseguenza, in determinate circostanze, avvii un movimento pericoloso. Questo svantaggio, dal punto di vista della sicurezza, dei componenti elettronici deve anch'esso essere controllato mediante un concetto circuitale idoneo. In particolare, quando vengono utilizzati moduli altamente integrati, può non essere possibile dimostrare che un dispositivo o un'apparecchiatura sia completamente esente da guasti anche all'inizio del proprio tempo di missione, ossia alla messa in servizio. Anche a livello di componente, i fabbricanti non sono più in grado di dimostrare l'assenza di guasti con una copertura di prova del 100\% per i circuiti integrati complessi. Una situazione analoga esiste per il software dell'elettronica programmabile.

A differenza dei circuiti elettromeccanici, i circuiti puramente elettronici presentano spesso il vantaggio che un cambiamento di stato può essere forzato dinamicamente. Ciò consente di raggiungere la DC richiesta a intervalli opportunamente brevi e senza alterazione dello stato dei segnali esterni (dinamica forzata). Tra i diversi canali sono necessarie misure di disaccoppiamento al fine di prevenire i guasti per causa comune. Tali misure consistono generalmente in contatti galvanicamente isolati, reti resistive o a diodi, circuiti di filtro, optoisolatori e trasformatori.

I guasti sistematici possono portare al guasto simultaneo di canali di elaborazione ridondanti, qualora ciò non venga impedito mediante una considerazione tempestiva, in particolare durante la fase di progettazione e integrazione. L'impiego di principi quali la corrente di riposo, la diversità o il sovradimensionamento consente di progettare circuiti elettronici robusti. Non dovrebbero essere trascurate le misure che rendono i canali di elaborazione insensibili alle influenze fisiche riscontrate, ad esempio, in un ambiente industriale. Tali influenze comprendono temperatura, umidità, polvere, vibrazioni, urti, atmosfere corrosive, influenze elettromagnetiche, interruzione della tensione, sovratensione e sottotensione.

Una SRP/CS di categoria 1 deve essere progettata e fabbricata con l'impiego di componenti ben collaudati e principi di sicurezza ben collaudati. Poiché componenti elettronici complessi quali PLC, microprocessori o ASIC non sono considerati ben collaudati nel senso della norma, il presente rapporto non contiene esempi corrispondenti di elettronica di categoria 1.
Gli esempi circuitali comprendono un'indicazione dell'efficacia, ossia del relativo Performance Level, delle misure richieste per l'evitamento/il controllo dei guasti per l'elettronica programmabile. Ulteriori dettagli si trovano al punto 6.3. Qualora vengano impiegati ASIC in uno sviluppo, sono richieste misure per l'evitamento dei guasti nel processo di sviluppo. Tali misure si trovano ad esempio nella IEC 61508-2 [48], che specifica un modello a V per lo sviluppo di un ASIC, basato sul modello a V già noto dallo sviluppo del software.

I seguenti punti meritano di essere menzionati, poiché tali questioni si presentano nella pratica:
\begin{itemize}
    \item in generale, due canali di una SRP/CS non devono essere instradati attraverso lo stesso circuito integrato. Per gli optoisolatori, questo requisito significa ad esempio che essi devono essere alloggiati in involucri separati quando vengono utilizzati per elaborare segnali provenienti da canali diversi.
    \item l'influenza dei sistemi operativi, ecc., deve essere presa in considerazione anche quando viene impiegata l'elettronica programmabile. Un PC standard e un tipico sistema operativo commerciale non sono idonei all'impiego in un sistema di comando legato alla sicurezza. La richiesta assenza di guasti (o, più realisticamente, la bassa incidenza di guasti) non può generalmente essere dimostrata con uno sforzo ragionevole, o non risulterà raggiungibile, su un sistema operativo che non è stato progettato per applicazioni legate alla sicurezza.
\end{itemize}

\section{Esempi di circuiti}
\label{sec:08-2}
La Tabella 8.1 mostra una panoramica degli esempi circuitali da 1 a 38. Ulteriori esempi si trovano in [22]. La Tabella 8.2 (pag. 105) contiene un elenco alfabetico delle principali abbreviazioni utilizzate negli esempi circuitali.

Nota: negli esempi contenenti più funzioni di sicurezza (17, 19, 23, 24), nello schema a blocchi legato alla sicurezza è rappresentata soltanto la prima funzione di sicurezza dell'esempio.

\begin{table}[ht]
    \centering
    \definecolor{tabblue}{RGB}{31,73,125}
    \setlength{\tabcolsep}{4pt}
    \renewcommand{\arraystretch}{1.2}
    \caption{Panoramica degli esempi circuitali.}
    \label{tab:08-1}
    \begin{tabularx}{\textwidth}{|>{\raggedright\arraybackslash}p{2.9cm}|>{\centering\arraybackslash}p{2.9cm}|C|C|C|}
        \hline
        \rowcolor{tabblue}
        \textcolor{white}{\textbf{PL raggiunto}} &
        \textcolor{white}{\textbf{Categoria\newline implementata}} &
        \multicolumn{3}{c|}{\textcolor{white}{\textbf{Tecnologia/esempio n.}}} \\
        \rowcolor{tabblue}
        & &
        \textcolor{white}{\textbf{Pneumatica}} &
        \textcolor{white}{\textbf{Idraulica}} &
        \textcolor{white}{\textbf{Elettrica}} \\
        \hline
        b & B &  &  & 1, 4 \\
        \hline
        \rowcolor{black!8}
        c & 1 & 2 & 3, 38 & 5, 6, 7 \\
        \hline
        c & 2 &  &  & 9 \\
        \hline
        \rowcolor{black!8}
        c & 3 &  &  & 10, 24 \\
        \hline
        d & 2 & 11 & 12 & 13 \\
        \hline
        \rowcolor{black!8}
        d & 3 & 14 & 15, 16 & 15, 16, 17, 18, 19, 20, 21, 22, 23, 24 \\
        \hline
        e & 3 & 25 & 27 & 28, 29, 30 \\
        \hline
        \rowcolor{black!8}
        e & 4 & 31 & 32, 33 & 33, 34, 35, 37 \\
        \hline
    \end{tabularx}
\end{table}

% Tabella lunga (longtable): puo' proseguire nella pagina successiva
\begingroup
\definecolor{tabblue}{RGB}{31,73,125}
\fontsize{9}{11}\selectfont
\setlength{\tabcolsep}{4pt}
\renewcommand{\arraystretch}{1.2}
\setlength{\LTcapwidth}{\textwidth}
\begin{longtable}{|p{2.4cm}|p{\dimexpr\textwidth-2.4cm-4\tabcolsep-3\arrayrulewidth\relax}|}
    \caption{Panoramica delle abbreviazioni utilizzate negli esempi circuitali.}
    \label{tab:08-2} \\
    \hline
    \rowcolor{tabblue}
    \textcolor{white}{\textbf{Abbreviazione}} & \textcolor{white}{\textbf{Forma estesa}} \\
    \hline
    \endfirsthead
    \hline
    \rowcolor{tabblue}
    \textcolor{white}{\textbf{Abbreviazione}} & \textcolor{white}{\textbf{Forma estesa}} \\
    \hline
    \endhead
    \hline
    \endfoot
    \rowcolor{black!8}
    {[D]} & Valori di B\textsubscript{10D} o \textit{MTTF}\textsubscript{D} tratti da banche dati (vedere ad esempio la sezione D.2.6) \\
    \hline
    {[E]} & Valori di B\textsubscript{10D} o \textit{MTTF}\textsubscript{D} stimati (vedere sopra) \\
    \hline
    \rowcolor{black!8}
    {[M]} & Valori di B\textsubscript{10D} o \textit{MTTF}\textsubscript{D} basati sulle indicazioni del fabbricante \\
    \hline
    {[S]} & Valori di B\textsubscript{10D} o \textit{MTTF}\textsubscript{D} basati sui dati riportati nella EN ISO 13849-1 (vedere ad esempio la Tabella D.2 del presente rapporto) \\
    \hline
    \rowcolor{black!8}
    µC & Microcontrollore \\
    \hline
    B\textsubscript{10} & Durata nominale: numero medio di cicli di commutazione (manovre) entro il quale si guasta il 10\% dei componenti considerati \\
    \hline
    \rowcolor{black!8}
    B\textsubscript{10D} & Durata nominale (pericolosa): numero medio di cicli di commutazione (manovre) entro il quale il 10\% dei componenti considerati subisce un guasto pericoloso \\
    \hline
    CBC & Combinazione frizione/freno \\
    \hline
    \rowcolor{black!8}
    CCF & Guasto di causa comune \\
    \hline
    CPU & Microprocessore (unità centrale di elaborazione) \\
    \hline
    \rowcolor{black!8}
    \textit{DC} & Copertura diagnostica \\
    \hline
    \textit{DC}\textsubscript{avg} & Copertura diagnostica media \\
    \hline
    \rowcolor{black!8}
    ESPE & Attrezzatura di protezione elettrosensibile \\
    \hline
    FIT & Numero di guasti in 10\textsuperscript{9} ore di funzionamento dei componenti (\textit{failures in time}) \\
    \hline
    \rowcolor{black!8}
    FMEA & Analisi dei modi e degli effetti dei guasti \\
    \hline
    FI & Convertitore di frequenza \\
    \hline
    \rowcolor{black!8}
    M & Motore \\
    \hline
    MPC & Comando multifunzione \\
    \hline
    \rowcolor{black!8}
    \textit{MTTF}\textsubscript{D} & Tempo medio al guasto pericoloso \\
    \hline
    n\textsubscript{op} & Numero medio annuo di manovre \\
    \hline
    \rowcolor{black!8}
    \textit{PFH}\textsubscript{D} & Probabilità media di un guasto pericoloso per ora \\
    \hline
    PL & Livello di prestazione (Performance Level) \\
    \hline
    \rowcolor{black!8}
    PL\textsubscript{r} & Livello di prestazione richiesto \\
    \hline
    PLC & Controllore logico programmabile \\
    \hline
    \rowcolor{black!8}
    RAM & Memoria ad accesso casuale (memoria variabile) \\
    \hline
    ROM & Memoria di sola lettura (memoria non modificabile) \\
    \hline
    \rowcolor{black!8}
    SBC & Comando sicuro del freno; fornisce un segnale di uscita per comandare un dispositivo di frenatura/bloccaggio \\
    \hline
    SDE & Diseccitazione sicura; scarico di una parte di un impianto \\
    \hline
    \rowcolor{black!8}
    SLS & Velocità limitata in sicurezza (vedere Tabella~\ref{tab:05-2}) \\
    \hline
    SRASW & Software applicativo legato alla sicurezza \\
    \hline
    \rowcolor{black!8}
    SRESW & Software integrato legato alla sicurezza \\
    \hline
    SRP/CS & Parte di un sistema di comando legata alla sicurezza \\
    \hline
    \rowcolor{black!8}
    SS1-r, SS1-t & Arresto sicuro 1 (vedere Tabella~\ref{tab:05-2}) \\
    \hline
    SS2-r, SS2-t & Arresto sicuro 2 (vedere Tabella~\ref{tab:05-2}) \\
    \hline
    \rowcolor{black!8}
    SSC & Arresto e chiusura sicuri; intrappolamento dell'aria compressa nelle camere del pistone senza regolazione della posizione ad anello chiuso \\
    \hline
    STO & Coppia disinserita in sicurezza (vedere Tabella~\ref{tab:05-2}) \\
    \hline
    \rowcolor{black!8}
    T\textsubscript{10D} & Tempo medio entro il quale il 10\% dei componenti considerati subisce un guasto pericoloso \\
    \hline
    THC & Comando bimanuale \\
\end{longtable}
\endgroup
\clearpage
\subsection{Esempio 1}
\label{sec:08-2-1}
\textit{Monitoraggio di posizione dei ripari mobili mediante interruttori di prossimità - categoria B - PL b}

% La minipage tiene insieme disegno e didascalia (figura non fluttuante)
\begin{center}
    \begin{minipage}{\textwidth}
        \centering
        \captionsetup{hypcap=false}
        \resizebox{0.6\textwidth}{!}{%==============================================================
%  Figura 8.3 - Monitoraggio di posizione dei ripari mobili
%  mediante interruttori di prossimita' (Esempio 1)
%  Adattamento in italiano della Figura 8.3 dell'IFA Report 2/2017
%
%  I simboli riproducono quelli dell'originale: le coordinate sono
%  ricavate dal disegno di partenza (1 cm = 70 px, origine
%  nell'angolo in alto a sinistra della cornice).
%
%  NOTA: nel disegno compaiono simboli matematici ($<$, $\sim$):
%  niente cambi di corpo del font (font=\large & c.), che farebbero
%  fallire la composizione; le etichette sono ingrandite con "scale".
%==============================================================
\begin{tikzpicture}[
  font=\normalsize,
  filo/.style={draw=black, line width=0.6pt},
  meccanico/.style={draw=black, line width=0.6pt, dash pattern=on 0.34cm off 0.23cm},
  etichetta/.style={inner sep=1pt, scale=1.55},
  grande/.style={inner sep=1pt, scale=2.1}
]

%--------------------------------------------------------------
% Cornice
%--------------------------------------------------------------
\draw[draw=black, line width=1.2pt] (0,0) rectangle (11.886,-12.071);

%--------------------------------------------------------------
% Linea di alimentazione del circuito di comando
%--------------------------------------------------------------
\draw[filo] (0.971,-0.514) -- (6.571,-0.514);
\fill[black] (3.10,-0.514) circle (0.07);

%--------------------------------------------------------------
% B1: interruttore di prossimita' con contatto di chiusura
%--------------------------------------------------------------
\draw[filo] (3.10,-0.514) -- (3.10,-2.729);
\draw[filo] (2.757,-2.643) -- (3.10,-3.60);
\draw[filo] (3.10,-3.60) -- (3.10,-11.629);
% collegamento di azionamento
\draw[filo] (2.243,-3.143) -- (2.486,-3.143);
\draw[filo] (2.700,-3.143) -- (2.936,-3.143);
% sensore
\draw[filo] (1.243,-3.143) -- (1.743,-2.643) -- (2.243,-3.143) -- (1.743,-3.643) -- cycle;
\draw[filo] (1.620,-2.766) -- (1.620,-3.520);
\draw[filo] (1.883,-2.783) -- (1.883,-3.503);
\node[etichetta] at (0.557,-3.143) {B1};

%--------------------------------------------------------------
% Q1: avviatore motore con sganciatore di minima tensione
%--------------------------------------------------------------
% organo di manovra con blocco meccanico
\draw[filo] (5.671,-5.186) rectangle (6.329,-5.820);
\draw[filo] (6.000,-5.186) -- (6.000,-5.820);
\draw[filo] (5.671,-5.503) -- (6.329,-5.503);
\draw[filo] (4.929,-5.357) -- (4.929,-5.629);
\draw[filo] (4.929,-5.503) -- (5.057,-5.503);
\draw[filo] (5.214,-5.503) -- (5.314,-5.643) -- (5.400,-5.503);
\draw[filo] (5.429,-5.503) -- (5.514,-5.503);
% collegamento meccanico verso il contatto principale
\draw[filo] (6.329,-5.503) -- (6.843,-5.503);
\draw[meccanico] (7.043,-5.503) -- (10.212,-5.503);
% collegamento meccanico verso lo sganciatore
\draw[meccanico] (6.000,-5.820) -- (6.000,-7.171);
% sganciatore di minima tensione
\draw[filo] (5.329,-7.171) rectangle (6.666,-7.943);
\node[etichetta] at (6.000,-7.557) {U\,$<$};
\draw[filo] (6.000,-7.943) -- (6.000,-11.629);
\draw[filo] (3.10,-11.629) -- (6.000,-11.629);
\node[etichetta] at (5.10,-6.557) {Q1};

%--------------------------------------------------------------
% Circuito di potenza: linea trifase L, contatto principale del
% Q1 e motore M
%--------------------------------------------------------------
\draw[filo] (10.42,-1.957) -- (10.42,-5.286);
\draw[filo] (10.004,-3.096) -- (10.804,-2.700);
\draw[filo] (10.004,-3.329) -- (10.804,-2.929);
\draw[filo] (10.004,-3.557) -- (10.804,-3.161);
\node[etichetta] at (10.39,-1.486) {L};
% contatto principale con apertura automatica
\draw[filo] (10.271,-5.286) -- (10.543,-5.286);
\draw[filo] (10.42,-5.381) circle (0.096);
\draw[filo] (10.147,-5.381) -- (10.42,-5.876);
\draw[filo] (10.42,-5.876) -- (10.42,-8.457);
% testo piu' lungo dell'originale inglese: allineato a sinistra ma
% spostato in modo da non toccare la linea L
\node[etichetta, anchor=base west] at (6.75,-4.45) {Avviatore};
\node[etichetta, anchor=base west] at (6.75,-5.08) {motore};
% motore trifase
\draw[filo] (10.443,-9.486) circle (1.029);
\node[grande] at (10.414,-9.00) {M};
\node[grande] at (10.53,-9.93) {3$\sim$};

\end{tikzpicture}
}
        \captionof{figure}{Monitoraggio di posizione dei ripari mobili mediante interruttori di prossimità.}
        \label{fig:08-03}
    \end{minipage}
\end{center}

\textbf{Funzioni di sicurezza}
\begin{itemize}
    \item funzione di arresto legata alla sicurezza, avviata da un mezzo di protezione: l'azionamento dell'interruttore di prossimità all'apertura del riparo mobile (riparo di sicurezza) avvia la funzione di sicurezza STO (Safe Torque Off).
\end{itemize}

\textbf{Descrizione funzionale}
\begin{itemize}
    \item l'apertura del riparo mobile (es. riparo di sicurezza) viene rilevata da un interruttore di prossimità B1, che agisce sul relè di minima tensione (under Voltage release) di un avviatore motore Q1. La diseccitazione di Q1 interrompe o impedisce movimenti o stati pericolosi;
    \item la funzione di sicurezza non può essere mantenuta a fronte di tutti i guasti dei componenti, ed è dipendente dall'affidabilità dei componenti.
    \item la rimozione del dispositivo di protezione viene rilevata;
    \item B1 non contiene misure di monitoraggio interne. Non sono realizzate ulteriori misure per il rilevamento dei guasti;
\end{itemize}

\textbf{Caratteristiche progettuali}
\begin{itemize}
    \item sono osservati i principi di sicurezza di base e sono soddisfatti i requisiti della categoria B. Sono realizzati circuiti di protezione (quali la protezione dei contatti), come descritto nei paragrafi iniziali del Capitolo 8. Il principio della corrente di riposo del relè di minima tensione è impiegato come principio di sicurezza di base;
    \item per l'azionamento dell'interruttore di prossimità è garantita una disposizione stabile del mezzo di protezione (schermo di sicurezza);
    \item a seconda della progettazione dell'interruttore di prossimità, potrebbe essere possibile un'elusione del funzionamento sicuro in modo ragionevolmente prevedibile. L'elusione può essere resa più difficoltosa, ad esempio mediante particolari condizioni di installazione, come l'installazione schermata/protetta (vedere anche EN ISO 14119);
    \item l'alimentazione dell'intera macchina viene disinserita (arresto di categoria 0 secondo la IEC 60204-1);
\end{itemize}

\begin{center}
    \begin{minipage}{\textwidth}
        \centering
        \captionsetup{hypcap=false}
        \resizebox{0.4\textwidth}{!}{%==============================================================
%  Figura 8.4 - Diagramma a blocchi (Esempio 1)
%  Adattamento in italiano della Figura 8.4 dell'IFA Report 2/2017
%
%  Coordinate ricavate dal disegno di partenza (1 cm = 70 px,
%  origine all'estremita' sinistra della linea di collegamento).
%==============================================================
\begin{tikzpicture}[
  font=\normalsize,
  filo/.style={draw=black, line width=0.6pt},
  blocco/.style={draw=black, line width=0.6pt, fill=black!10},
  etichetta/.style={inner sep=1pt, scale=1.1}
]

% Linea di collegamento
\draw[filo] (0,0) -- (1.043,0);
\draw[filo] (2.314,0) -- (3.371,0);
\draw[filo] (4.614,0) -- (5.500,0);

% Blocchi
\draw[blocco] (1.043,-0.586) rectangle (2.314,0.586);
\draw[blocco] (3.371,-0.586) rectangle (4.614,0.586);

\node[etichetta] at (1.679,0) {B1};
\node[etichetta] at (3.993,0) {Q1};

\end{tikzpicture}
}
        \captionof{figure}{Diagramma a blocchi.}
        \label{fig:08-04}
    \end{minipage}
\end{center}

\textbf{Calcolo della probabilità di guasto}\\
MTTF\textsubscript{D}: B1 è un convenzionale interruttore di prossimità su uno schermo di sicurezza, con un MTTF\textsubscript{D} di 1.100 anni [M]. Per il relè di minima tensione dell'avviatore motore Q1, il valore di B10 si approssima alla durata elettrica di 10.000 cicli di commutazione [M]. Nell'ipotesi che il 50\% dei guasti sia pericoloso, il valore di B\textsubscript{10D} si ottiene raddoppiando il valore di B10. Assumendo un azionamento giornaliero dell'interruttore di prossimità, un nop di 365 cicli all'anno per Q1 produce un MTTF\textsubscript{D} di 548 anni. Per la combinazione di B1 e Q1, l'MTTF\textsubscript{D} per ciascun canale è pari a 365 anni. Questo valore viene limitato al valore massimo aritmetico per la categoria B, ossia 27 anni ("medio").

La DC\textsubscript{avg} e le misure contro i guasti per causa comune non sono rilevanti nella categoria B.
Il sistema di comando elettromeccanico soddisfa la categoria B con un MTTF\textsubscript{D} medio (27 anni). Ne risulta una probabilità media di guasto pericoloso di 4,2 10\textsuperscript{-6} all'ora. Ciò soddisfa il PL b.
\clearpage
\subsection{Esempio 2}
\label{sec:08-2-2}
\textit{Valvola pneumatica (sottosistema) - categoria 1 - PL c}

% La minipage tiene insieme disegno e didascalia (figura non fluttuante)
\begin{center}
    \begin{minipage}{\textwidth}
        \centering
        \captionsetup{hypcap=false}
        \resizebox{0.8\textwidth}{!}{%==============================================================
%  Figura 8.5 - Valvola pneumatica per il comando di movimenti
%  pericolosi (Esempio 2)
%  Adattamento in italiano della Figura 8.5 dell'IFA Report 2/2017
%
%  I simboli dei componenti riproducono quelli dell'originale
%  (ISO 1219): le coordinate sono ricavate dal disegno di partenza
%  (1 cm = 70 px, origine nell'angolo in alto a sinistra della
%  cornice). Nessun cambio di corpo del font: le etichette sono
%  ingrandite con "scale" (vedere la nota nella Figura 8.3).
%==============================================================
\begin{tikzpicture}[
  font=\normalsize,
  linea/.style={draw=black, line width=0.8pt},
  simbolo/.style={draw=black, line width=0.8pt},
  sottile/.style={draw=black, line width=0.6pt},
  freccia/.style={sottile, -{Latex[length=8pt,width=3.6pt]}},
  doppia/.style={sottile, {Latex[length=8pt,width=3.6pt]}-{Latex[length=8pt,width=3.6pt]}},
  pilota/.style={draw=black, line width=0.6pt, dash pattern=on 3.6pt off 2pt},
  zona/.style={draw={rgb,255:red,22;green,67;blue,38}, line width=0.6pt,
               dash pattern=on 0.636cm off 0.1cm on 0.4pt off 0.1cm},
  etichetta/.style={inner sep=1pt, scale=1.2},
  testo/.style={align=center, inner sep=1pt, scale=1.2,
               execute at begin node=\setstretch{0.9}},
  nodo/.style={circle, fill=black, inner sep=0pt, minimum size=5.2pt}
]

%--------------------------------------------------------------
% Cornice
%--------------------------------------------------------------
\draw[draw=black, line width=1.2pt] (0.000,-0.000) rectangle (15.643,-18.664);

%--------------------------------------------------------------
% Cilindro 1A: camicia, pistone e stelo a doppia linea, blocco di
% attacco con le due bocche
%--------------------------------------------------------------
\draw[simbolo] (2.764,-0.750) rectangle (4.921,-1.943);
\draw[simbolo] (3.007,-0.750) -- (3.007,-1.943);
\draw[simbolo] (3.479,-0.750) -- (3.479,-1.943);
\draw[simbolo, fill=white] (3.479,-1.236) rectangle (5.629,-1.464);
\draw[linea] (2.821,-1.943) -- (2.821,-2.786) -- (3.586,-2.786) -- (3.586,-4.990);
\draw[linea] (4.843,-1.943) -- (4.843,-2.786) -- (4.061,-2.786) -- (4.061,-4.990);
\node[etichetta] at (2.086,-0.943) {1A};

% movimento pericoloso
\draw[sottile, {Latex[length=6pt,width=3.4pt]}-{Latex[length=6pt,width=3.4pt]}]
  (6.479,-1.347) -- (9.800,-1.347);
\draw[sottile] (8.150,-1.229) -- (8.150,-1.469);
\node[testo] at (8.143,-0.664) {movimento\\ pericoloso};

%--------------------------------------------------------------
% 1V1: valvola 4/3 a centro chiuso con scarichi, azionata da
% elettromagneti su entrambi i lati e centrata da molle
%--------------------------------------------------------------
\draw[simbolo] (2.396,-4.704) rectangle (5.261,-5.661);
\draw[simbolo] (3.357,-4.704) -- (3.357,-5.661);
\draw[simbolo] (4.310,-4.704) -- (4.310,-5.661);
% elettromagnete sinistro
\draw[simbolo] (0.476,-5.190) rectangle (2.396,-5.661);
\draw[simbolo] (1.429,-5.190) -- (1.429,-5.661);
\draw[sottile] (0.824,-5.661) -- (1.076,-5.190);
\draw[sottile] (1.429,-5.200) -- (1.847,-5.421) -- (1.429,-5.657);
% molla sinistra
\draw[sottile] (1.666,-4.943) -- (1.729,-5.190) -- (1.847,-4.704) -- (1.970,-5.190)
  -- (2.086,-4.704) -- (2.209,-5.190) -- (2.326,-4.704) -- (2.396,-4.943);
% elettromagnete destro
\draw[simbolo] (5.261,-5.190) rectangle (7.181,-5.661);
\draw[simbolo] (6.204,-5.190) -- (6.204,-5.661);
\draw[sottile] (6.586,-5.190) -- (6.810,-5.661);
\draw[sottile] (6.204,-5.200) -- (5.800,-5.421) -- (6.204,-5.657);
% molla destra
\draw[sottile] (5.991,-4.943) -- (5.929,-5.190) -- (5.810,-4.704) -- (5.687,-5.190)
  -- (5.571,-4.704) -- (5.449,-5.190) -- (5.331,-4.704) -- (5.261,-4.943);
% posizione sinistra
\draw[freccia] (2.857,-5.661) -- (2.621,-4.704);
\draw[freccia] (3.100,-4.704) -- (3.100,-5.661);
\draw[sottile] (2.619,-5.661) -- (2.619,-5.381);
\draw[sottile] (2.543,-5.381) -- (2.694,-5.381);
% posizione centrale (tutte le bocche chiuse)
\draw[sottile] (3.504,-4.990) -- (3.667,-4.990);
\draw[sottile] (3.981,-4.990) -- (4.143,-4.990);
\draw[sottile] (3.504,-5.381) -- (3.667,-5.381);
\draw[sottile] (3.743,-5.381) -- (3.906,-5.381);
\draw[sottile] (3.981,-5.381) -- (4.143,-5.381);
% posizione destra
\draw[freccia] (4.539,-4.704) -- (4.539,-5.661);
\draw[freccia] (4.776,-5.661) -- (5.014,-4.704);
\draw[sottile] (5.014,-5.661) -- (5.014,-5.381);
\draw[sottile] (4.939,-5.381) -- (5.090,-5.381);
% scarichi
\draw[linea] (3.586,-5.381) -- (3.586,-6.143);
\draw[linea] (4.061,-5.381) -- (4.061,-6.143);
\draw[sottile] (3.467,-6.143) -- (3.704,-6.143) -- (3.586,-6.347) -- cycle;
\draw[sottile] (3.943,-6.143) -- (4.180,-6.143) -- (4.061,-6.347) -- cycle;
\node[etichetta] at (0.871,-4.186) {1V1};
% valvola collaudata
\draw[sottile] (4.771,-4.700) -- (5.957,-4.091);
\node[testo, anchor=west] at (6.114,-3.807) {valvola collaudata per\\ applicazioni legate alla sicurezza};

%--------------------------------------------------------------
% Linea di alimentazione, derivazione verso ulteriori utenze e
% pressostato 0S1
%--------------------------------------------------------------
\draw[linea] (3.824,-5.381) -- (3.824,-11.157);
\node[nodo] at (3.824,-7.824) {};
\draw[linea] (3.824,-7.824) -- (5.261,-7.824);
\draw[sottile] (5.261,-7.696) -- (5.467,-7.824) -- (5.261,-7.933) -- cycle;
\node[testo, anchor=west] at (5.729,-7.764) {ulteriori utenze e\\ sistemi di comando};

\node[nodo] at (3.824,-9.490) {};
\draw[linea] (3.824,-9.490) -- (5.261,-9.490);
\draw[simbolo] (5.261,-9.019) rectangle (6.219,-9.976);
\draw[sottile] (5.261,-9.976) -- (6.219,-9.019);
\node[etichetta] at (5.557,-9.343) {P};
\draw[sottile] (5.600,-9.847) -- (5.800,-9.847) -- (5.800,-9.729) -- (5.967,-9.729);
\node[etichetta] at (4.629,-8.943) {0S1};

%--------------------------------------------------------------
% 0V1: valvola 3/2 con elettromagnete e ritorno a molla
%--------------------------------------------------------------
\draw[simbolo] (2.624,-11.157) rectangle (4.539,-12.114);
\draw[simbolo] (3.576,-11.157) -- (3.576,-12.114);
% elettromagnete
\draw[simbolo] (0.704,-11.647) rectangle (2.624,-12.114);
\draw[simbolo] (1.667,-11.647) -- (1.667,-12.114);
\draw[sottile] (1.076,-12.114) -- (1.300,-11.647);
\draw[sottile] (1.667,-11.647) -- (2.086,-11.886) -- (1.667,-12.114);
% molla
\draw[sottile] (4.539,-11.886) -- (4.600,-11.647) -- (4.719,-12.114) -- (4.847,-11.647)
  -- (4.957,-12.114) -- (5.086,-11.647) -- (5.204,-12.114) -- (5.261,-11.871);
% posizione di passaggio
\draw[freccia] (2.861,-12.114) -- (2.861,-11.157);
\draw[sottile] (3.357,-12.114) -- (3.357,-11.833);
\draw[sottile] (3.267,-11.833) -- (3.429,-11.833);
% posizione di scarico
\draw[freccia] (3.824,-11.157) -- (4.300,-12.114);
\draw[sottile] (3.733,-11.833) -- (3.900,-11.833);
\draw[linea] (4.300,-12.114) -- (4.300,-12.633);
\draw[sottile] (4.181,-12.633) -- (4.419,-12.633) -- (4.300,-12.833) -- cycle;
\node[etichetta] at (1.086,-10.757) {0V1};

%--------------------------------------------------------------
% Linea di alimentazione verso l'unita' di manutenzione 0Z
%--------------------------------------------------------------
\draw[linea] (3.824,-11.833) -- (3.824,-16.229) -- (7.926,-16.229);
\draw[linea] (8.879,-16.229) -- (9.647,-16.229);
\draw[linea] (11.000,-16.229) -- (12.239,-16.229);
\draw[linea] (13.190,-16.229) -- (14.976,-16.229);
% sorgente di aria compressa
\draw[sottile] (14.976,-16.229) -- (15.386,-15.990) -- (15.386,-16.476) -- cycle;

\draw[zona] (5.379,-13.593) rectangle (14.307,-17.721);
\node[etichetta] at (4.643,-13.800) {0Z};

% manometro
\draw[simbolo] (6.490,-15.271) circle (0.479);
\draw[sottile, -{Latex[length=6pt,width=3.6pt]}] (6.704,-15.490) -- (6.276,-15.061);
\draw[linea] (6.490,-15.750) -- (6.490,-16.229);
\node[nodo] at (6.490,-16.229) {};

% regolatore di pressione con molla regolabile, sfiato secondario
% e pilotaggio
\draw[simbolo] (7.926,-15.521) rectangle (8.879,-16.479);
\draw[doppia] (7.933,-16.229) -- (8.871,-16.229);
\draw[sottile] (8.879,-15.650) -- (9.129,-15.757) -- (8.879,-15.879);
\draw[sottile] (8.657,-15.521) -- (8.504,-15.450) -- (8.879,-15.364) -- (8.504,-15.264)
  -- (8.879,-15.179) -- (8.504,-15.079) -- (8.879,-14.986) -- (8.693,-14.936);
\draw[sottile, -{Latex[length=6pt,width=3.6pt]}] (9.121,-15.321) -- (8.040,-15.079);
\draw[pilota] (7.926,-16.229) -- (7.686,-16.497) -- (7.686,-16.897) -- (8.664,-16.897)
  -- (8.664,-16.479);

% filtro con separatore d'acqua e indicatore di intasamento
\draw[simbolo] (10.333,-15.547) -- (11.000,-16.229) -- (10.333,-16.910) -- (9.647,-16.229) -- cycle;
\draw[pilota] (10.333,-15.547) -- (10.333,-16.467);
\draw[sottile] (9.886,-16.467) -- (10.776,-16.467);
\draw[sottile] (10.014,-16.467) -- (10.333,-16.786) -- (10.657,-16.467);
\draw[linea] (10.333,-16.910) -- (10.333,-18.229);
\draw[simbolo] (9.753,-14.833) rectangle (10.933,-15.547);
\draw[sottile, -{Latex[length=6pt,width=3.6pt]}] (10.286,-15.404) -- (9.857,-14.976);
\draw[sottile] (10.581,-15.204) circle (0.186);
\draw[sottile] (10.450,-15.073) -- (10.713,-15.336);
\draw[sottile] (10.450,-15.336) -- (10.713,-15.073);

% valvola di intercettazione manuale 3/2 con scarico
\draw[simbolo] (12.239,-15.029) rectangle (13.190,-16.953);
\draw[simbolo] (12.239,-15.990) -- (13.190,-15.990);
% posizione di passaggio (bocca chiusa a destra)
\draw[doppia] (12.246,-15.276) -- (13.183,-15.276);
\draw[sottile] (12.910,-15.681) -- (12.910,-15.839);
\draw[sottile] (12.910,-15.753) -- (13.190,-15.753);
% posizione di scarico (bocca chiusa sulla linea)
\draw[doppia] (12.246,-16.231) -- (13.183,-16.701);
\draw[sottile] (12.910,-16.143) -- (12.910,-16.310);
\draw[sottile] (12.910,-16.229) -- (13.190,-16.229);
\draw[linea] (13.190,-16.704) -- (13.667,-16.704);
\draw[sottile] (13.667,-16.586) -- (13.881,-16.704) -- (13.667,-16.833) -- cycle;
% azionamento manuale con arresto
\draw[linea] (12.896,-15.029) -- (12.896,-14.076);
\draw[linea] (12.647,-15.029) -- (12.647,-14.667) -- (12.776,-14.553) -- (12.647,-14.443)
  -- (12.647,-14.076);
\draw[sottile] (12.529,-14.553) -- (12.647,-14.553);
\draw[simbolo] (12.529,-14.076) -- (13.014,-14.076);

\end{tikzpicture}
}
        \captionof{figure}{Valvola pneumatica per il comando di movimenti pericolosi.}
        \label{fig:08-05}
    \end{minipage}
\end{center}

\textbf{Funzioni di sicurezza}
\begin{itemize}
    \item funzione di arresto legata alla sicurezza: arresto del movimento pericoloso e prevenzione dell'avviamento inatteso dalla posizione di riposo, realizzata mediante la sottofunzione di sicurezza SSC.
    \item qui è mostrata soltanto la parte pneumatica del sistema di comando, sotto forma di sottosistema. Ulteriori parti dei sistemi di comando legate alla sicurezza (ad es. mezzi di protezione ed elementi logici elettrici) devono essere aggiunte sotto forma di sottosistemi per completare la funzione di sicurezza.
\end{itemize}

\textbf{Descrizione funzionale}
\begin{itemize}
    \item i movimenti pericolosi sono controllati da una valvola direzionale 1V1 ben collaudata per applicazioni di sicurezza.
    \item il guasto della valvola direzionale può comportare la perdita della funzione di sicurezza. Il guasto dipende dall'affidabilità della valvola direzionale.
    \item non sono realizzate misure per il rilevamento dei guasti.
    \item qualora l'aria compressa intrappolata costituisca un ulteriore pericolo, sono necessarie misure aggiuntive.
\end{itemize}

\textbf{Caratteristiche progettuali}
\begin{itemize}
    \item sono osservati i principi di sicurezza di base e ben collaudati, e sono soddisfatti i requisiti della categoria B.
    \item 1V1 è una valvola direzionale con posizione centrale chiusa, sovrapposizione sufficiente, posizione centrale centrata a molla e molle resistenti a fatica;
    \item la posizione di commutazione orientata alla sicurezza viene raggiunta mediante l'annullamento del segnale di comando.
    \item il fabbricante/utilizzatore deve confermare che la valvola direzionale è un componente ben collaudato per applicazioni di sicurezza (di affidabilità sufficientemente elevata);
    \item la funzione di sicurezza può essere ottenuta anche mediante una disposizione logica di valvole idonee;
\end{itemize}

\textbf{Calcolo della probabilità di guasto}\\
MTTF\textsubscript{D}: per la valvola direzionale 1V1 si assume un valore di B10D pari a 20.000.000 di cicli di commutazione [S]. Con 240 giorni lavorativi, 16 ore lavorative e un tempo di ciclo di 10 secondi, nop è pari a 1.382.400 cicli all'anno e l'MTTF\textsubscript{D} è di 145 anni. Questo è anche il valore di MTTF\textsubscript{D} per canale, che viene limitato a 100 anni ("alto").

La DC\textsubscript{avg} e le misure contro i guasti per causa comune non sono rilevanti nella categoria 1.

Il comando pneumatico soddisfa la categoria 1 con un MTTF\textsubscript{D} alto (100 anni). Ne risulta una probabilità media di guasto pericoloso di 1,1·10\textsuperscript{-6} all'ora. Ciò soddisfa il PL c. A seguito dell'aggiunta di ulteriori parti dei sistemi di comando legate alla sicurezza, sotto forma di sottosistemi, per il completamento della funzione di sicurezza, il PL potrebbe in determinate circostanze risultare inferiore. Tenendo conto della stima a favore della sicurezza sopra descritta, si ottiene un valore di 14 anni per il tempo di funzionamento (T\textsubscript{10D}) prima che la valvola direzionale 1V1, soggetta a usura, debba essere sostituita.
\clearpage
\subsection{Esempio 3}
\label{sec:08-2-3}
\textit{Valvola idraulica per il controllo di movimenti pericolosi}

% La minipage tiene insieme disegno e didascalia (figura non fluttuante)
\begin{center}
    \begin{minipage}{\textwidth}
        \centering
        \captionsetup{hypcap=false}
        \resizebox{0.85\textwidth}{!}{%==============================================================
%  Figura 8.7 - Valvola idraulica per il controllo di movimenti
%  pericolosi (Esempio 3)
%  Adattamento in italiano della Figura 8.7 dell'IFA Report 2/2017
%
%  I simboli dei componenti riproducono quelli dell'originale
%  (ISO 1219): le coordinate sono ricavate dal disegno di partenza
%  (1 cm = 70 px, origine nell'angolo in alto a sinistra della
%  cornice). Nessun cambio di corpo del font: le etichette sono
%  ingrandite con "scale" (vedere la nota nella Figura 8.3).
%==============================================================
\begin{tikzpicture}[
  font=\normalsize,
  linea/.style={draw=black, line width=0.8pt},
  simbolo/.style={draw=black, line width=0.8pt},
  sottile/.style={draw=black, line width=0.6pt},
  freccia/.style={sottile, -{Latex[length=9pt,width=4pt]}},
  pilota/.style={draw=black, line width=0.6pt, dash pattern=on 3.6pt off 2.4pt},
  elemento/.style={draw={rgb,255:red,87;green,120;blue,99}, line width=0.6pt,
                   dash pattern=on 0.25cm off 0.121cm},
  etichetta/.style={inner sep=1pt, scale=1.25},
  testo/.style={align=center, inner sep=1pt, scale=1.25,
                execute at begin node=\setstretch{0.9}},
  nodo/.style={circle, fill=black, inner sep=0pt, minimum size=4.8pt},
  nodogrande/.style={circle, fill=black, inner sep=0pt, minimum size=8pt}
]

%--------------------------------------------------------------
% Cornice
%--------------------------------------------------------------
\draw[draw=black, line width=1.2pt] (0.000,-0.000) rectangle (18.307,-19.671);

%--------------------------------------------------------------
% Cilindro 1A: camicia, pistone e stelo a doppia linea, blocco di
% attacco con le due bocche
%--------------------------------------------------------------
\draw[simbolo] (4.790,-1.157) rectangle (7.143,-2.443);
\draw[simbolo] (5.057,-1.157) -- (5.057,-2.443);
\draw[simbolo] (5.581,-1.157) -- (5.581,-2.443);
\draw[simbolo, fill=white] (5.581,-1.671) rectangle (7.924,-1.933);
\draw[linea] (4.853,-2.443) -- (4.853,-3.371) -- (5.700,-3.371) -- (5.700,-5.047);
\draw[linea] (7.081,-2.443) -- (7.081,-3.371) -- (6.210,-3.371) -- (6.210,-5.047);
\node[etichetta] at (4.171,-1.329) {1A};
% richiamo (*)
\draw[sottile] (7.243,-1.600) -- (7.910,-0.957);
\node[inner sep=0pt, scale=1.5] at (8.129,-0.814) {$\ast$};

% movimento pericoloso
\draw[sottile, {Latex[length=6pt,width=3.4pt]}-{Latex[length=6pt,width=3.4pt]}]
  (8.457,-1.804) -- (12.096,-1.804);
\draw[sottile] (10.271,-1.671) -- (10.271,-1.933);
\node[testo] at (10.271,-0.814) {movimento\\ pericoloso};

%--------------------------------------------------------------
% 1V3: valvola 4/3 a centro chiuso (P e T collegati), azionata da
% elettromagneti a e b e centrata da molle
%--------------------------------------------------------------
\draw[simbolo] (4.390,-4.786) rectangle (7.529,-5.833);
\draw[simbolo] (5.433,-4.786) -- (5.433,-5.833);
\draw[simbolo] (6.481,-4.786) -- (6.481,-5.833);
% elettromagnete a
\draw[simbolo] (2.814,-5.314) rectangle (4.390,-5.833);
\draw[simbolo] (3.600,-5.314) -- (3.600,-5.833);
\draw[sottile] (3.090,-5.833) -- (3.339,-5.314);
\fill[black] (3.600,-5.314) -- (3.876,-5.571) -- (3.600,-5.833) -- cycle;
% molla sinistra
\draw[sottile] (3.713,-5.059) -- (3.807,-4.791) -- (3.900,-5.309) -- (3.994,-4.791)
  -- (4.087,-5.309) -- (4.180,-4.791) -- (4.276,-5.309) -- (4.383,-5.057);
% elettromagnete b
\draw[simbolo] (7.529,-5.314) rectangle (9.081,-5.833);
\draw[simbolo] (8.304,-5.314) -- (8.304,-5.833);
\draw[sottile] (8.557,-5.314) -- (8.829,-5.833);
\fill[black] (8.304,-5.314) -- (8.033,-5.571) -- (8.304,-5.833) -- cycle;
% molla destra
\draw[sottile] (8.207,-5.059) -- (8.113,-4.791) -- (8.020,-5.309) -- (7.926,-4.791)
  -- (7.833,-5.309) -- (7.740,-4.791) -- (7.644,-5.309) -- (7.537,-5.057);
% posizione a: collegamenti incrociati
\draw[freccia] (4.647,-4.786) -- (5.171,-5.833);
\draw[freccia] (4.647,-5.833) -- (5.171,-4.786);
% posizione centrale: utilizzi chiusi, P collegato a T
\draw[sottile] (5.567,-5.047) -- (5.829,-5.047);
\draw[sottile] (6.076,-5.047) -- (6.353,-5.047);
\draw[sottile] (5.700,-5.833) -- (5.700,-5.314) -- (6.210,-5.314) -- (6.210,-5.833);
% posizione b: collegamenti paralleli
\draw[freccia] (6.733,-5.833) -- (6.733,-4.786);
\draw[freccia] (7.257,-4.786) -- (7.257,-5.833);
\node[etichetta] at (2.057,-4.843) {1V3};
\node[etichetta] at (2.257,-5.557) {a};
\node[etichetta] at (9.671,-5.557) {b};
% valvola collaudata
\draw[sottile] (8.394,-5.220) -- (9.086,-4.229);
\node[testo] at (10.600,-3.529) {valvola collaudata per\\ applicazioni legate alla sicurezza};

%--------------------------------------------------------------
% Linea di pressione P: filtro 1Z3, valvola limitatrice 1V2,
% valvola di non ritorno 1V1, pompa 1P
%--------------------------------------------------------------
\draw[linea] (5.700,-5.833) -- (5.700,-9.479);
\draw[linea] (5.700,-10.957) -- (5.700,-15.063);
\draw[linea] (5.700,-15.407) -- (5.700,-16.433);

% 1Z3: filtro con indicatore di intasamento
\node[nodo] at (5.700,-9.221) {};
\draw[linea] (5.700,-9.221) -- (6.421,-9.221);
\draw[simbolo] (6.421,-8.829) rectangle (7.743,-9.607);
\draw[sottile, -{Latex[length=6pt,width=3.6pt]}] (7.029,-9.457) -- (6.557,-8.979);
\draw[sottile] (7.343,-9.214) circle (0.2);
\draw[sottile] (7.201,-9.073) -- (7.484,-9.356);
\draw[sottile] (7.201,-9.356) -- (7.484,-9.073);
\draw[simbolo] (5.700,-9.479) -- (6.436,-10.214) -- (5.700,-10.957) -- (4.950,-10.214) -- cycle;
\draw[elemento] (4.950,-10.214) -- (6.436,-10.214);
\node[etichetta] at (4.314,-10.229) {1Z3};

% 1V2: valvola limitatrice di pressione con molla regolabile
\node[nodogrande] at (5.700,-13.569) {};
\draw[linea] (5.700,-13.569) -- (7.257,-13.569);
\draw[simbolo] (7.257,-13.043) rectangle (8.304,-14.086);
\draw[sottile, -{Latex[length=8pt,width=3.6pt]}] (7.257,-13.829) -- (8.304,-13.829);
\draw[sottile] (7.464,-12.419) -- (7.257,-12.461) -- (7.686,-12.576) -- (7.257,-12.676)
  -- (7.686,-12.764) -- (7.257,-12.871) -- (7.686,-12.990) -- (7.543,-13.043);
\draw[sottile, -{Latex[length=8pt,width=4pt]}] (7.007,-12.843) -- (8.193,-12.569);
\draw[pilota] (7.257,-13.569) -- (6.740,-14.079) -- (6.740,-14.614) -- (7.526,-14.614)
  -- (7.526,-14.086);
\draw[linea] (8.304,-13.569) -- (9.079,-13.569) -- (9.079,-19.053);
\node[etichetta] at (6.443,-13.186) {1V2};

% 1V1: valvola di non ritorno con molla
\draw[sottile] (5.700,-14.550) -- (5.971,-14.627) -- (5.419,-14.711) -- (5.971,-14.773)
  -- (5.419,-14.857) -- (5.971,-14.926) -- (5.419,-15.011) -- (5.700,-15.069);
\draw[simbolo] (5.700,-15.201) circle (0.139);
\draw[simbolo] (5.499,-15.201) -- (5.700,-15.407) -- (5.886,-15.201);
\node[etichetta] at (4.671,-14.800) {1V1};

% 1P: pompa a cilindrata fissa con senso di rotazione
\draw[simbolo] (5.700,-17.219) circle (0.790);
\fill[black] (5.700,-16.433) -- (5.971,-16.886) -- (5.424,-16.886) -- cycle;
\draw[sottile, -{Latex[length=8pt,width=4pt]}] (5.700,-17.219) ++(-30:1.171) arc[start angle=-30, end angle=30, radius=1.171];
\draw[linea] (5.700,-18.009) -- (5.700,-19.053);
\node[etichetta] at (7.629,-17.157) {1P};

% 1M: motore elettrico trifase e giunto
\draw[simbolo] (2.329,-17.219) circle (0.786);
\node[etichetta] at (2.329,-16.919) {M};
\node[etichetta] at (2.281,-17.396) {3\,$\sim$};
\draw[sottile] (3.104,-17.096) -- (3.890,-17.096);
\draw[sottile] (3.104,-17.357) -- (3.890,-17.357);
\draw[simbolo] (3.890,-16.957) -- (3.890,-17.481);
\draw[simbolo] (4.139,-16.957) -- (4.139,-17.481);
\draw[sottile] (4.139,-17.096) -- (4.914,-17.096);
\draw[sottile] (4.139,-17.357) -- (4.914,-17.357);
\node[etichetta] at (0.961,-17.204) {1M};

%--------------------------------------------------------------
% Linea di ritorno T: filtro 1Z2 con valvola di bypass
%--------------------------------------------------------------
\draw[linea] (6.210,-5.833) -- (6.210,-7.200) -- (11.257,-7.200) -- (11.257,-15.743);
\draw[linea] (11.257,-17.219) -- (11.257,-19.053);

\draw[simbolo] (9.424,-15.214) rectangle (12.339,-17.739);
\node[nodo] at (11.257,-15.433) {};
\node[nodo] at (11.257,-17.524) {};
\draw[simbolo] (11.257,-15.743) -- (12.004,-16.481) -- (11.257,-17.219) -- (10.519,-16.481) -- cycle;
\draw[elemento] (10.519,-16.481) -- (12.004,-16.481);
% bypass con valvola di non ritorno a molla
\draw[linea] (11.257,-15.433) -- (9.963,-15.433) -- (9.963,-16.067);
\draw[simbolo] (9.766,-16.184) -- (9.963,-15.997) -- (10.141,-16.184);
\draw[simbolo] (9.963,-16.201) circle (0.134);
\draw[sottile] (9.963,-16.309) -- (10.230,-16.407) -- (9.686,-16.500) -- (10.230,-16.594)
  -- (9.686,-16.690) -- (10.230,-16.783) -- (9.686,-16.871) -- (9.963,-16.961);
\draw[linea] (9.963,-16.336) -- (9.963,-17.524) -- (11.257,-17.524);
% indicatore di intasamento
\draw[linea] (11.257,-15.433) -- (12.890,-15.433);
\draw[simbolo] (12.890,-15.047) rectangle (14.186,-15.833);
\draw[sottile, -{Latex[length=6pt,width=3.6pt]}] (13.481,-15.671) -- (13.004,-15.196);
\draw[sottile] (13.796,-15.443) circle (0.207);
\draw[sottile] (13.649,-15.296) -- (13.943,-15.590);
\draw[sottile] (13.649,-15.590) -- (13.943,-15.296);
\node[etichetta] at (9.743,-14.729) {1Z2};

%--------------------------------------------------------------
% Serbatoio con indicatore di livello 1S1, termometro 1S2 e
% filtro di aerazione 1Z1
%--------------------------------------------------------------
\draw[simbolo] (1.004,-18.214) -- (1.004,-19.324) -- (17.767,-19.324) -- (17.767,-18.204);

% 1S1
\draw[simbolo] (13.581,-17.833) circle (0.526);
\draw[sottile] (13.581,-17.433) -- (13.581,-18.167);
\node[nodo, minimum size=4.5pt] at (13.581,-18.167) {};
\draw[linea] (13.581,-18.357) -- (13.581,-19.047);
\node[etichetta] at (13.533,-16.886) {1S1};

% 1S2
\draw[simbolo] (15.153,-17.833) circle (0.529);
\draw[sottile] (14.624,-17.833) -- (15.681,-17.833);
\fill[black] (15.004,-17.596) -- (15.290,-17.596) -- (15.153,-17.833) -- cycle;
\draw[linea] (15.153,-18.361) -- (15.153,-19.047);
\node[etichetta] at (15.104,-16.886) {1S2};

% 1Z1
\draw[simbolo] (16.971,-16.896) -- (17.710,-17.633) -- (16.971,-18.371) -- (16.224,-17.633) -- cycle;
\draw[elemento] (16.224,-17.633) -- (17.710,-17.633);
\draw[sottile] (16.843,-16.671) -- (17.100,-16.671) -- (16.971,-16.896) -- cycle;
\draw[sottile] (16.971,-16.671) -- (16.971,-16.410);
\draw[sottile] (16.843,-16.410) -- (17.100,-16.410) -- (16.971,-16.181) -- cycle;
\draw[linea] (16.971,-18.371) -- (16.971,-19.047);
\node[etichetta] at (16.210,-16.204) {1Z1};

\end{tikzpicture}
}
        \captionof{figure}{Valvola idraulica per il controllo di movimenti pericolosi.}
        \label{fig:08-07}
    \end{minipage}
\end{center}

\textbf{Funzioni di sicurezza}
\begin{itemize}
    \item funzione di arresto legata alla sicurezza: arresto del movimento pericoloso e prevenzione dell'avviamento inatteso dalla posizione di riposo
    \item qui è mostrata soltanto la parte idraulica del sistema di comando, sotto forma di sottosistema. Ulteriori parti dei sistemi di comando legate alla sicurezza (ad es. mezzi di protezione ed elementi logici elettrici) devono essere aggiunte sotto forma di sottosistemi per completare la funzione di sicurezza.
\end{itemize}

\textbf{Descrizioone funzionale}
\begin{itemize}
    \item i movimenti pericolosi sono controllati da una valvola direzionale 1V3 ben collaudata per applicazioni di sicurezza.
    \item il guasto della valvola direzionale può comportare la perdita della funzione di sicurezza. Il guasto dipende dall'affidabilità della valvola direzionale.
    \item non sono realizzate misure per il rilevamento dei guasti.
\end{itemize}

\textbf{Caratteristiche progettuali}
\begin{itemize}
    \item sono osservati i principi di sicurezza di base e ben collaudati, e sono soddisfatti i requisiti della categoria B.
    \item 1V3 è una valvola direzionale con posizione centrale chiusa, sovrapposizione sufficiente, posizione centrale centrata a molla e molle resistenti a fatica.
    \item la posizione di commutazione orientata alla sicurezza viene raggiunta mediante l'annullamento del segnale di comando.
    \item ove necessario, il fabbricante/utilizzatore deve confermare che la valvola direzionale è un componente ben collaudato per applicazioni di sicurezza.
    \item sono realizzate le seguenti misure specifiche per aumentare l'affidabilità della valvola direzionale: un filtro in pressione 1Z3 a monte della valvola direzionale, e misure idonee sul cilindro per impedire che lo sporco venga trascinato all'interno dallo stelo del pistone (es. raschiatore efficace sullo stelo del pistone, vedere * nella Figura 8.7).
\end{itemize}

\textbf{Calcolo della probabilità di guasto}
MTTF\textsubscript{D}: per la valvola direzionale 1V3 si assume un MTTF\textsubscript{D} di 150 anni [M]. Questo è anche il valore di MTTF\textsubscript{D} per canale, che viene limitato a 100 anni ("alto").
La DC\textsubscript{avg} e le misure contro i guasti per causa comune non sono rilevanti nella categoria 1.
Il comando idraulico soddisfa la categoria 1 con un MTTF\textsubscript{D} alto (100 anni). Ne risulta una probabilità media di guasto pericoloso di 1,1 10\textsuperscript{-6} all'ora. Ciò soddisfa il PL c. A seguito dell'aggiunta di ulteriori parti del sistema di comando legate alla sicurezza, sotto forma di sottosistemi, per il completamento della funzione di sicurezza, il PL potrebbe in determinate circostanze risultare inferiore.
\clearpage

\subsection{Esempio 4}
\label{sec:08-2-4}
\textbf{Esempio 4}
\textit{Arresto di una macchina per la lavorazione del legno - categoria B - PL b}

% La minipage tiene insieme disegno e didascalia (figura non fluttuante)
\begin{center}
    \begin{minipage}{\textwidth}
        \centering
        \captionsetup{hypcap=false}
        \resizebox{0.85\textwidth}{!}{%==============================================================
%  Figura 8.9 - Combinazione di un equipaggiamento di comando
%  elettromeccanico e di un semplice dispositivo di frenatura
%  elettronico per l'arresto di macchine per la lavorazione del
%  legno (Esempio 4)
%  Adattamento in italiano della Figura 8.9 dell'IFA Report 2/2017
%
%  I simboli riproducono quelli dell'originale: le coordinate sono
%  ricavate dal disegno di partenza (1 cm = 70 px, origine
%  nell'angolo in alto a sinistra della cornice). Nessun cambio di
%  corpo del font: le etichette sono ingrandite con "scale" (vedere
%  la nota nella Figura 8.3).
%==============================================================
\definecolor{verdemotore}{RGB}{22,67,38}
\begin{tikzpicture}[
  font=\normalsize,
  filo/.style={draw=black, line width=0.8pt},
  sottile/.style={draw=black, line width=0.6pt},
  etichetta/.style={inner sep=1pt, scale=1.3},
  testo/.style={align=left, inner sep=1pt, scale=1.3,
                execute at begin node=\setstretch{0.8}},
  nodo/.style={circle, fill=black, inner sep=0pt, minimum size=7.3pt}
]

%--------------------------------------------------------------
% Cornice
%--------------------------------------------------------------
\draw[draw=black, line width=1.2pt] (0.000,-0.000) rectangle (19.657,-18.093);

%--------------------------------------------------------------
% Circuito di comando: linea di alimentazione superiore
%--------------------------------------------------------------
\draw[filo] (2.193,-1.150) -- (7.371,-1.150);
\node[nodo] at (4.139,-1.150) {};

% S1/Off: pulsante di arresto (contatto di apertura) con simbolo
% di azionamento
\draw[filo] (4.139,-1.150) -- (4.139,-7.957) -- (4.529,-7.957);
\draw[filo] (4.137,-8.686) -- (4.437,-7.834);
\draw[filo] (4.139,-8.686) -- (4.139,-10.207);
\draw[filo] (3.157,-8.080) -- (3.157,-8.457);
\draw[filo] (3.157,-8.080) -- (3.357,-8.080);
\draw[filo] (3.157,-8.457) -- (3.357,-8.457);
\draw[filo] (3.157,-8.271) -- (3.257,-8.271);
\draw[filo] (3.343,-8.271) -- (3.691,-8.271);
\draw[filo] (3.763,-8.271) -- (4.023,-8.271);
\draw[filo] (4.137,-8.271) -- (4.283,-8.271);
\draw[sottile] (3.014,-7.214) circle (0.381);
\draw[sottile] (2.700,-7.214) -- (3.323,-7.214);
\draw[sottile] (3.006,-6.977) -- (3.323,-7.214) -- (3.006,-7.443);
\node[etichetta] at (1.843,-8.229) {S1/Off};

% S2/On (contatto di chiusura) e contatto di autoritenuta Q1
\node[nodo] at (4.139,-10.207) {};
\draw[filo] (4.139,-10.207) -- (4.139,-10.968);
\draw[filo] (3.611,-10.832) -- (4.139,-11.700);
\draw[filo] (4.139,-11.700) -- (4.139,-14.333);
\draw[filo] (2.746,-11.075) -- (2.746,-11.457);
\draw[filo] (2.746,-11.075) -- (2.954,-11.075);
\draw[filo] (2.746,-11.457) -- (2.954,-11.457);
\draw[filo] (2.746,-11.271) -- (2.843,-11.271);
\draw[filo] (2.932,-11.271) -- (3.157,-11.271);
\draw[filo] (3.354,-11.271) -- (3.621,-11.271);
\draw[filo] (3.729,-11.271) -- (3.879,-11.271);
\node[etichetta] at (1.457,-11.243) {S2/On};

\draw[filo] (4.139,-10.207) -- (6.014,-10.207) -- (6.014,-10.968);
\draw[filo] (5.486,-10.832) -- (6.014,-11.700);
\draw[filo] (6.014,-11.700) -- (6.014,-12.450);
\node[etichetta] at (5.000,-11.186) {Q1};

\node[nodo] at (4.139,-12.450) {};
\node[nodo] at (6.014,-12.450) {};
\draw[filo] (4.139,-12.450) -- (8.257,-12.450) -- (8.257,-5.153) -- (10.143,-5.153);

% bobina di Q1
\draw[filo] (3.381,-14.333) rectangle (4.886,-15.081);
\draw[filo] (4.139,-15.081) -- (4.139,-16.581);
\node[etichetta] at (2.686,-14.686) {Q1};

% linea di alimentazione inferiore e messa a terra
\draw[filo] (0.753,-16.581) -- (5.647,-16.581);
\node[nodo] at (4.139,-16.581) {};
\node[nodo] at (2.010,-16.581) {};
\draw[filo] (2.010,-16.581) -- (2.010,-17.139);
\draw[filo] (1.633,-17.139) -- (2.396,-17.139);
\draw[filo] (1.719,-17.343) -- (2.310,-17.343);
\draw[filo] (1.810,-17.539) -- (2.214,-17.539);

%--------------------------------------------------------------
% K1: dispositivo di frenatura elettronico
%--------------------------------------------------------------
\draw[filo] (12.014,-1.200) -- (12.014,-3.453);
\draw[sottile] (11.629,-2.504) -- (12.397,-2.119);
\draw[sottile] (11.629,-2.719) -- (12.397,-2.333);
\node[etichetta] at (12.043,-0.771) {L};
\node[testo, anchor=east] at (11.271,-2.507) {Dispositivo di\\ frenatura elettronico};

\draw[filo] (10.143,-3.453) rectangle (13.881,-6.833);
\node[etichetta] at (9.414,-3.657) {K1};
\node[inner sep=1pt, scale=1.1, rotate=90] at (10.624,-5.143) {Avvio frenatura};
% tiristore
\draw[sottile] (12.014,-4.200) -- (12.014,-6.090);
\draw[sottile] (11.633,-4.581) -- (12.396,-4.581) -- (12.014,-5.333) -- cycle;
\draw[sottile] (11.633,-5.333) -- (12.396,-5.333);
\draw[sottile] (12.276,-5.333) -- (12.381,-5.714) -- (12.767,-5.714);

% contatto di apertura di Q1
\draw[filo] (12.014,-6.833) -- (12.014,-8.714) -- (12.400,-8.714);
\draw[filo] (12.320,-8.509) -- (12.014,-9.443);
\node[nodo, minimum size=6.7pt] at (12.320,-8.509) {};
\draw[filo] (12.014,-9.443) -- (12.014,-11.329) -- (18.014,-11.329);
\draw[sottile] (14.693,-11.686) -- (15.157,-10.857);
\draw[sottile] (14.857,-11.786) -- (15.321,-10.971);
\node[etichetta] at (11.086,-9.014) {Q1};

%--------------------------------------------------------------
% Circuito di potenza: linea trifase L, contatto principale di Q1
% e motore M
%--------------------------------------------------------------
\draw[filo] (18.014,-1.200) -- (18.014,-4.961);
\draw[sottile] (17.624,-3.043) -- (18.396,-2.661);
\draw[sottile] (17.624,-3.257) -- (18.396,-2.876);
\draw[sottile] (17.624,-3.461) -- (18.396,-3.081);
\node[etichetta] at (18.029,-0.800) {L};
% contatto principale con apertura automatica
\draw[sottile] (18.014,-4.571) arc[start angle=90, end angle=270, radius=0.196];
\draw[filo] (17.481,-4.833) -- (18.014,-5.700);
\draw[filo] (18.014,-5.700) -- (18.014,-13.579);
\node[nodo] at (18.014,-11.329) {};
\node[etichetta] at (17.114,-5.457) {Q1};
% motore trifase
\draw[draw=verdemotore, line width=0.8pt] (17.943,-14.757) circle (1.179);
\node[inner sep=1pt, scale=1.9, text=verdemotore, font=\sffamily\bfseries\boldmath]
  at (17.943,-14.257) {M};
\node[inner sep=1pt, scale=1.9, text=verdemotore, font=\sffamily\bfseries\boldmath]
  at (18.029,-15.164) {3$\sim$};

\end{tikzpicture}
}
        \captionof{figure}{Combinazione di un equipaggiamento di comando elettromeccanico e di un semplice dispositivo di frenatura elettronico per l'arresto di macchine per la lavorazione del legno.}
        \label{fig:08-09}
    \end{minipage}
\end{center}

\textbf{Funzioni di sicurezza}
\begin{itemize}
    \item l'azionamento del pulsante di arresto (Off) determina l'SS1-t (safe stop 1, time controlled — arresto sicuro 1, controllato a tempo), ossia un arresto controllato del motore entro un tempo massimo ammissibile.
\end{itemize}

\textbf{Descrizione funzionale}
\begin{itemize}
    \item l'arresto del motore viene avviato dall'azionamento del pulsante di arresto S1. Il contattore motore Q1 si diseccita e viene avviata la funzione di frenatura. Il motore viene frenato mediante una corrente continua, generata nel dispositivo di frenatura K1 da un controllo dell'angolo di fase con tiristori, che genera una coppia frenante nell'avvolgimento del motore.
    \item il tempo di arresto non deve superare un valore massimo (ad es. 10 secondi). Il livello di corrente di frenatura necessario a tale scopo può essere impostato mediante un potenziometro sul dispositivo di frenatura.
    \item trascorso il tempo massimo di frenatura, il tiristore non viene più attivato e il percorso di corrente per la corrente di frenatura viene interrotto. Il processo di arresto corrisponde a un arresto di categoria 1 secondo la IEC 60204-1.
    \item la funzione di sicurezza non può essere mantenuta a fronte di tutti i guasti dei componenti, ed è dipendente dall'affidabilità dei componenti.
    \item non sono realizzate misure per il rilevamento dei guasti
\end{itemize}

\textbf{Caratteristiche progettuali}
\begin{itemize}
    \item sono osservati i principi di sicurezza di base e ben collaudati, e sono soddisfatti i requisiti della categoria B. Sono realizzati circuiti di protezione (quali la protezione dei contatti), come descritto nei paragrafi iniziali del Capitolo 8. Il principio di diseccitazione (principio della corrente di riposo) è applicato come principio di sicurezza di base. Per la protezione contro l'avviamento inatteso a seguito del ripristino dell'alimentazione, il sistema di comando è dotato di un'autoritenuta.
    \item S1 è un pulsante con apertura ad azione positiva secondo la IEC 60947-5-1, Allegato K. S1 è pertanto considerato un componente ben collaudato.
    \item il contattore Q1 è un componente ben collaudato, tenendo conto delle condizioni aggiuntive secondo la Tabella D.3 della EN ISO 13849-2.
    \item il dispositivo di frenatura K1 è costruito interamente con componenti elettronici semplici, quali transistor, condensatori, diodi, resistori e tiristori. Il comportamento legato alla sicurezza è determinato dalla selezione dei componenti. Non sono realizzate misure interne per il rilevamento dei guasti.
\end{itemize}

\textbf{Applicazione}
Su macchine per la lavorazione del legno o macchine analoghe, sulle quali un arresto non frenato comporterebbe una durata di arresto per inerzia inammissibilmente lunga dei movimenti pericolosi dell'utensile. Il comando della funzione di frenatura sulle macchine per la lavorazione del legno deve essere progettato in modo tale da raggiungere almeno il PL b (prEN ISO 19085-1:2014).

\textbf{Calcolo della probabilità di guasto}
\begin{itemize}
    \item il pulsante S1 e il contattore Q1 sono combinati, ai fini del calcolo in SISTEMA, in un sottosistema che soddisfa i requisiti della categoria 1. Il dispositivo di frenatura K1 costituisce un sottosistema separato in categoria B.
    \item S1 è un pulsante con apertura ad azione positiva secondo la IEC 60947-5-1, Allegato K.
    \item MTTF\textsubscript{D}: per il pulsante S1 è specificato un valore di B10D pari a 20 10⁶ cicli di commutazione [M]. Per il contattore Q1 si assume un valore di B10D pari a 1.300.000 cicli di commutazione [S] a carico nominale. Con 300 giorni lavorativi, 8 ore lavorative e un tempo di ciclo di 2 minuti, nop è pari a 72.000 cicli all'anno. L'MTTF\textsubscript{D} è di 2.777 anni per il pulsante S1 e di 180 anni per Q1. Complessivamente, ne risulta un MTTF\textsubscript{D} di 169 anni, che conformemente alla norma viene ridotto a 100 anni ("alto") per il sottosistema. Il contattore Q1 ha un tempo di funzionamento limitato (T10D) di 18 anni. Se ne raccomanda la sostituzione tempestiva. L'MTTF\textsubscript{D} per il dispositivo di frenatura K1 è stato determinato mediante il metodo del conteggio dei componenti. Le informazioni sui componenti tratte dalla distinta dei componenti e i valori dalla banca dati SN 29500 [48] producono un MTTF\textsubscript{D} di 518 anni [D]. Anche questo viene ridotto a 100 anni ("alto").
    \item la DC\textsubscript{avg} e le misure contro i guasti per causa comune non sono rilevanti nella categoria B e nella categoria 1.
    \item il sottosistema S1/Q1 soddisfa la categoria 1 con un MTTFD alto (100 anni). Ne risulta una probabilità media di guasto pericoloso di 1,1 10\textsuperscript{-6} all'ora. Ciò soddisfa il PL c.
    \item il sottosistema K1 soddisfa la categoria B con un MTTFD alto (100 anni). Ne risulta una probabilità media di guasto pericoloso di 4,2 10\textsuperscript{-6} all'ora. Ciò soddisfa il PL b.
    \item per la funzione di arresto legata alla sicurezza, la probabilità media di guasto pericoloso risultante è di 5,4 10\textsuperscript{-6} all'ora. Ciò soddisfa il PL b
\end{itemize}
\clearpage

\subsection{Esempio 5}
\label{sec:08-2-5}
\textbf{Esempio 5}
\textit{Controllo della posizione dei ripari mobili - Categoria 1 - PL c}

% La minipage tiene insieme disegno e didascalia (figura non fluttuante)
\begin{center}
    \begin{minipage}{\textwidth}
        \centering
        \captionsetup{hypcap=false}
        \resizebox{0.6\textwidth}{!}{%==============================================================
%  Figura 8.11 - Controllo della posizione del riparo mobile per
%  la prevenzione di movimenti pericolosi (STO) (Esempio 5)
%  Adattamento in italiano della Figura 8.11 dell'IFA Report 2/2017
%
%  I simboli riproducono quelli dell'originale: le coordinate sono
%  ricavate dal disegno di partenza (1 cm = 70 px, origine
%  nell'angolo in alto a sinistra della cornice). Nessun cambio di
%  corpo del font: le etichette sono ingrandite con "scale" (vedere
%  la nota nella Figura 8.3).
%==============================================================
\definecolor{verdepotenza}{RGB}{22,67,38}
\begin{tikzpicture}[
  font=\normalsize,
  filo/.style={draw=black, line width=0.6pt},
  potenza/.style={draw=verdepotenza, line width=0.6pt},
  meccanico/.style={draw=black, line width=0.6pt, dash pattern=on 0.286cm off 0.21cm},
  etichetta/.style={inner sep=1pt, scale=1.5},
  nodo/.style={circle, fill=black, inner sep=0pt, minimum size=4.2pt}
]

%--------------------------------------------------------------
% Cornice
%--------------------------------------------------------------
\draw[draw=black, line width=1.2pt] (0.000,-0.000) rectangle (12.121,-14.807);

%--------------------------------------------------------------
% Riparo mobile con posizioni aperto/chiuso e direzione di apertura
%--------------------------------------------------------------
\draw[draw={rgb,255:red,77;green,77;blue,77}, line width=0.8pt, fill={rgb,255:red,216;green,216;blue,216}]
  (2.207,-3.500) -- (4.171,-3.500) -- (5.029,-4.300) -- (5.029,-8.200) -- (4.171,-8.964)
  -- (2.207,-8.964) -- cycle;
\draw[filo] (0.707,-2.329) -- (3.050,-2.329);
\draw[filo] (0.707,-8.964) -- (2.207,-8.964);
\node[etichetta] at (1.207,-1.614) {Aperto};
\node[etichetta, anchor=east] at (2.064,-8.357) {Chiuso};
\draw[filo, -{Latex[length=0.464cm,width=0.286cm]}] (1.264,-7.414) -- (1.264,-3.786);

%--------------------------------------------------------------
% B1: interruttore di posizione con contatto di apertura ad
% azionamento meccanico positivo, azionato dal riparo tramite rotella
%--------------------------------------------------------------
\draw[filo] (4.800,-3.800) circle (0.2);
\draw[filo] (5.000,-3.800) -- (5.274,-3.800);
\draw[meccanico] (5.479,-3.800) -- (7.200,-3.800);
% simbolo di apertura positiva
\draw[filo] (6.297,-3.239) circle (0.421);
\draw[filo] (5.946,-3.239) -- (6.636,-3.239);
\draw[filo] (6.297,-2.990) -- (6.636,-3.239) -- (6.297,-3.500);
% alimentazione (+) e contatto
\draw[filo] (6.893,-1.190) -- (7.183,-1.190);
\draw[filo] (7.040,-1.047) -- (7.040,-1.343);
\fill[black] (6.907,-1.396) -- (7.183,-1.396) -- (7.040,-2.000) -- cycle;
\draw[filo] (7.040,-2.000) -- (7.040,-3.490) -- (7.446,-3.490);
\draw[filo] (7.040,-4.276) -- (7.350,-3.357);
\draw[filo] (7.040,-4.276) -- (7.040,-11.107);
\node[etichetta] at (6.226,-4.453) {B1};

%--------------------------------------------------------------
% Q1: bobina del contattore, linea di ritorno e messa a terra
%--------------------------------------------------------------
\draw[potenza] (6.064,-11.107) rectangle (7.964,-12.014);
\draw[filo] (7.036,-12.014) -- (7.036,-12.979);
\node[etichetta] at (5.300,-11.393) {Q1};
\draw[filo] (4.907,-12.979) -- (8.693,-12.979);
\node[nodo] at (7.036,-12.979) {};
\node[nodo] at (6.207,-12.979) {};
\draw[filo] (6.207,-12.979) -- (6.207,-13.879);
\draw[filo] (5.571,-13.879) -- (6.836,-13.879);
\draw[filo] (5.743,-14.207) -- (6.679,-14.207);
\draw[filo] (5.886,-14.536) -- (6.536,-14.536);

%--------------------------------------------------------------
% Circuito di potenza: linea trifase L, contatto principale di Q1
% e motore M
%--------------------------------------------------------------
\draw[potenza] (10.300,-1.229) -- (10.300,-5.857);
\draw[potenza] (9.764,-2.750) -- (10.821,-2.214);
\draw[potenza] (9.764,-3.057) -- (10.821,-2.521);
\draw[potenza] (9.764,-3.357) -- (10.821,-2.821);
\node[etichetta, text=verdepotenza] at (10.279,-0.614) {L};
% contatto principale con apertura automatica
\draw[potenza] (10.300,-5.586) arc[start angle=90, end angle=270, radius=0.136];
\draw[potenza] (9.943,-5.771) -- (10.300,-6.414);
\draw[potenza] (10.300,-6.414) -- (10.300,-9.821);
\node[etichetta] at (9.321,-6.143) {Q1};
% motore trifase
\draw[potenza] (10.389,-10.839) circle (1.018);
\node[inner sep=1pt, scale=1.8, text=verdepotenza, font=\sffamily] at (10.321,-10.393) {M};
\node[inner sep=1pt, scale=1.8, text=verdepotenza, font=\sffamily] at (10.300,-11.321) {3$\sim$};

\end{tikzpicture}
}
        \captionof{figure}{Controllo della posizione del riparo mobile per la prevenzione di movimenti pericolosi (STO -- Safe Torque Off).}
        \label{fig:08-11}
    \end{minipage}
\end{center}

\textbf{Funzione di sicurezza}
\begin{itemize}
    \item funzione di arresto attinente alla sicurezza, avviata da un riparo: l'apertura del riparo mobile attiva la funzione di sicurezza STO (safe torque off).
\end{itemize}

\textbf{Descrizione funzionale}
\begin{itemize}
    \item l'apertura del riparo mobile (ad es. riparo di sicurezza) è rilevata da un interruttore di posizione B1 con contatto ad apertura positiva, che comanda un contattore Q1. La diseccitazione di Q1 interrompe o impedisce movimenti o stati pericolosi.
    \item la funzione di sicurezza non può essere mantenuta a fronte di tutti i guasti dei componenti e dipende dall'affidabilità dei componenti stessi.
    \item non sono attuate misure per il rilevamento dei guasti.
    \item la rimozione del dispositivo di protezione non viene rilevata.
\end{itemize}

\textbf{Caratteristiche progettuali}
\begin{itemize}
    \item sono osservati i principi di sicurezza di base e i principi di sicurezza ben collaudati e sono soddisfatti i requisiti della Categoria B. Sono attuati i circuiti di protezione (quali la protezione dei contatti) descritti nei paragrafi iniziali del Capitolo 8. Il principio della corrente di riposo è impiegato come principio di sicurezza di base. La messa a terra del circuito di comando è considerata un principio di sicurezza ben collaudato.
    \item l'interruttore B1 è un interruttore di posizione con contatto ad apertura positiva conforme alla IEC 60947-5-1, Allegato K, ed è pertanto considerato un componente ben collaudato. Il contatto di apertura interrompe il circuito direttamente per via meccanica quando il dispositivo di protezione non si trova nella posizione di sicurezza.
    \item il contattore Q1 è un componente ben collaudato a condizione che siano soddisfatte le condizioni aggiuntive di cui alla Tabella D.3 della EN ISO 13849-2.
    \item per il controllo della posizione è impiegato un interruttore di posizione. È assicurata una disposizione stabile del dispositivo di protezione ai fini dell'azionamento dell'interruttore di posizione. Gli elementi di azionamento dell'interruttore di posizione sono protetti contro lo spostamento. Sono impiegate esclusivamente parti meccaniche rigide (nessun elemento elastico che agisca nella direzione della forza di azionamento)
    \item la corsa di azionamento dell'interruttore di posizione è conforme alle specifiche del costruttore.
\end{itemize}

\textbf{Calcolo della probabilità di guasto}
\begin{itemize}
    \item MTTF\textsubscript{D}: per B1 è dichiarato un MTTF\textsubscript{D} di 20 10\textsuperscript{6} cicli di manovra [M]. Con 365 giorni lavorativi, 16 ore lavorative al giorno e un tempo di ciclo di 10 minuti, il nop di questi componenti è pari a 35.040 cicli all'anno e l'MTTF\textsubscript{D} è di 5.707 anni. Per il contattore Q1, il valore B10 corrisponde, con carico induttivo (AC-3), a una durata elettrica di 1.300.000 cicli di manovra [M]. Ipotizzando che il 50\% dei guasti sia pericoloso, il valore B\textsubscript{10D} si ottiene raddoppiando il valore B\textsubscript{10}. Il valore di n\textsubscript{op} sopra ipotizzato porta a un MTTF\textsubscript{D} di 742 anni per Q1. La combinazione di B1 e Q1 dà un MTTF\textsubscript{D} di 656 anni per ciascun canale. Tale valore è limitato a 100 anni ("alto").
    Il DC\textsubscript{avg} e le misure contro i guasti per causa comune non sono rilevanti nella Categoria 1.
    \item il sistema di comando elettromeccanico soddisfa la Categoria 1 con un MTTF\textsubscript{D} alto (100 anni). Ne risulta una probabilità media di guasto pericoloso di 1,1 · 10\textsubscript{-6} all'ora. Ciò soddisfa il PL c. Il PL\textsubscript{r} pari a b è pertanto superato.
\end{itemize}
\clearpage

\subsection{Esempio 6}
\label{sec:08-2-6}
\textbf{Esempio 6}
\textit{Dispositivo di avvio/arresto con dispositivo di arresto di emergenza - Categoria 1 - PL c}

% La minipage tiene insieme disegno e didascalia (figura non fluttuante)
\begin{center}
    \begin{minipage}{\textwidth}
        \centering
        \captionsetup{hypcap=false}
        \resizebox{0.6\textwidth}{!}{%==============================================================
%  Figura 8.13 - Dispositivo combinato di avvio/arresto con
%  dispositivo di arresto di emergenza (Esempio 6)
%  Adattamento in italiano della Figura 8.13 dell'IFA Report 2/2017
%
%  I simboli riproducono quelli dell'originale: le coordinate sono
%  ricavate dal disegno di partenza (1 cm = 70 px, origine
%  nell'angolo in alto a sinistra della cornice). Nessun cambio di
%  corpo del font: le etichette sono ingrandite con "scale" (vedere
%  la nota nella Figura 8.3).
%==============================================================
\definecolor{verdepotenza}{RGB}{22,67,38}
\begin{tikzpicture}[
  font=\normalsize,
  filo/.style={draw=black, line width=0.6pt},
  meccanico/.style={draw=black, line width=0.6pt, dash pattern=on 0.264cm off 0.136cm},
  meccanicocorto/.style={draw=black, line width=0.6pt, dash pattern=on 0.143cm off 0.114cm},
  etichetta/.style={inner sep=1pt, scale=1.15},
  nodo/.style={circle, fill=black, inner sep=0pt, minimum size=0.271cm}
]

%--------------------------------------------------------------
% Cornice
%--------------------------------------------------------------
\draw[draw=black, line width=1.2pt] (0.000,-0.000) rectangle (11.879,-13.879);

%--------------------------------------------------------------
% Linea di alimentazione superiore
%--------------------------------------------------------------
\draw[filo] (2.864,-0.471) -- (8.350,-0.471);
\draw[filo] (4.929,-0.471) -- (4.929,-2.757);
\node[nodo] at (4.929,-0.471) {};

%--------------------------------------------------------------
% S1: dispositivo di arresto di emergenza (pulsante a fungo con
% ritenuta, contatto di apertura ad azionamento meccanico positivo)
%--------------------------------------------------------------
% simbolo di apertura positiva
\draw[filo] (4.140,-2.311) circle (0.406);
\draw[filo] (3.826,-2.311) -- (4.479,-2.311);
\draw[filo] (4.140,-2.071) -- (4.479,-2.311) -- (4.140,-2.547);
% pulsante a fungo
\draw[filo] (3.464,-2.800) -- (3.464,-3.514);
\draw[filo] (3.464,-2.800) arc[start angle=90, end angle=270, x radius=0.271, y radius=0.357];
% collegamento meccanico con ritenuta
\draw[meccanicocorto] (3.464,-3.161) -- (4.121,-3.161);
\draw[filo] (4.121,-3.161) -- (4.300,-3.514) -- (4.493,-3.161);
\draw[meccanicocorto] (4.493,-3.161) -- (5.079,-3.161);
% contatto di apertura
\draw[filo] (4.929,-2.757) -- (5.336,-2.757);
\draw[filo] (4.929,-3.557) -- (5.236,-2.636);
\draw[filo] (4.929,-3.557) -- (4.929,-5.143);
\node[etichetta] at (2.450,-3.086) {S1};
\node[etichetta, anchor=base west] at (0.736,-3.814) {Arresto di};
\node[etichetta, anchor=base west] at (0.736,-4.357) {emergenza};

%--------------------------------------------------------------
% S2: pulsante di START (contatto di chiusura) in parallelo al
% contatto di autoritenuta di Q1
%--------------------------------------------------------------
\node[nodo] at (4.929,-4.357) {};
\draw[filo] (4.929,-4.357) -- (8.100,-4.357) -- (8.100,-5.143);
% pulsante
\draw[filo] (3.311,-5.271) -- (3.111,-5.271) -- (3.111,-5.671) -- (3.311,-5.671);
\draw[filo] (3.111,-5.471) -- (3.193,-5.471);
\draw[meccanico] (3.336,-5.471) -- (4.650,-5.471);
% contatto di chiusura
\draw[filo] (4.929,-5.929) -- (4.360,-5.014);
\draw[filo] (4.929,-5.929) -- (4.929,-7.536);
\node[etichetta] at (2.479,-5.357) {S2};
\node[etichetta, anchor=west] at (0.736,-5.900) {START};
% contatto di autoritenuta Q1
\draw[filo] (8.100,-5.929) -- (7.550,-5.014);
\draw[filo] (8.100,-5.929) -- (8.100,-6.729) -- (4.929,-6.729);
\node[nodo] at (4.929,-6.729) {};
\node[etichetta] at (6.750,-5.400) {Q1};

%--------------------------------------------------------------
% S3: pulsante di STOP (contatto di apertura)
%--------------------------------------------------------------
\draw[filo] (3.319,-7.661) -- (3.111,-7.661) -- (3.111,-8.061) -- (3.319,-8.061);
\draw[filo] (3.111,-7.861) -- (3.193,-7.861);
\draw[meccanico] (3.336,-7.861) -- (5.086,-7.861);
\draw[filo] (4.929,-7.536) -- (5.336,-7.536);
\draw[filo] (4.929,-8.314) -- (5.236,-7.400);
\draw[filo] (4.929,-8.314) -- (4.929,-9.929);
\node[etichetta] at (2.479,-7.757) {S3};
\node[etichetta, anchor=west] at (0.736,-8.286) {STOP};

%--------------------------------------------------------------
% Q1: bobina del contattore, linea di ritorno e messa a terra
%--------------------------------------------------------------
\draw[filo] (4.129,-9.929) rectangle (5.721,-10.714);
\draw[filo] (4.929,-10.714) -- (4.929,-12.300);
\node[etichetta] at (3.336,-10.214) {Q1};
\draw[filo] (1.336,-12.300) -- (6.507,-12.300);
\node[nodo] at (4.929,-12.300) {};
\node[nodo] at (2.657,-12.300) {};
\draw[filo] (2.657,-12.300) -- (2.657,-12.900);
\draw[filo] (2.264,-12.900) -- (3.050,-12.900);
\draw[filo] (2.350,-13.107) -- (2.964,-13.107);
\draw[filo] (2.450,-13.314) -- (2.864,-13.314);

%--------------------------------------------------------------
% Circuito di potenza: linea trifase L, contatto principale di Q1
% e motore M
%--------------------------------------------------------------
\draw[filo] (10.483,-2.371) -- (10.483,-5.633);
\draw[filo] (10.083,-3.510) -- (10.883,-3.119);
\draw[filo] (10.083,-3.734) -- (10.883,-3.343);
\draw[filo] (10.083,-3.957) -- (10.883,-3.567);
\node[etichetta] at (10.450,-1.814) {L};
% contatto principale con apertura automatica
\draw[filo] (10.483,-5.253) arc[start angle=90, end angle=270, radius=0.190];
\draw[filo] (9.993,-5.524) -- (10.483,-6.343);
\draw[filo] (10.483,-6.343) -- (10.483,-8.714);
\node[etichetta] at (9.136,-5.714) {Q1};
% motore trifase
\draw[filo] (10.486,-9.714) circle (1.000);
\node[inner sep=1pt, scale=1.8, text=verdepotenza, font=\sffamily] at (10.479,-9.271) {M};
\node[inner sep=1pt, scale=1.8, text=verdepotenza, font=\sffamily] at (10.479,-10.129) {3$\sim$};

\end{tikzpicture}
}
        \captionof{figure}{Dispositivo combinato di avvio/arresto con dispositivo di arresto di emergenza.}
        \label{fig:08-13}
    \end{minipage}
\end{center}

\textbf{Funzione di sicurezza}
\begin{itemize}
    \item funzione di arresto di emergenza, STO – coppia disinserita in sicurezza mediante azionamento del dispositivo di arresto di emergenza
\end{itemize}

\textbf{Descrizione funzionale}
\begin{itemize}
    \item i movimenti o gli stati pericolosi vengono disattivati mediante l'interruzione della tensione di comando del contattore Q1 all'azionamento del dispositivo di arresto di emergenza S1.
    \item la funzione di sicurezza non può essere mantenuta in presenza di tutti i possibili guasti dei componenti e dipende dall'affidabilità dei componenti stessi.
    \item non sono implementate misure per la rilevazione dei guasti.
\end{itemize}

\textbf[Caratteristiche di progetto]
\begin{itemize}
    \item sono rispettati i principi di sicurezza di base e i principi di sicurezza ben collaudati e sono soddisfatti i requisiti della Categoria B. Sono realizzati i circuiti di protezione (per es. la protezione dei contatti) descritti nei paragrafi iniziali del Capitolo 8. Il principio della corrente di riposo è impiegato come principio di sicurezza di base. Il circuito di comando è inoltre collegato a terra, in quanto principio di sicurezza ben collaudato.
    \item il dispositivo di arresto di emergenza S1 è un interruttore con modo di azionamento positivo conforme alla IEC 60947-5-5 ed è pertanto un componente ben collaudato secondo la Tabella D.3 della EN ISO 13849-2.
    \item il segnale è elaborato da un contattore (arresto di Categoria 0 secondo la IEC 60204-1).
    \item il contattore Q1 è un componente ben collaudato a condizione che siano rispettate le condizioni supplementari indicate nella Tabella D.3 della EN ISO 13849-2.
\end{itemize}

\textbf{Note}
\begin{itemize}
    \item la funzione di arresto in caso di emergenza è una misura di protezione che integra le funzioni di sicurezza per la protezione delle zone pericolose.
\end{itemize}

\textbf{Calcolo della probabilità di guasto}
\begin{itemize}
    \item MTTF\textsubscript{D}: S1 è un dispositivo di arresto di emergenza standard secondo la EN ISO 13850. È costruito conformemente alla IEC 60947-5-5. Secondo la EN ISO 13849-1, Tabella C.1, per i dispositivi di arresto di emergenza in questo caso può essere assunto un valore B10D di 100.000 cicli di commutazione, indipendentemente dal carico [S]. Per il contattore Q1, il valore B10 corrisponde, con carico induttivo (AC 3), a una durata elettrica di 1.300.000 cicli di commutazione [M]. Nell'ipotesi che il 50\% dei guasti sia pericoloso, il valore B10D si ottiene raddoppiando il valore B\textsubscript{10}. Se si ipotizza che il dispositivo di avvio/arresto venga azionato due volte al giorno per 365 giorni lavorativi e che il dispositivo di arresto di emergenza venga azionato dodici volte all'anno, con un nop risultante di 742 cicli all'anno Q1 presenta un MTTF\textsubscript{D} di 35.040 anni. Questo è anche l'MTTF\textsubscript{D} del canale, che viene limitato a 100 anni ("alto").
    \item il DC\textsubscript{avg} e le misure contro i guasti per causa comune non sono rilevanti nella Categoria 1.
    \item il sistema di comando elettromeccanico soddisfa la Categoria 1 con un MTTF\textsubscript{D} elevato (100 anni). Ne risulta una probabilità media di guasto pericoloso di 1,1 10\textsubscript{-6} all'ora. Ciò soddisfa il PL c.
\end{itemize}
\clearpage

\subsection{Esempio 7}
\label{sec:08-2-7}
\textbf{Esempio 7}
\textit{Sganciatore di minima tensione azionato mediante dispositivo di arresto di emergenza - Categoria 1 - PL c}

% La minipage tiene insieme disegno e didascalia (figura non fluttuante)
\begin{center}
    \begin{minipage}{\textwidth}
        \centering
        \captionsetup{hypcap=false}
        \resizebox{0.6\textwidth}{!}{%==============================================================
%  Figura 8.15 - Dispositivo di arresto di emergenza che agisce
%  sullo sganciatore di minima tensione del dispositivo di
%  sezionamento dell'alimentazione (avviatore motore) (Esempio 7)
%  Adattamento in italiano della Figura 8.15 dell'IFA Report 2/2017
%
%  I simboli riproducono quelli dell'originale: le coordinate sono
%  ricavate dal disegno di partenza (1 cm = 70 px, origine
%  nell'angolo in alto a sinistra della cornice). Nessun cambio di
%  corpo del font: le etichette sono ingrandite con "scale" (vedere
%  la nota nella Figura 8.3). Nell'originale tutti i collegamenti
%  sono in verde scuro e i testi in nero.
%==============================================================
\definecolor{verdepotenza}{RGB}{22,67,38}
\begin{tikzpicture}[
  font=\normalsize,
  filo/.style={draw=verdepotenza, line width=0.6pt},
  meccanico/.style={draw=verdepotenza, line width=0.6pt, dash pattern=on 0.286cm off 0.223cm},
  etichetta/.style={inner sep=1pt, scale=1.15},
  nodo/.style={circle, fill=verdepotenza, inner sep=0pt, minimum size=0.129cm}
]

%--------------------------------------------------------------
% Cornice
%--------------------------------------------------------------
\draw[draw=black, line width=1.2pt] (0.000,-0.000) rectangle (11.850,-11.086);

%--------------------------------------------------------------
% Linea di alimentazione superiore
%--------------------------------------------------------------
\draw[filo] (2.179,-0.371) -- (7.207,-0.371);
\draw[filo] (4.093,-0.371) -- (4.093,-2.390);
\node[nodo] at (4.093,-0.371) {};

%--------------------------------------------------------------
% S1: dispositivo di arresto di emergenza (pulsante a fungo con
% ritenuta, contatto di apertura ad azionamento meccanico positivo)
%--------------------------------------------------------------
% simbolo di apertura positiva
\draw[filo] (3.254,-1.836) circle (0.411);
\draw[filo] (3.011,-1.836) -- (3.521,-1.836);
\draw[filo] (3.254,-1.629) -- (3.521,-1.836) -- (3.254,-2.040);
% pulsante a fungo
\draw[filo] (2.729,-2.421) -- (2.729,-3.076);
\draw[filo] (2.729,-2.421) arc[start angle=90, end angle=270, x radius=0.250, y radius=0.327];
% collegamento meccanico con ritenuta
\draw[filo] (2.729,-2.736) -- (2.933,-2.736);
\draw[filo] (3.136,-2.736) -- (3.314,-2.736) -- (3.476,-3.040) -- (3.676,-2.736) -- (3.819,-2.736);
\draw[filo] (4.007,-2.736) -- (4.200,-2.736);
% contatto di apertura
\draw[filo] (4.093,-2.390) -- (4.433,-2.390);
\draw[filo] (4.064,-3.129) -- (4.350,-2.271);
\node[etichetta] at (1.793,-2.757) {S1};
\node[etichetta, anchor=base west] at (0.493,-3.729) {Arresto di};
\node[etichetta, anchor=base west] at (0.493,-4.286) {emergenza};

%--------------------------------------------------------------
% Linea di ritorno verso lo sganciatore di minima tensione
%--------------------------------------------------------------
\draw[filo] (4.064,-3.129) -- (4.064,-10.407) -- (6.693,-10.407) -- (6.693,-7.057);

%--------------------------------------------------------------
% Q1: avviatore motore con sganciatore di minima tensione
%--------------------------------------------------------------
\draw[filo] (6.107,-6.386) rectangle (7.293,-7.057);
\node[inner sep=1pt, scale=1.3, font=\sffamily] at (6.693,-6.714) {U\,\textless};
% collegamento meccanico sganciatore - serratura
\draw[filo] (6.693,-6.386) -- (6.693,-6.100);
\draw[filo] (6.693,-5.943) -- (6.693,-5.657);
\draw[filo] (6.693,-5.443) -- (6.693,-4.600);
% serratura
\draw[filo] (6.407,-4.600) rectangle (6.979,-5.157);
\draw[filo] (6.321,-4.871) -- (7.450,-4.871);
% ritenuta a sinistra della serratura
\draw[filo] (5.736,-4.757) -- (5.736,-4.986);
\draw[filo] (5.736,-4.871) -- (5.821,-4.871);
\draw[filo] (5.907,-4.871) -- (5.993,-4.871) -- (6.071,-5.014) -- (6.150,-4.871) -- (6.236,-4.871);
% collegamento meccanico verso il contatto principale
\draw[meccanico] (7.664,-4.871) -- (10.493,-4.871);
\node[etichetta] at (5.893,-5.814) {Q1};
\node[etichetta] at (9.021,-3.786) {Avviatore};
\node[etichetta] at (9.021,-4.343) {motore};

%--------------------------------------------------------------
% Circuito di potenza: linea trifase L, contatto principale di Q1
% e motore M
%--------------------------------------------------------------
\draw[filo] (10.679,-1.671) -- (10.679,-4.771);
\draw[filo] (10.321,-2.700) -- (11.036,-2.343);
\draw[filo] (10.321,-2.907) -- (11.036,-2.550);
\draw[filo] (10.321,-3.114) -- (11.036,-2.757);
\node[etichetta] at (10.679,-1.243) {L};
% contatto principale con apertura automatica
\draw[filo] (10.679,-4.843) circle (0.071);
\draw[filo] (10.436,-4.743) -- (10.686,-5.229);
\draw[filo] (10.686,-5.229) -- (10.686,-7.536);
% motore trifase
\draw[filo] (10.686,-8.179) circle (0.643);
\node[inner sep=1pt, scale=1.3, font=\sffamily] at (10.679,-7.857) {M};
\node[inner sep=1pt, scale=1.3, font=\sffamily] at (10.679,-8.471) {3$\sim$};

\end{tikzpicture}
}
        \captionof{figure}{Dispositivo di arresto di emergenza che agisce sullo sganciatore di minima tensione del dispositivo di sezionamento dell'alimentazione (avviatore motore).}
        \label{fig:08-15}
    \end{minipage}
\end{center}

\textbf{Funzioni di sicurezza}
\begin{itemize}
    \item funzione di arresto di emergenza, STO (coppia disinserita in sicurezza), ottenuta mediante azionamento del dispositivo di arresto di emergenza che agisce sullo sganciatore di minima tensione di un avviatore motore o, se del caso, del dispositivo di sezionamento dell'alimentazione.
\end{itemize}

\textbf{Descrizione funzionale}
\begin{itemize}
    \item i movimenti o gli stati pericolosi vengono interrotti, all'azionamento del dispositivo di arresto di emergenza S1, mediante lo sganciatore di minima tensione del dispositivo di sezionamento dell'alimentazione, in questo caso costituito da un avviatore motore Q1.
    \item la funzione di sicurezza non può essere mantenuta in presenza di tutti i possibili guasti dei componenti e dipende dall'affidabilità dei componenti stessi.
    \item non sono implementate misure per la rilevazione dei guasti.
\end{itemize}

\textbf{Caratteristiche di progetto}
\begin{itemize}
    \item sono rispettati i principi di sicurezza di base e i principi di sicurezza ben collaudati e sono soddisfatti i requisiti della Categoria B. Sono realizzati i circuiti di protezione (per es. la protezione dei contatti) descritti nei paragrafi iniziali del Capitolo 8. Il principio della corrente di riposo dello sganciatore di minima tensione è impiegato come principio di sicurezza di base.
    \item il dispositivo di arresto di emergenza S1 è un interruttore con modo di azionamento positivo conforme alla IEC 60947-5-5 ed è pertanto un componente ben collaudato secondo la Tabella D.3 della EN ISO 13849-2.
    \item l'avviatore motore Q1 deve essere considerato equivalente a un interruttore automatico secondo la Tabella D.3 della EN ISO 13849-2. Q1 può pertanto essere considerato un componente ben collaudato.
    \item l'alimentazione dell'intera macchina viene disinserita (arresto di Categoria 0 secondo la IEC 60204-1).
\end{itemize}

\textbf{Note}
La funzione di arresto in caso di emergenza è una misura di protezione che integra le funzioni di sicurezza per la protezione delle zone pericolose.

\textbf{Calcolo della probabilità di guasto}
\begin{itemize}
    \item MTTF\textsubscript{D}: S1 è un dispositivo di arresto di emergenza standard secondo la EN ISO 13850. È costruito conformemente alla IEC 60947-5-5. Secondo la EN ISO 13849-1, Tabella C.1, per i dispositivi di arresto di emergenza in questo caso può essere assunto un valore B\textsubscript{10D}D di 100.000 cicli di commutazione, indipendentemente dal carico [S]. Per lo sganciatore di minima tensione dell'avviatore motore Q1, il valore B\textsubscript{10} corrisponde approssimativamente alla durata elettrica di 10.000 cicli di commutazione [M]. Nell'ipotesi che il 50\% dei guasti sia pericoloso, il valore B\textsubscript{10D} si ottiene raddoppiando il valore B\textsubscript{10}. Con un azionamento del dispositivo di arresto di emergenza dodici volte all'anno e un nop risultante di 12 cicli all'anno, Q1 presenta un MTTF\textsubscript{D} di 16.666 anni. Questo è anche l'MTTF\textsubscript{D} del canale, che viene limitato a 100 anni ("alto").
    \item il DC\textsubscript{avg} e le misure contro i guasti per causa comune non sono rilevanti nella Categoria 1.
    \item il sistema di comando elettromeccanico soddisfa la Categoria 1 con un MTTF\textsubscript{D} elevato (100 anni). Ne risulta una probabilità media di guasto pericoloso di 1,1 10\textsubscript{-6} all'ora. Ciò soddisfa il PL c.
\end{itemize}
\clearpage

\subsection{Esempio 8}
\label{sec:08-2-8}
\textbf{Esempio 8}
Questo esempio è stato eliminato, in quanto la tecnologia non è più attuale.
\clearpage

\subsection{Esempio 9}
\label{sec:08-2-9}
\textbf{Esempio 9}
\textit{Barriere fotoelettriche testate - Categoria 2 - PL c con dispositivo di commutazione del segnale di uscita di Categoria 1 posto a valle}

% La minipage tiene insieme disegno e didascalia (figura non fluttuante)
\begin{center}
    \begin{minipage}{\textwidth}
        \centering
        \captionsetup{hypcap=false}
        \resizebox{\textwidth}{!}{%==============================================================
%  Figura 8.17 - Test delle barriere fotoelettriche con un PLC
%  standard (Esempio 9)
%  Adattamento in italiano della Figura 8.17 dell'IFA Report 2/2017
%
%  I simboli riproducono quelli dell'originale: le coordinate sono
%  ricavate dal disegno di partenza (1 cm = 70 px, origine
%  nell'angolo in alto a sinistra della cornice). Nessun cambio di
%  corpo del font: le etichette sono ingrandite con "scale" (vedere
%  la nota nella Figura 8.3). Le etichette START e STOP restano in
%  inglese, come nella Figura 8.13.
%==============================================================
\begin{tikzpicture}[
  font=\normalsize,
  filo/.style={draw=black, line width=0.6pt},
  sottile/.style={draw=black, line width=0.4pt},
  meccanico/.style={draw=black, line width=0.6pt, dash pattern=on 0.200cm off 0.079cm},
  meccanicocorto/.style={draw=black, line width=0.6pt, dash pattern=on 0.100cm off 0.063cm},
  etichetta/.style={inner sep=0pt, scale=0.95},
  medio/.style={inner sep=0pt, scale=1.2},
  grande/.style={inner sep=0pt, scale=1.75},
  motore/.style={inner sep=0pt, scale=1.6, font=\sffamily},
  nodo/.style={circle, fill=black, inner sep=0pt, minimum size=0.186cm}
]

%--------------------------------------------------------------
% Cornice
%--------------------------------------------------------------
\draw[draw=black, line width=1.2pt] (0.000,-0.000) rectangle (19.636,-15.079);

%--------------------------------------------------------------
% Linea di alimentazione +U_B
%--------------------------------------------------------------
\draw[filo] (2.421,-1.000) -- (18.236,-1.000);
\node[nodo] at (3.407,-1.000) {};
\node[nodo] at (5.393,-1.000) {};
\node[nodo] at (9.350,-1.000) {};
\node[nodo] at (17.579,-1.000) {};
\node[etichetta, anchor=base west] at (18.679,-1.043) {+U\textsubscript{B}};

%--------------------------------------------------------------
% S2: pulsante di START (contatto di chiusura) -> I1.0
%--------------------------------------------------------------
\draw[filo] (3.407,-1.000) -- (3.407,-2.486);
\draw[filo] (3.064,-2.414) -- (3.407,-2.971);
\draw[filo] (2.807,-2.614) -- (2.693,-2.614) -- (2.693,-2.857) -- (2.807,-2.857);
\draw[filo] (2.693,-2.743) -- (2.779,-2.743);
\draw[meccanicocorto, dash pattern=on 0.100cm off 0.064cm] (2.850,-2.743) -- (3.264,-2.743);
\draw[filo] (3.407,-2.971) -- (3.407,-4.214);
\fill (3.321,-4.214) -- (3.493,-4.214) -- (3.407,-4.471) -- cycle;
\node[nodo] at (3.407,-3.486) {};
\node[etichetta, anchor=base] at (2.021,-2.829) {S2};
\node[etichetta, anchor=base] at (2.021,-3.143) {START};

%--------------------------------------------------------------
% K2: contatto di chiusura in serie a R1 (ritorno su I1.0)
% e diodo R1
%--------------------------------------------------------------
\draw[filo] (3.407,-3.486) -- (1.436,-3.486) -- (1.436,-7.929);
\draw[filo] (1.436,-7.929) -- (1.679,-7.929);
\draw[filo] (1.621,-7.857) -- (1.436,-8.400);
\draw[filo] (1.436,-8.400) -- (1.436,-10.400) -- (3.407,-10.400);
\draw[filo] (1.186,-9.171) -- (1.686,-9.171) -- (1.436,-9.657) -- cycle;
\draw[filo] (1.186,-9.657) -- (1.686,-9.657);
\node[nodo] at (3.407,-10.400) {};
\node[etichetta, anchor=base east] at (1.064,-8.200) {K2};
\node[etichetta, anchor=base east] at (0.750,-9.443) {R1};

%--------------------------------------------------------------
% Q1: contatto ausiliario di apertura e K2: contatto di
% apertura -> I1.1
%--------------------------------------------------------------
\draw[filo] (5.393,-1.000) -- (5.393,-2.000);
\draw[filo] (5.393,-2.000) -- (5.636,-2.000);
\draw[filo] (5.621,-1.871) -- (5.393,-2.471);
\fill (5.636,-1.800) circle (0.071);
\draw[filo] (5.393,-2.471) -- (5.393,-3.486);
\draw[filo] (5.393,-3.486) -- (5.636,-3.486);
\draw[filo] (5.579,-3.414) -- (5.393,-3.957);
\draw[filo] (5.393,-3.957) -- (5.393,-4.214);
\fill (5.307,-4.214) -- (5.479,-4.214) -- (5.393,-4.471) -- cycle;
\node[etichetta, anchor=base east] at (5.064,-2.257) {Q1};
\node[etichetta, anchor=base east] at (5.093,-3.800) {K2};

%--------------------------------------------------------------
% K2: contatto di chiusura -> I1.2 (collegamento a R2)
%--------------------------------------------------------------
\draw[filo] (13.286,-8.829) -- (13.286,-1.500) -- (7.364,-1.500) -- (7.364,-3.486);
\draw[filo] (7.021,-3.400) -- (7.364,-3.971);
\draw[filo] (7.364,-3.971) -- (7.364,-4.214);
\fill (7.279,-4.214) -- (7.450,-4.214) -- (7.364,-4.471) -- cycle;
\draw[meccanicocorto] (5.564,-3.729) -- (7.221,-3.729);

%--------------------------------------------------------------
% S1: pulsante di STOP (contatto di apertura ad azionamento
% meccanico positivo) -> I1.3
%--------------------------------------------------------------
\draw[filo] (9.350,-1.000) -- (9.350,-2.486);
\draw[filo] (9.350,-2.486) -- (9.593,-2.486);
\draw[filo] (9.607,-2.414) -- (9.350,-2.957);
\draw[sottile] (9.021,-2.286) circle (0.179);
\draw[sottile] (8.879,-2.286) -- (9.193,-2.286);
\draw[sottile] (9.007,-2.171) -- (9.193,-2.286) -- (9.021,-2.393);
\draw[filo] (8.821,-2.543) -- (8.693,-2.543) -- (8.693,-2.786) -- (8.821,-2.786);
\draw[filo] (8.693,-2.657) -- (8.779,-2.657);
\draw[meccanicocorto, dash pattern=on 0.100cm off 0.071cm] (8.850,-2.657) -- (9.464,-2.657);
\draw[filo] (9.350,-2.957) -- (9.350,-4.214);
\fill (9.264,-4.214) -- (9.436,-4.214) -- (9.350,-4.471) -- cycle;
\node[nodo] at (9.350,-3.486) {};
\draw[filo] (9.350,-3.486) -- (11.821,-3.486) -- (11.821,-6.943) -- (10.336,-6.943);
\node[etichetta, anchor=base] at (8.164,-2.786) {S1};
\node[etichetta, anchor=base] at (8.164,-3.143) {STOP};

%--------------------------------------------------------------
% K3: PLC standard
%--------------------------------------------------------------
\draw[filo] (2.421,-4.471) rectangle (10.336,-7.429);
\draw[filo] (2.421,-5.457) -- (10.336,-5.457);
\draw[filo] (2.421,-6.443) -- (10.336,-6.443);
\node[etichetta, anchor=base east] at (2.179,-4.686) {K3};
\node[etichetta, anchor=base] at (3.421,-4.757) {I1.0};
\node[etichetta, anchor=base] at (5.421,-4.757) {I1.1};
\node[etichetta, anchor=base] at (7.393,-4.757) {I1.2};
\node[etichetta, anchor=base] at (9.386,-4.757) {I1.3};
\node[medio, anchor=base] at (6.236,-5.257) {Ingressi};
\node[grande, anchor=base] at (6.236,-6.100) {PLC};
\node[medio, anchor=base] at (6.164,-6.814) {Uscite};
\node[etichetta, anchor=base] at (3.414,-7.329) {O1.0};
\node[etichetta, anchor=base] at (5.379,-7.329) {O1.1};
\node[etichetta, anchor=base] at (8.350,-7.329) {O1.2};

%--------------------------------------------------------------
% R2 e bobina di K2
%--------------------------------------------------------------
\draw[filo] (13.036,-6.200) -- (13.536,-6.200) -- (13.286,-5.700) -- cycle;
\draw[filo] (13.036,-5.700) -- (13.536,-5.700);
\node[nodo] at (13.286,-7.929) {};
\draw[filo] (13.286,-7.929) -- (19.221,-7.929) -- (19.221,-12.386) -- (18.264,-12.386);
\fill (19.129,-10.529) -- (19.314,-10.529) -- (19.221,-10.271) -- cycle;
\draw[filo] (12.793,-8.829) rectangle (13.779,-9.314);
\draw[filo] (13.286,-9.314) -- (13.286,-9.714);
\draw[filo] (13.043,-9.714) -- (13.529,-9.714);
\node[etichetta, anchor=base east] at (12.664,-6.029) {R2};
\node[etichetta, anchor=base east] at (12.536,-8.886) {K2};

%--------------------------------------------------------------
% O1.0: contatto di chiusura di K1 e bobina di K1
%--------------------------------------------------------------
\draw[filo] (3.407,-7.429) -- (3.407,-8.914);
\draw[filo] (3.064,-8.843) -- (3.407,-9.400);
\draw[filo] (3.407,-9.400) -- (3.407,-11.643);
\draw[filo] (2.921,-11.643) rectangle (3.907,-12.129);
\draw[filo] (3.407,-12.129) -- (3.407,-12.529);
\draw[filo] (3.164,-12.529) -- (3.650,-12.529);
\node[etichetta, anchor=base east] at (2.779,-9.157) {K1};
\node[etichetta, anchor=base east] at (2.607,-11.943) {K1};

%--------------------------------------------------------------
% O1.1: contatto di apertura di K1, contatto di chiusura di K2
% e bobina di Q1
%--------------------------------------------------------------
\draw[filo] (5.393,-7.429) -- (5.393,-8.914);
\draw[filo] (5.393,-8.914) -- (5.636,-8.914);
\draw[filo] (5.579,-8.843) -- (5.393,-9.386);
\draw[filo] (5.393,-9.386) -- (5.393,-10.400);
\draw[filo] (5.050,-10.329) -- (5.393,-10.900);
\draw[filo] (5.393,-10.900) -- (5.393,-11.643);
\draw[filo] (4.893,-11.643) rectangle (5.879,-12.129);
\draw[filo] (5.393,-12.129) -- (5.393,-12.529);
\draw[filo] (5.150,-12.529) -- (5.636,-12.529);
\node[etichetta, anchor=base east] at (4.793,-10.686) {K2};
\node[etichetta, anchor=base east] at (4.621,-11.943) {Q1};

%--------------------------------------------------------------
% O1.2: contatti di chiusura di K1 e K2 in parallelo,
% segnale di test verso le barriere fotoelettriche
%--------------------------------------------------------------
\draw[filo] (8.350,-7.429) -- (8.350,-8.914);
\draw[filo] (8.007,-8.843) -- (8.350,-9.400);
\draw[filo] (8.350,-9.400) -- (8.350,-12.386) -- (9.093,-12.386);
\fill (9.093,-12.300) -- (9.093,-12.471) -- (9.350,-12.386) -- cycle;
\node[nodo] at (8.350,-7.929) {};
\node[nodo] at (8.350,-10.400) {};
\draw[filo] (8.350,-7.929) -- (9.836,-7.929) -- (9.836,-8.914);
\draw[filo] (9.493,-8.843) -- (9.836,-9.400);
\draw[filo] (9.836,-9.400) -- (9.836,-10.400) -- (8.350,-10.400);
\draw[meccanico] (3.407,-9.100) -- (8.179,-9.100);
\draw[meccanicocorto] (9.636,-9.071) -- (12.793,-9.071);

%--------------------------------------------------------------
% Circuito di potenza: linea trifase L, contatto principale
% di Q1 e motore M
%--------------------------------------------------------------
\draw[filo] (7.050,-9.871) -- (7.050,-11.529);
\draw[filo] (6.800,-10.771) -- (7.300,-10.529);
\draw[filo] (6.800,-10.914) -- (7.300,-10.671);
\draw[filo] (6.800,-11.057) -- (7.300,-10.814);
\draw[filo] (7.050,-11.529) arc[start angle=90, end angle=270, radius=0.107];
\draw[filo] (6.750,-11.686) -- (7.050,-12.171);
\draw[filo] (7.050,-12.171) -- (7.050,-13.371);
\draw[meccanico, dash pattern=on 0.214cm off 0.086cm] (5.879,-11.886) -- (6.879,-11.886);
\draw[filo] (7.036,-14.000) circle (0.629);
\node[etichetta, anchor=base] at (7.057,-9.671) {L};
\node[motore, anchor=base] at (7.050,-13.900) {M};
\node[motore, anchor=base] at (7.050,-14.414) {3$\sim$};

%--------------------------------------------------------------
% Barriere fotoelettriche F1 ... Fn (emettitore e ricevitore)
%--------------------------------------------------------------
\draw[filo] (9.350,-12.129) rectangle (10.707,-13.614);
\draw[filo] (9.721,-12.386) -- (9.721,-13.371);
\draw[filo] (9.464,-12.629) -- (9.964,-12.629) -- (9.721,-13.129) -- cycle;
\draw[filo] (9.464,-13.129) -- (9.964,-13.129);
\draw[filo] (10.136,-12.786) -- (10.407,-12.786);
\fill (10.379,-12.750) -- (10.379,-12.821) -- (10.579,-12.786) -- cycle;
\draw[filo] (10.136,-12.971) -- (10.407,-12.971);
\fill (10.379,-12.936) -- (10.379,-13.007) -- (10.579,-12.971) -- cycle;
\draw[filo] (10.021,-12.129) -- (10.021,-10.900);
\draw[filo] (10.021,-13.614) -- (10.021,-14.014);
\draw[filo] (9.779,-14.014) -- (10.264,-14.014);
\draw[meccanicocorto, dash pattern=on 0.100cm off 0.071cm] (10.707,-12.871) -- (11.264,-12.871);
\fill (11.264,-12.814) -- (11.264,-12.929) -- (11.479,-12.871) -- cycle;
\draw[filo] (11.479,-12.129) rectangle (12.821,-13.614);
\draw[filo] (11.821,-12.629) -- (11.821,-13.129);
\draw[filo] (11.821,-12.743) -- (12.293,-12.386) -- (12.693,-12.386);
\draw[filo] (11.821,-13.014) -- (12.293,-13.371) -- (12.693,-13.371);
\draw[filo] (11.964,-13.221) -- (12.143,-13.257) -- (12.079,-13.071);
\draw[filo] (12.136,-12.129) -- (12.136,-10.900);
\draw[filo] (12.136,-13.614) -- (12.136,-14.014);
\draw[filo] (11.893,-14.014) -- (12.379,-14.014);
\draw[filo] (14.764,-12.129) rectangle (16.136,-13.614);
\draw[filo] (15.136,-12.386) -- (15.136,-13.371);
\draw[filo] (14.879,-12.629) -- (15.379,-12.629) -- (15.136,-13.129) -- cycle;
\draw[filo] (14.879,-13.129) -- (15.379,-13.129);
\draw[filo] (15.550,-12.786) -- (15.821,-12.786);
\fill (15.793,-12.750) -- (15.793,-12.821) -- (15.993,-12.786) -- cycle;
\draw[filo] (15.550,-12.971) -- (15.821,-12.971);
\fill (15.793,-12.936) -- (15.793,-13.007) -- (15.993,-12.971) -- cycle;
\draw[filo] (15.450,-12.129) -- (15.450,-10.900);
\draw[filo] (15.450,-13.614) -- (15.450,-14.014);
\draw[filo] (15.207,-14.014) -- (15.693,-14.014);
\draw[meccanicocorto, dash pattern=on 0.100cm off 0.071cm] (16.136,-12.871) -- (16.679,-12.871);
\fill (16.679,-12.814) -- (16.679,-12.929) -- (16.893,-12.871) -- cycle;
\draw[filo] (16.893,-12.129) rectangle (18.264,-13.614);
\draw[filo] (17.236,-12.629) -- (17.236,-13.129);
\draw[filo] (17.236,-12.743) -- (17.707,-12.386) -- (18.107,-12.386);
\draw[filo] (17.236,-13.014) -- (17.707,-13.371) -- (18.107,-13.371);
\draw[filo] (17.379,-13.221) -- (17.557,-13.257) -- (17.493,-13.071);
\draw[filo] (17.579,-12.129) -- (17.579,-10.900);
\draw[filo] (17.579,-13.614) -- (17.579,-14.014);
\draw[filo] (17.336,-14.014) -- (17.821,-14.014);
\draw[filo] (17.579,-1.000) -- (17.579,-10.900);
\draw[filo] (10.021,-10.900) -- (17.579,-10.900);
\node[nodo] at (12.136,-10.900) {};
\node[nodo] at (15.450,-10.900) {};
\node[nodo] at (17.579,-10.900) {};
\draw[meccanicocorto, dash pattern=on 0.086cm off 0.071cm] (12.821,-12.386) -- (14.521,-12.386);
\fill (14.507,-12.300) -- (14.507,-12.471) -- (14.764,-12.386) -- cycle;
\draw[filo] (18.264,-12.386) -- (19.221,-12.386);
\node[etichetta, anchor=base] at (11.043,-11.843) {F1};
\node[etichetta, anchor=base] at (16.479,-11.843) {Fn};
\node[etichetta, anchor=base] at (11.086,-14.429) {1ª barriera};
\node[etichetta, anchor=base] at (11.086,-14.914) {fotoelettrica};
\node[etichetta, anchor=base] at (16.514,-14.429) {n-esima barriera};
\node[etichetta, anchor=base] at (16.514,-14.914) {fotoelettrica};
\draw[meccanicocorto, dash pattern=on 0.086cm off 0.071cm] (12.593,-14.600) -- (14.879,-14.600);

\end{tikzpicture}
}
        \captionof{figure}{Test delle barriere fotoelettriche con un PLC standard.}
        \label{fig:08-17}
    \end{minipage}
\end{center}

\textbf{funzioni di sicurezza}
\begin{itemize}
    \item funzione di arresto legata alla sicurezza, attivata da un dispositivo di protezione: quando il fascio luminoso viene interrotto, un movimento pericoloso viene arrestato (STO - coppia disinserita in sicurezza).
\end{itemize}

\textbf{Descrizione funzionale}
\begin{itemize}
    \item l'interruzione di un fascio luminoso delle n barriere fotoelettriche F1,..Fn collegate in cascata determina un comando di disinserzione sia per via elettromeccanica, con la diseccitazione del contattore ausiliario K2, sia attraverso l'uscita del PLC (O1.1) del canale di test. Il movimento pericoloso viene quindi arrestato per mezzo del contattore principale Q1.
    \item le barriere fotoelettriche vengono testate prima di ogni avvio del movimento pericoloso, a seguito della pressione del pulsante di avvio S2. A tale scopo l'uscita O1.2 del PLC diseccita l'emettitore della barriera fotoelettrica in risposta a un comando software. La reazione del ricevitore (K2 si diseccita nuovamente) è controllata sugli ingressi I1.1 e I1.2 del PLC. Se il comportamento risulta esente da guasti, K2 si automantiene tramite O1.2 e il movimento pericoloso può essere avviato rilasciando S2. K1 viene diseccitato tramite O1.0 e il contattore principale Q1 viene attivato tramite O1.1.
    \item se il test rileva un guasto in una delle barriere fotoelettriche o in K2, le uscite O1.1 e O1.2 vengono disattivate e al contattore principale Q1 non viene più applicato alcun segnale di attivazione.
    \item in caso di guasto globale del PLC (uscita O1.0 a potenziale basso, uscite O1.1 e O1.2 a potenziale alto), l'interruzione di un fascio luminoso determina la diseccitazione di K2 indipendentemente dal PLC. Per garantire questa indipendenza, le uscite delle barriere fotoelettriche sono disaccoppiate dal PLC mediante il diodo di disaccoppiamento R2. In circostanze sfavorevoli, le barriere fotoelettriche possono essere riattivate da K2 azionando il pulsante di avvio, con conseguente attivazione del contattore principale Q1. In questo caso risulterebbe guasta (soltanto) l'apparecchiatura di test. Il guasto dell'apparecchiatura di test viene rilevato in ragione della probabilità che, in tali circostanze, il processo presenti un funzionamento difettoso.
    \item durante il test, l'attivazione di Q1 da parte di K1 e O1.1 è inibita.
\end{itemize}

\textbf{Caratteristiche di progetto}
\begin{itemize}
    \item sono rispettati i principi di sicurezza di base e i principi di sicurezza ben collaudati e sono soddisfatti i requisiti della Categoria B. Sono realizzati i circuiti di protezione (per es. la protezione dei contatti) descritti nei paragrafi iniziali del Capitolo 8.
    \item vengono impiegate barriere fotoelettriche speciali con caratteristiche ottiche adeguate (angolo di apertura, immunità alla luce estranea, ecc.) secondo la IEC 61496-2.
    \item più barriere fotoelettriche possono essere collegate in cascata e monitorate con soli due ingressi del PLC e un relè o un contattore ausiliario.
    \item i contattori ausiliari K1 e K2 dispongono di elementi di contatto collegati meccanicamente secondo la IEC 60947-5-1, Allegato L. Il contattore principale Q1 dispone di un contatto a specchio secondo la IEC 60947-4-1, Allegato F.
    \item i componenti standard F1..Fn e K3 sono impiegati conformemente alle indicazioni del punto 6.3.10.
    \item il software (SRASW) è programmato secondo i requisiti per il PL b (riduzione dei requisiti nel canale di test grazie alla diversità) e le indicazioni del punto 6.3.10.
    \item il pulsante di avvio S2 deve essere collocato all'esterno della zona pericolosa e in una posizione dalla quale la zona pericolosa sia visibile.
    \item il numero, la disposizione e l'altezza dei fasci luminosi devono essere conformi alla EN ISO 13855 e alla IEC 62046.
    \item se una disposizione per la protezione delle zone pericolose consente di oltrepassare il campo di rilevamento e di sostare al di là di esso, sono necessarie ulteriori misure, come un interblocco di riavvio. A tale scopo può essere utilizzato il pulsante di avvio S2: il PLC K3 confronta la durata della pressione del pulsante con un valore massimo e un valore minimo. Solo se le condizioni sono soddisfatte il comando di avvio è considerato valido.
\end{itemize}

\textbf{Note}
\begin{itemize}
    \item l'esempio è destinato all'impiego in applicazioni con richiesta poco frequente della funzione di sicurezza. Ciò consente di soddisfare il requisito previsto per l'architettura designata della Categoria 2, ossia "test molto più frequente della richiesta della funzione di sicurezza" (cfr. Allegato G).
    \item dopo l'attivazione di un arresto, le barriere fotoelettriche restano disattivate fino al successivo avvio. In questo modo è possibile, per esempio, accedere a una zona pericolosa senza che ciò venga "registrato" dal circuito. Tale comportamento può essere modificato adattando opportunamente il circuito.
\end{itemize}

\textbf{Calcolo della probabilità di guasto}
\begin{itemize}
    \item a titolo di esempio, per il calcolo della probabilità di guasto vengono considerate tre barriere fotoelettriche F1…F3. La protezione di una seconda zona pericolosa costituisce un'ulteriore funzione di sicurezza, il cui calcolo viene eseguito separatamente.
    \item per il calcolo della probabilità di guasto, il sistema complessivo è suddiviso in due sottosistemi: "barriere fotoelettriche" e "contattore principale" (Q1).
\end{itemize}

Per il sottosistema "barriere fotoelettriche":
\begin{itemize}
    \item F1, F2, F3 e K2 costituiscono il canale funzionale della struttura circuitale di Categoria 2; il PLC K3 (compreso il diodo di disaccoppiamento R2) costituisce l'apparecchiatura di test. S2 e K1 hanno la funzione di attivare il test della barriera fotoelettrica e non rientrano nel calcolo della probabilità di guasto.
    \item MTTF\textsubscript{D}: per ciascuna delle barriere da F1 a F3 si assume un MTTF\textsubscript{D} di 100 anni [E]. Il valore B\textsubscript{10D} di K2 è pari a 20.000.000 di cicli [S]. Con 240 giorni lavorativi, 16 ore lavorative giornaliere e un tempo di ciclo di 180 secondi, n\textsubscript{op} risulta pari a 76.800 cicli all'anno. Il test descritto sopra raddoppia tale valore, portandolo a un nop di 153.600 cicli all'anno con un MTTFD di 1.302 anni per K2.
\end{itemize}

Da questi valori risulta un MTTF\textsubscript{D} di 32 anni ("alto") per il canale funzionale. Per K3 si assume un MTTF\textsubscript{D} di 50 anni [E]. Rispetto a tale valore, l'MTTF\textsubscript{D} di 228.311 anni [S] del diodo di disaccoppiamento R2 è irrilevante.

\begin{itemize}
    \item DC\textsubscript{avg}: la giustificazione del DC del 60\% per F1…F3 è il test funzionale descritto sopra. Il DC del 99\% per K2 deriva dal monitoraggio diretto in K3 mediante elementi di contatto collegati meccanicamente. La formula di media restituisce per il DC\textsubscript{avg} un risultato del 61\% ("basso").
    \item misure adeguate contro i guasti per causa comune (85 punti): separazione (15), diversità (20), protezione contro le sovratensioni ecc. (15) e condizioni ambientali (25 + 10)
    \item la combinazione degli elementi di comando del sottosistema "barriere fotoelettriche" soddisfa la Categoria 2 con un MTTFD elevato del canale funzionale (32 anni) e un DCavg basso (61\%). Ne risulta una probabilità media di guasto pericoloso PFHD di 1,9 10\textsuperscript{-6} all'ora.
\end{itemize}

Per il sottosistema "contattore principale" si assumono le ipotesi seguenti:
\begin{itemize}
    \item B\textsubscript{10D} = 1.300.000 cicli [S] con un n\textsubscript{op} di 76.800 cicli all'anno. Ne consegue un MTTF\textsubscript{D} di 169 anni, che secondo la norma viene limitato a 100 anni. La struttura soddisfa la Categoria 1; il DC\textsubscript{avg} e i guasti per causa comune non sono pertanto rilevanti. La probabilità media di guasto pericoloso risultante è di 1,1 10\textsubscript{-6} all'ora.
    \item la somma delle probabilità medie di guasto pericoloso dei due sottosistemi dà un PFHD di 3,0 10⁻⁶ all'ora. Ciò soddisfa il PL c.
    \item se si prevede che la funzione di sicurezza venga richiesta più frequentemente di quanto ipotizzato per l'architettura designata di Categoria 2 (rapporto inferiore a 100:1, cioè più di una volta ogni 5 ore), se ne può tenere conto, conformemente alla Nota 1 dell'Allegato K della norma, applicando una penalizzazione aggiuntiva del 10\% fino a un rapporto di 25:1. Nel caso qui considerato, con tre barriere fotoelettriche, il sottosistema "barriere fotoelettriche" raggiunge comunque un PFHD di 2,1 10\textsuperscript{-6} all'ora. La probabilità media di guasto pericoloso PFH\textsubscript{D} di 3,2 10\textsubscript{-6} all'ora consente però di raggiungere soltanto il PL b. Per raggiungere il PL c occorrerebbe, per esempio, ridurre il numero di barriere fotoelettriche oppure impiegare componenti con un MTTF\textsubscript{D} più elevato.
    \item tenuto conto della stima cautelativa descritta sopra, per la sostituzione programmata del componente soggetto a usura Q1 risulta un tempo di funzionamento (T\textsubscript{10D}) di 17 anni.
\end{itemize}
\clearpage

\textbf{Esempio 10}

% La minipage tiene insieme disegno e didascalia (figura non fluttuante)
\begin{center}
    \begin{minipage}{\textwidth}
        \centering
        \captionsetup{hypcap=false}
        \resizebox{\textwidth}{!}{%==============================================================
%  Figura 8.19 - Arresto di un azionamento con convertitore di
%  frequenza comandato da PLC a seguito di un comando di arresto
%  di emergenza (Esempio 10)
%  Adattamento in italiano della Figura 8.19 dell'IFA Report 2/2017
%
%  I simboli riproducono quelli dell'originale: le coordinate sono
%  ricavate dal disegno di partenza (1 cm = 70 px, origine
%  nell'angolo in alto a sinistra della cornice). Nessun cambio di
%  corpo del font: le etichette sono ingrandite con "scale" (vedere
%  la nota nella Figura 8.3). Le etichette START e STOP restano in
%  inglese, come nella Figura 8.13.
%==============================================================
\begin{tikzpicture}[
  font=\normalsize,
  filo/.style={draw=black, line width=0.6pt},
  sottile/.style={draw=black, line width=0.4pt},
  meccanico/.style={draw=black, line width=0.6pt},
  etichetta/.style={inner sep=0pt, scale=0.9},
  testo/.style={inner sep=0pt, scale=1.0},
  medio/.style={inner sep=0pt, scale=1.65},
  grande/.style={inner sep=0pt, scale=2.5},
  motore/.style={inner sep=0pt, scale=1.4, font=\sffamily},
  nodo/.style={circle, fill=black, inner sep=0pt, minimum size=0.186cm}
]

%--------------------------------------------------------------
% Cornice
%--------------------------------------------------------------
\draw[draw=black, line width=1.2pt] (0.000,-0.000) rectangle (23.757,-17.214);

%--------------------------------------------------------------
% Linea di alimentazione +U_B
%--------------------------------------------------------------
\draw[filo] (3.750,-0.536) -- (20.307,-0.536);
\node[nodo] at (4.764,-0.536) {};
\node[nodo] at (6.793,-0.536) {};
\node[nodo] at (7.821,-0.536) {};
\node[nodo] at (9.836,-0.536) {};
\node[nodo] at (15.143,-0.536) {};
\node[nodo] at (16.114,-0.536) {};
\node[nodo] at (18.064,-0.536) {};
\node[etichetta, anchor=base west] at (20.393,-0.636) {+\,U\textsubscript{B}};

%--------------------------------------------------------------
% S5: pulsante di START (contatto di chiusura) -> I1
%--------------------------------------------------------------
\draw[filo] (4.764,-0.536) -- (4.764,-1.036);
\draw[filo] (4.421,-0.964) -- (4.764,-1.536);
\draw[filo] (4.764,-1.536) -- (4.764,-3.321);
\fill (4.679,-3.321) -- (4.850,-3.321) -- (4.764,-3.564) -- cycle;
\draw[filo] (3.964,-1.121) -- (3.836,-1.121) -- (3.836,-1.379) -- (3.964,-1.379);
\draw[filo] (3.836,-1.250) -- (3.907,-1.250);
\draw[meccanico, dash pattern=on 0.157cm off 0.114cm] (3.979,-1.250) -- (4.593,-1.250);
\node[etichetta, anchor=base] at (3.321,-1.393) {S5};
\node[etichetta, anchor=base] at (3.236,-1.793) {START};

%--------------------------------------------------------------
% K5 -> I0
%--------------------------------------------------------------
\draw[filo] (1.736,-2.564) -- (3.750,-2.564) -- (3.750,-3.321);
\fill (3.664,-3.321) -- (3.836,-3.321) -- (3.750,-3.564) -- cycle;
\node[etichetta, anchor=base] at (2.907,-3.264) {K5};

%--------------------------------------------------------------
% Q1: contatto ausiliario di chiusura -> I2
%--------------------------------------------------------------
\draw[filo] (6.793,-0.536) -- (6.793,-2.564);
\draw[filo] (6.440,-2.483) -- (6.793,-3.064);
\draw[filo] (6.793,-3.064) -- (6.793,-3.321);
\fill (6.707,-3.321) -- (6.879,-3.321) -- (6.793,-3.564) -- cycle;
\node[etichetta, anchor=base] at (6.107,-2.936) {Q1};

%--------------------------------------------------------------
% K3: contatto di apertura e Q1: contatto a specchio -> I3
%--------------------------------------------------------------
\draw[filo] (7.821,-0.536) -- (7.821,-1.036) -- (8.064,-1.036);
\draw[filo] (8.007,-0.964) -- (7.821,-1.550);
\draw[filo] (7.821,-1.550) -- (7.821,-2.564) -- (8.064,-2.564);
\fill (8.050,-2.364) circle (0.086);
\draw[filo] (8.050,-2.364) -- (7.821,-3.050);
\draw[filo] (7.821,-3.050) -- (7.821,-3.321);
\fill (7.736,-3.321) -- (7.907,-3.321) -- (7.821,-3.564) -- cycle;
\draw[meccanico, dash pattern=on 0.186cm off 0.143cm] (6.664,-2.807) -- (7.907,-2.807);
\node[etichetta, anchor=base] at (7.443,-1.336) {K3};

%--------------------------------------------------------------
% K1: contatto di chiusura -> I4
%--------------------------------------------------------------
\draw[filo] (9.836,-0.536) -- (9.836,-1.550);
\draw[filo] (9.479,-1.464) -- (9.836,-2.050);
\draw[filo] (9.836,-2.050) -- (9.836,-3.321);
\fill (9.750,-3.321) -- (9.921,-3.321) -- (9.836,-3.564) -- cycle;
\node[etichetta, anchor=base] at (9.050,-1.893) {K1};

%--------------------------------------------------------------
% S2: secondo contatto di chiusura -> I5
%--------------------------------------------------------------
\draw[filo] (13.207,-9.636) -- (13.207,-2.564) -- (11.850,-2.564) -- (11.850,-3.321);
\fill (11.764,-3.321) -- (11.936,-3.321) -- (11.850,-3.564) -- cycle;

%--------------------------------------------------------------
% PLC
%--------------------------------------------------------------
\draw[filo] (2.750,-3.564) rectangle (12.864,-8.121);
\draw[filo] (2.750,-5.093) -- (12.864,-5.093);
\draw[filo] (2.750,-6.607) -- (12.864,-6.607);
\node[etichetta, anchor=base] at (3.793,-4.021) {I0};
\node[etichetta, anchor=base] at (4.793,-4.021) {I1};
\node[etichetta, anchor=base] at (6.750,-4.021) {I2};
\node[etichetta, anchor=base] at (7.793,-4.021) {I3};
\node[etichetta, anchor=base] at (9.793,-4.021) {I4};
\node[etichetta, anchor=base] at (11.850,-4.021) {I5};
\node[medio, anchor=base] at (7.607,-4.750) {Ingressi};
\node[grande, anchor=base] at (7.629,-6.179) {PLC};
\node[medio, anchor=base] at (7.586,-7.236) {Uscite};
\node[etichetta, anchor=base] at (5.293,-7.950) {O0};
\node[etichetta, anchor=base] at (8.779,-7.950) {O1};
\node[etichetta, anchor=base] at (11.850,-7.950) {O2};

%--------------------------------------------------------------
% S1: pulsante di STOP (due contatti di apertura ad azionamento
% meccanico positivo)
%--------------------------------------------------------------
\draw[filo] (1.736,-2.564) -- (1.736,-9.636) -- (1.979,-9.636);
\draw[filo] (1.931,-9.564) -- (1.746,-10.136);
\draw[filo] (1.736,-10.136) -- (1.736,-10.650);
\draw[sottile] (1.336,-9.450) circle (0.243);
\draw[sottile] (1.179,-9.450) -- (1.479,-9.450);
\draw[sottile] (1.369,-9.341) -- (1.479,-9.450) -- (1.369,-9.559);
\node[etichetta, anchor=base] at (1.721,-11.079) {+\,U\textsubscript{B}};
\draw[filo] (5.279,-8.121) -- (5.279,-9.636) -- (5.536,-9.636);
\draw[filo] (5.474,-9.564) -- (5.283,-10.136);
\draw[filo] (5.279,-10.136) -- (5.279,-11.164);
\draw[sottile] (4.889,-9.450) circle (0.243);
\draw[sottile] (4.731,-9.450) -- (5.031,-9.450);
\draw[sottile] (4.922,-9.341) -- (5.031,-9.450) -- (4.922,-9.559);
\draw[filo] (1.064,-9.693) -- (0.936,-9.693) -- (0.936,-9.950) -- (1.064,-9.950);
\draw[filo] (0.936,-9.821) -- (1.007,-9.821);
\draw[meccanico, dash pattern=on 0.186cm off 0.114cm] (1.121,-9.821) -- (5.407,-9.821);
\node[etichetta, anchor=base] at (0.600,-9.964) {S1};
\node[etichetta, anchor=base] at (0.593,-10.393) {STOP};

%--------------------------------------------------------------
% O0: contatto di chiusura di K2, bobina di K3 con ritardo alla
% diseccitazione e condensatore C1
%--------------------------------------------------------------
\draw[filo] (4.926,-11.074) -- (5.283,-11.660);
\draw[filo] (5.279,-11.664) -- (5.279,-13.693);
\draw[filo] (4.764,-13.693) rectangle (5.779,-14.193);
\draw[filo] (5.279,-14.193) -- (5.279,-16.221);
\node[nodo] at (5.279,-12.679) {};
\draw[filo] (5.279,-12.679) -- (7.293,-12.679) -- (7.293,-13.893);
\draw[filo] (7.107,-13.893) -- (7.507,-13.893);
\draw[filo] (7.107,-14.021) -- (7.507,-14.021);
\draw[filo] (7.293,-14.021) -- (7.293,-16.221);
\node[etichetta, anchor=base] at (4.529,-11.321) {K2};
\node[etichetta, anchor=base] at (4.443,-14.079) {K3};
\node[etichetta, anchor=base] at (6.636,-14.050) {C1};
\node[etichetta, anchor=base] at (3.893,-14.779) {con ritardo alla};
\node[etichetta, anchor=base] at (3.893,-15.179) {diseccitazione};

%--------------------------------------------------------------
% O1: pulsante S2 (ON), contatto ausiliario di Q1 (automantenimento),
% contatto di chiusura di K3 e bobina di Q1
%--------------------------------------------------------------
\draw[filo] (8.821,-8.121) -- (8.821,-9.636);
\node[nodo] at (8.821,-9.136) {};
\draw[filo] (8.821,-9.136) -- (10.336,-9.136) -- (10.336,-10.650);
\draw[filo] (8.474,-9.564) -- (8.821,-10.136);
\draw[filo] (8.821,-10.136) -- (8.821,-12.421) -- (10.336,-12.421);
\node[nodo] at (10.336,-12.421) {};
\draw[filo] (8.021,-9.721) -- (7.893,-9.721) -- (7.893,-9.979) -- (8.021,-9.979);
\draw[filo] (7.893,-9.850) -- (7.964,-9.850);
\draw[meccanico, dash pattern=on 0.186cm off 0.143cm] (8.093,-9.850) -- (13.093,-9.850);
\draw[filo] (12.864,-9.564) -- (13.211,-10.136);
\draw[filo] (13.207,-10.136) -- (13.207,-10.650);
\node[etichetta, anchor=base] at (12.879,-11.079) {+\,U\textsubscript{B}};
\node[etichetta, anchor=base] at (7.364,-9.964) {S2};
\node[etichetta, anchor=base] at (7.379,-10.379) {ON};
\draw[filo] (9.989,-10.579) -- (10.340,-11.160);
\draw[filo] (10.336,-11.164) -- (10.336,-13.179);
\draw[filo] (9.989,-13.103) -- (10.340,-13.697);
\draw[filo] (10.336,-13.693) -- (10.336,-14.693);
\draw[filo] (9.836,-14.693) rectangle (10.836,-15.207);
\draw[filo] (10.336,-15.207) -- (10.336,-16.221);
\node[etichetta, anchor=base] at (9.689,-10.993) {Q1};
\node[etichetta, anchor=base] at (9.546,-13.579) {K3};
\node[etichetta, anchor=base] at (9.414,-15.064) {Q1};
\node[etichetta, anchor=base] at (13.350,-14.921) {2\textsuperscript{o} percorso di disinserzione};
\node[etichetta, anchor=base] at (13.350,-15.336) {ritardato};

%--------------------------------------------------------------
% O2: contatto ausiliario di Q1 -> ingresso di avvio/arresto di T1
% (1o percorso di disinserzione)
%--------------------------------------------------------------
\draw[filo] (11.850,-8.121) -- (11.850,-10.650);
\draw[filo] (11.511,-10.579) -- (11.854,-11.160);
\draw[filo] (11.850,-11.164) -- (11.850,-12.164) -- (19.493,-12.164) -- (19.493,-10.150) -- (20.264,-10.150);
\fill (20.264,-10.064) -- (20.264,-10.236) -- (20.507,-10.150) -- cycle;
\draw[meccanico, dash pattern=on 0.200cm off 0.129cm] (10.193,-10.864) -- (11.664,-10.864);
\node[etichetta, anchor=base] at (15.229,-12.036) {1\textsuperscript{o} percorso di disinserzione};

%--------------------------------------------------------------
% S3, S4: dispositivi di arresto di emergenza (due contatti di
% apertura ad azionamento meccanico positivo)
%--------------------------------------------------------------
\draw[filo] (15.143,-0.536) -- (15.143,-1.607) -- (15.443,-1.607);
\draw[filo] (15.386,-1.521) -- (15.143,-2.564);
\draw[filo] (16.114,-0.536) -- (16.114,-1.607) -- (16.414,-1.607);
\draw[filo] (16.357,-1.521) -- (16.114,-2.564);
\draw[filo] (14.021,-1.686) arc[start angle=90, end angle=270, x radius=0.171, y radius=0.221] -- cycle;
\draw[filo] (14.021,-1.907) -- (14.107,-1.907);
\draw[filo] (14.379,-1.907) -- (14.479,-1.907) -- (14.621,-2.164) -- (14.750,-1.907) -- (14.893,-1.907);
\draw[meccanico, dash pattern=on 0.214cm off 0.143cm] (15.036,-1.907) -- (16.264,-1.907);
\draw[sottile] (14.460,-2.550) circle (0.286);
\draw[sottile] (14.260,-2.550) -- (14.646,-2.550);
\draw[sottile] (14.517,-2.421) -- (14.646,-2.550) -- (14.517,-2.679);
\draw[sottile] (14.746,-2.550) -- (16.114,-2.550);
\draw[sottile] (15.014,-2.421) -- (15.143,-2.550) -- (15.014,-2.679);
\draw[sottile] (15.986,-2.421) -- (16.114,-2.550) -- (15.986,-2.679);
\node[etichetta, anchor=base] at (13.521,-2.021) {S3};
\node[etichetta, anchor=base] at (14.664,-3.207) {S3.1};
\node[etichetta, anchor=base] at (15.636,-3.207) {S3.2};
\draw[filo] (15.386,-3.936) -- (15.143,-4.979);
\draw[filo] (16.357,-3.936) -- (16.114,-4.979);
\draw[filo] (14.021,-4.100) arc[start angle=90, end angle=270, x radius=0.171, y radius=0.221] -- cycle;
\draw[filo] (14.021,-4.321) -- (14.107,-4.321);
\draw[filo] (14.379,-4.321) -- (14.479,-4.321) -- (14.621,-4.579) -- (14.750,-4.321) -- (14.893,-4.321);
\draw[meccanico, dash pattern=on 0.214cm off 0.143cm] (15.036,-4.321) -- (16.264,-4.321);
\draw[sottile] (14.460,-4.964) circle (0.286);
\draw[sottile] (14.260,-4.964) -- (14.646,-4.964);
\draw[sottile] (14.517,-4.836) -- (14.646,-4.964) -- (14.517,-5.093);
\draw[sottile] (14.746,-4.964) -- (16.114,-4.964);
\draw[sottile] (15.014,-4.836) -- (15.143,-4.964) -- (15.014,-5.093);
\draw[sottile] (15.986,-4.836) -- (16.114,-4.964) -- (15.986,-5.093);
\node[etichetta, anchor=base] at (13.521,-4.436) {S4};
\node[etichetta, anchor=base] at (14.664,-5.621) {S4.1};
\node[etichetta, anchor=base] at (15.636,-5.621) {S4.2};
\draw[filo] (15.143,-2.564) -- (15.143,-4.021) -- (15.443,-4.021);
\draw[filo] (16.114,-2.564) -- (16.114,-4.021) -- (16.414,-4.021);
\draw[filo] (15.143,-4.979) -- (15.143,-6.293);
\draw[filo] (16.114,-4.979) -- (16.114,-6.293);

%--------------------------------------------------------------
% K4: dispositivo di comando di sicurezza
%--------------------------------------------------------------
\draw[filo] (14.336,-6.293) rectangle (19.893,-8.121);
\draw[filo] (16.736,-6.293) -- (16.736,-8.121);
\node[etichetta, anchor=base] at (14.000,-6.579) {K4};
\draw[filo] (14.964,-8.121) -- (14.964,-8.379) -- (15.221,-8.379);
\draw[filo] (15.164,-8.307) -- (14.964,-8.879);
\draw[filo] (14.964,-8.879) -- (14.964,-9.136) -- (15.221,-9.136);
\draw[filo] (15.164,-9.064) -- (14.964,-9.636);
\draw[filo] (14.964,-9.636) -- (14.964,-9.893) -- (16.193,-9.893) -- (16.193,-8.121);
\node[etichetta, anchor=base] at (14.636,-8.793) {K1};
\node[etichetta, anchor=base] at (14.650,-9.464) {K2};

%--------------------------------------------------------------
% Contatti di uscita di K4 e bobine di K1 e K2
%--------------------------------------------------------------
\draw[filo] (18.064,-0.536) -- (18.064,-7.121);
\draw[filo] (17.707,-7.036) -- (18.064,-7.607);
\draw[filo] (18.064,-7.607) -- (18.064,-14.693);
\draw[filo] (17.550,-14.693) rectangle (18.564,-15.207);
\draw[filo] (18.064,-15.207) -- (18.064,-16.221);
\node[nodo] at (18.064,-4.007) {};
\draw[filo] (18.064,-4.007) -- (19.079,-4.007) -- (19.079,-7.121);
\draw[filo] (18.721,-7.036) -- (19.079,-7.607);
\draw[filo] (19.079,-7.607) -- (19.079,-13.179);
\draw[filo] (18.564,-13.179) rectangle (19.579,-13.693);
\draw[filo] (19.079,-13.693) -- (19.079,-16.221);
\draw[meccanico, dash pattern=on 0.186cm off 0.129cm] (17.921,-7.364) -- (18.936,-7.364);
\node[etichetta, anchor=base] at (18.700,-13.050) {K1};
\node[etichetta, anchor=base] at (17.207,-15.079) {K2};

%--------------------------------------------------------------
% Circuito di potenza: linea L, contatto principale di Q1,
% convertitore di frequenza T1 e motore M
%--------------------------------------------------------------
\node[etichetta, anchor=base] at (22.536,-0.893) {L};
\draw[filo] (22.540,-1.036) -- (22.540,-4.536);
\draw[filo] (22.283,-2.553) -- (22.800,-2.297);
\draw[filo] (22.283,-2.696) -- (22.800,-2.440);
\draw[filo] (22.283,-2.839) -- (22.800,-2.583);
\draw[filo] (22.540,-4.136) arc[start angle=90, end angle=270, x radius=0.179, y radius=0.207];
\draw[filo] (22.093,-4.336) -- (22.540,-5.100);
\draw[filo] (22.540,-5.100) -- (22.540,-6.850);
\node[etichetta, anchor=base] at (21.764,-4.807) {Q1};
\node[etichetta, anchor=base] at (20.643,-6.650) {T1};
\draw[filo] (20.507,-6.850) rectangle (23.550,-12.164);
\draw[filo] (20.507,-7.621) -- (23.550,-7.621);
\draw[filo] (20.507,-11.407) -- (23.550,-11.407);
\draw[sottile] (20.507,-7.607) -- (23.550,-6.864);
\draw[sottile] (20.507,-12.164) -- (23.550,-11.421);
\node[testo, anchor=base] at (22.036,-8.579) {Convertitore};
\node[testo, anchor=base] at (22.036,-9.079) {di frequenza};
\node[testo, anchor=base] at (22.036,-10.264) {Avvio / arresto};
\draw[filo] (22.540,-12.164) -- (22.540,-14.450);
\draw[filo] (22.293,-13.279) -- (22.793,-13.050);
\draw[filo] (22.293,-13.421) -- (22.793,-13.193);
\draw[filo] (22.293,-13.564) -- (22.793,-13.336);
\draw[sottile] (22.536,-15.207) circle (0.757);
\node[motore, anchor=base] at (22.536,-15.107) {M};
\node[motore, anchor=base] at (22.536,-15.750) {3$\sim$};

%--------------------------------------------------------------
% Linea di massa
%--------------------------------------------------------------
\draw[filo] (0.721,-16.221) -- (20.436,-16.221);
\node[nodo] at (2.236,-16.221) {};
\node[nodo] at (5.279,-16.221) {};
\node[nodo] at (7.293,-16.221) {};
\node[nodo] at (10.336,-16.221) {};
\node[nodo] at (18.064,-16.221) {};
\node[nodo] at (19.079,-16.221) {};
\draw[filo] (2.236,-16.221) -- (2.236,-16.593);
\draw[filo] (1.979,-16.593) -- (2.507,-16.593);
\draw[filo] (2.050,-16.721) -- (2.450,-16.721);
\draw[filo] (2.107,-16.864) -- (2.379,-16.864);

\end{tikzpicture}
}
        \captionof{figure}{Arresto di un azionamento con convertitore di frequenza comandato da PLC a seguito di un comando di arresto di emergenza.}
        \label{fig:08-19}
    \end{minipage}
\end{center}

\textbf{Funzioni di sicurezza}
\begin{itemize}
    \item il movimento pericoloso viene arrestato azionando il pulsante di arresto S1 oppure uno dei dispositivi di arresto di emergenza S3 o S4. In questo esempio si considera unicamente l'azionamento del dispositivo di arresto di emergenza S3. In caso di emergenza l'azionamento viene arrestato in risposta alla manovra di S3: dapprima mediante la disattivazione del modulo di sicurezza per arresto di emergenza K4, con conseguente diseccitazione dei contattori ausiliari K1 e K2. L'apertura del contatto di chiusura K1 sull'ingresso I4 del PLC K5 determina l'annullamento del segnale di avvio sul convertitore di frequenza (FI) T1 tramite l'uscita O2 del PLC. In modo ridondante rispetto alla catena K1-K5-T1, l'apertura del contatto di chiusura K2 a monte del contattore ausiliario con ritardo alla diseccitazione K3 avvia un temporizzatore di frenatura; allo scadere del temporizzatore di frenatura viene interrotto il segnale di comando del contattore di rete Q1. L'impostazione del temporizzatore è scelta in modo tale che, in condizioni di funzionamento sfavorevoli, il movimento della macchina si arresti prima che il contattore di rete Q1 si sia diseccitato.
    \item l'arresto funzionale dell'azionamento a seguito di un comando di arresto è avviato dall'apertura dei due contatti di apertura del pulsante di arresto S1. Come per l'arresto in caso di emergenza, lo stato viene dapprima interrogato dal PLC K5 tramite l'ingresso I0 e il convertitore di frequenza viene disattivato mediante il reset dell'uscita O2 del PLC. In modo ridondante rispetto a questo processo, il contattore ausiliario K3 viene diseccitato - con il ritardo alla diseccitazione fornito dal condensatore C1 - e, allo scadere del tempo di frenatura impostato, viene interrotto il segnale di attivazione del contattore di rete Q1.
    \item in caso di guasto del PLC K5, del convertitore di frequenza T1, del contattore di rete Q1, dei contattori ausiliari K1/K2 o del contattore ausiliario con ritardo alla diseccitazione K3, l'arresto dell'azionamento è comunque garantito, poiché sono sempre presenti due percorsi di disinserzione indipendenti l'uno dall'altro. La mancata diseccitazione dei contattori ausiliari K1 e K2 viene rilevata - al più tardi dopo il ripristino del dispositivo di arresto di emergenza azionato - mediante il monitoraggio dei contatti di apertura collegati meccanicamente all'interno del modulo di sicurezza per arresto di emergenza K4. La mancata diseccitazione del contattore ausiliario K3 viene rilevata - al più tardi prima del nuovo avvio del movimento della macchina - tramite la retroazione all'ingresso I3 del PLC del contatto di apertura collegato meccanicamente. La mancata diseccitazione del contattore di rete Q1 viene rilevata dal contatto a specchio letto sull'ingresso I3 del PLC. L'incollaggio di questo contatto a specchio viene rilevato dal contatto ausiliario di chiusura collegato meccanicamente sull'ingresso I2 del PLC. In caso di guasto del condensatore C1, il tempo di diseccitazione misurato del contattore ausiliario K3 differisce dal tempo impostato nel PLC. Il guasto viene rilevato e determina lo spegnimento della macchina e il blocco del suo funzionamento. Misure organizzative garantiscono che ogni dispositivo di arresto di emergenza venga azionato almeno una volta all'anno.
\end{itemize}

\textbf{Caratteristiche di progetto}
\begin{itemize}
    \item sono rispettati i principi di sicurezza di base e i principi di sicurezza ben collaudati e sono soddisfatti i requisiti della Categoria B. Sono realizzati i circuiti di protezione (per es. la protezione dei contatti) descritti nei paragrafi iniziali del Capitolo 8.
    \item •	I contattori ausiliari K1, K2 e K3 dispongono di contatti collegati meccanicamente secondo la IEC 60947-5-1, Allegato L.
    \item i pulsanti S1, S3 e S4 dispongono di contatti ad apertura positiva secondo la IEC 60947-5-1, Allegato K.
    \item il contattore Q1 dispone di un contatto a specchio secondo la IEC 60947-4-1, Allegato F.
    \item i componenti standard K5 e T1 sono impiegati conformemente alle indicazioni del punto 6.3.10.
    \item il software (SRASW) è programmato secondo i requisiti per il PL b (declassato in ragione della diversità) e le indicazioni del punto 6.3.10.
    \item il raggiungimento ritardato dell'arresto per mezzo del solo secondo percorso di disinserzione, in caso di guasto, non deve comportare un rischio residuo inaccettabilmente elevato.
    \item Le SRP/CS del modulo di sicurezza per arresto di emergenza K4 soddisfano tutti i requisiti per la Categoria 3 e il PL d.
\end{itemize}

\textbf{Calcolo della probabilità di guasto}
Viene calcolata unicamente la probabilità di guasto della funzione di arresto di emergenza.
\begin{itemize}
    \item il dispositivo di arresto di emergenza S3 è dotato di due contatti di apertura S3.1 e S3.2. Il costruttore indica un B\textsubscript{10D} di 127.500 cicli per ciascuno dei blocchi S3.1 e S3.2. Con un azionamento annuale e un nop risultante di 1 ciclo all'anno, l'MTTF\textsubscript{D} di ciascun contatto è di 1.275.000 anni. Il modulo di sicurezza per arresto di emergenza K4 è un componente di sicurezza sottoposto a prove. La sua probabilità di guasto è di 3,0 10\textsubscript{-7} all'ora [M] e viene sommata al termine del calcolo.
\end{itemize}

Per la probabilità di guasto della struttura a due canali posta a valle vale quanto segue:
\begin{itemize}
    \item MTTF\textsubscript{D}: il PLC K5 ha un MTTFD di dieci anni [S]. Il convertitore di frequenza ha un MTTF\textsubscript{D} di 35 anni [M]. Il condensatore C1 è considerato nel calcolo con un MTTF\textsubscript{D} di 45.662 anni [D]. Con un valore B\textsubscript{10D} di 5.000.000 di cicli [M] e una frequenza di manovra pari a un'eccitazione giornaliera su 240 giorni lavorativi, si ottiene un MTTFD di 208.333 anni per K1 e K2. Con un valore B10D di 2.000.000 di cicli [M] e con 240 giorni lavorativi, 16 ore lavorative giornaliere e un tempo di ciclo di 3 minuti, nop è pari a 76.800 cicli all'anno e l'MTTFD di K3 è di 260 anni. Con un valore B10D di 600.000 cicli [M] e con 240 giorni lavorativi, 16 ore lavorative giornaliere e un tempo di ciclo di 3 minuti, nop è pari a 76.800 cicli all'anno e l'MTTFD di Q1 è di 7,8 anni. Da questi valori risulta per il canale un MTTFD simmetrizzato di 60 anni ("alto").
    \item DC\textsubscript{avg}: è garantita una frequenza di prova adeguata dei dispositivi di arresto di emergenza (cfr. le indicazioni dei punti 6.2.14 e D.2.5.1). Il rilevamento dei guasti dei blocchi S3.1 e S3.2 è ottenuto mediante monitoraggio incrociato in K4 (DC = 90\%). Il rilevamento dei guasti da parte del processo in caso di mancata attuazione della rampa di decelerazione porta a un DC del 60\% per K5. Anche per T1 il DC è del 60\%, sempre per effetto del rilevamento dei guasti da parte del processo. K1 e K2 presentano un DC del 99\% grazie al rilevamento dei guasti integrato in K4. Per K3 il DC è del 99\% grazie al rilevamento dei guasti da parte di K5. Per C1 il DC è del 60\% grazie alla verifica nel PLC dell'elemento temporizzatore, in occasione della disinserzione del convertitore di frequenza, attraverso il tempo di diseccitazione del contattore ausiliario K3. Per Q1 il DC è quindi del 99\% grazie al monitoraggio diretto in K5. La formula di media per il DC\textsubscript{avg} fornisce un risultato del 64\% ("basso").
    \item misure adeguate contro i guasti per causa comune (75 punti): separazione (15), diversità (20), FMEA (5) e condizioni ambientali (25 + 10)
    \item la combinazione a due canali degli elementi di comando soddisfa la Categoria 3. Ne risulta una probabilità media di guasto pericoloso PFHD di 3,9 10\textsuperscript{-7} all'ora. Ciò soddisfa il PL d. Sommando la probabilità di guasto pericoloso di K4 e S3 si ottiene una probabilità di guasto complessiva di 7,4 10\textsuperscript{-7} all'ora. Anche in questo caso è soddisfatto il PL d.
    \item il contattore Q1, soggetto a usura, dovrebbe essere sostituito dopo circa 7,8 anni (T\textsubscript{10D}).
\end{itemize}
\clearpage

\textbf{Esempio 11}
\textit{Valvola pneumatica testata (sottosistema) - Categoria 2 - PL d}

% La minipage tiene insieme disegno e didascalia (figura non fluttuante)
\begin{center}
    \begin{minipage}{\textwidth}
        \centering
        \captionsetup{hypcap=false}
        \resizebox{0.84\textwidth}{!}{%==============================================================
%  Figura 8.21 - Valvola pneumatica con test elettronico per il
%  comando di movimenti pericolosi (Esempio 11)
%  Adattamento in italiano della Figura 8.21 dell'IFA Report 2/2017
%
%  I simboli dei componenti riproducono quelli dell'originale
%  (ISO 1219): le coordinate sono ricavate dal disegno di partenza
%  (1 cm = 84 px, origine nell'angolo in alto a sinistra della
%  cornice); a questa scala i simboli hanno le stesse dimensioni
%  della Figura 8.5, da cui sono ripresi 0V1 e l'unita' di
%  manutenzione 0Z. Nessun cambio di corpo del font: le etichette
%  sono ingrandite con "scale" (vedere la nota nella Figura 8.3).
%==============================================================
\begin{tikzpicture}[
  font=\normalsize,
  linea/.style={draw=black, line width=0.8pt},
  simbolo/.style={draw=black, line width=0.8pt},
  sottile/.style={draw=black, line width=0.6pt},
  freccia/.style={sottile, -{Latex[length=8pt,width=3.6pt]}},
  doppia/.style={sottile, {Latex[length=8pt,width=3.6pt]}-{Latex[length=8pt,width=3.6pt]}},
  segnale/.style={sottile, -{Latex[length=8pt,width=4pt]}},
  ingresso/.style={linea, -{Latex[length=5.5pt,width=3.6pt]}},
  pilota/.style={draw=black, line width=0.6pt, dash pattern=on 3.6pt off 2pt},
  zona/.style={draw={rgb,255:red,22;green,67;blue,38}, line width=0.6pt,
               dash pattern=on 0.636cm off 0.1cm on 0.4pt off 0.1cm},
  etichetta/.style={inner sep=0pt, scale=1.2},
  testo/.style={inner sep=0pt, scale=0.95},
  piccolo/.style={inner sep=0pt, scale=0.8},
  grande/.style={inner sep=0pt, scale=1.5},
  nodo/.style={circle, fill=black, inner sep=0pt, minimum size=5.2pt}
]

%--------------------------------------------------------------
% Cornice
%--------------------------------------------------------------
\draw[draw=black, line width=1.2pt] (0.000,-0.000) rectangle (16.369,-20.155);

%--------------------------------------------------------------
% Cilindro 1A: camicia, pistone e stelo a doppia linea, blocco di
% attacco con le due bocche
%--------------------------------------------------------------
\draw[simbolo] (3.470,-2.125) rectangle (5.625,-3.327);
\draw[simbolo] (3.708,-2.125) -- (3.708,-3.327);
\draw[simbolo] (4.185,-2.125) -- (4.185,-3.327);
\draw[simbolo, fill=white] (4.185,-2.601) rectangle (6.351,-2.839);
\draw[linea] (3.530,-3.327) -- (3.530,-4.173) -- (4.292,-4.173) -- (4.292,-6.381);
\draw[linea] (5.565,-3.327) -- (5.565,-4.173) -- (4.768,-4.173) -- (4.768,-6.381);
\node[etichetta, anchor=base] at (2.702,-2.452) {1A};

%--------------------------------------------------------------
% Sensore di posizione 1S1 sull'estremita' dello stelo, con uscita
% verso K1
%--------------------------------------------------------------
\draw[linea] (6.351,-2.839) -- (6.351,-1.887);
\draw[simbolo] (6.351,-0.923) rectangle (7.304,-1.887);
\draw[sottile] (6.369,-0.940) -- (7.286,-1.869);
\draw[sottile] (6.827,-1.161) -- (7.012,-1.161) -- (7.012,-1.048) -- (7.199,-1.048);
\node[etichetta, anchor=base] at (6.595,-1.774) {G};
\draw[segnale] (7.304,-0.923) -- (7.304,-0.298) -- (7.917,-0.643);
\node[etichetta, anchor=base] at (8.345,-0.857) {K1};
\node[etichetta, anchor=base] at (5.601,-1.250) {1S1};

%--------------------------------------------------------------
% Movimento pericoloso
%--------------------------------------------------------------
\draw[sottile, {Latex[length=6pt,width=3.4pt]}-{Latex[length=6pt,width=3.4pt]}] (7.333,-2.720) -- (10.643,-2.720);
\draw[sottile] (8.994,-2.607) -- (8.994,-2.845);
\node[testo, anchor=base] at (9.012,-1.893) {movimento};
\node[testo, anchor=base] at (9.012,-2.345) {pericoloso};

%--------------------------------------------------------------
% 1V1: valvola 4/3 a centro chiuso con scarichi, azionata da
% elettromagneti su entrambi i lati (a, b) e centrata da molle
%--------------------------------------------------------------
\draw[simbolo] (3.089,-6.089) rectangle (5.970,-7.042);
\draw[simbolo] (4.054,-6.089) -- (4.054,-7.042);
\draw[simbolo] (5.018,-6.089) -- (5.018,-7.042);
% elettromagnete sinistro (a) e collegamento con K1
\draw[simbolo] (1.173,-6.565) rectangle (3.089,-7.042);
\draw[simbolo] (2.137,-6.565) -- (2.137,-7.042);
\draw[sottile] (1.545,-7.042) -- (1.783,-6.565);
\draw[sottile] (2.137,-6.583) -- (2.562,-6.810) -- (2.137,-7.030);
\draw[segnale] (1.783,-6.565) -- (1.783,-6.065) -- (2.274,-5.789);
% molla sinistra
\draw[sottile] (2.378,-6.333) -- (2.446,-6.565) -- (2.562,-6.095) -- (2.690,-6.565) -- (2.804,-6.095) -- (2.923,-6.565) -- (3.042,-6.095) -- (3.089,-6.250);
% elettromagnete destro (b) e collegamento con K1
\draw[simbolo] (5.970,-6.565) rectangle (7.893,-7.042);
\draw[simbolo] (6.932,-6.565) -- (6.932,-7.042);
\draw[sottile] (7.298,-6.565) -- (7.527,-7.042);
\draw[sottile] (6.932,-6.583) -- (6.524,-6.810) -- (6.932,-7.030);
\draw[segnale] (7.298,-6.565) -- (7.298,-6.065) -- (7.789,-5.789);
% molla destra
\draw[sottile] (6.688,-6.333) -- (6.634,-6.565) -- (6.515,-6.095) -- (6.396,-6.565) -- (6.277,-6.095) -- (6.158,-6.565) -- (6.039,-6.095) -- (5.970,-6.280);
% posizione sinistra
\draw[freccia] (3.568,-7.042) -- (3.345,-6.095);
\draw[freccia] (3.815,-6.089) -- (3.815,-7.042);
\draw[sottile] (3.351,-7.042) -- (3.351,-6.762);
\draw[sottile] (3.250,-6.762) -- (3.440,-6.762);
% posizione centrale (tutte le bocche chiuse)
\draw[sottile] (4.214,-6.381) -- (4.387,-6.381);
\draw[sottile] (4.690,-6.381) -- (4.863,-6.381);
\draw[sottile] (4.214,-6.762) -- (4.381,-6.762);
\draw[sottile] (4.452,-6.762) -- (4.619,-6.762);
\draw[sottile] (4.690,-6.762) -- (4.857,-6.762);
% posizione destra
\draw[freccia] (5.259,-6.089) -- (5.259,-7.042);
\draw[freccia] (5.503,-7.042) -- (5.750,-6.095);
\draw[sottile] (5.741,-7.042) -- (5.741,-6.762);
\draw[sottile] (5.655,-6.762) -- (5.833,-6.762);
% scarichi
\draw[linea] (4.292,-6.762) -- (4.292,-7.536);
\draw[linea] (4.768,-6.762) -- (4.768,-7.536);
\draw[sottile] (4.179,-7.536) -- (4.417,-7.536) -- (4.292,-7.738) -- cycle;
\draw[sottile] (4.655,-7.536) -- (4.893,-7.536) -- (4.768,-7.738) -- cycle;
\node[etichetta, anchor=base] at (0.738,-5.893) {1V1};
\node[etichetta, anchor=base] at (2.327,-5.524) {K1};
\node[etichetta, anchor=base] at (7.875,-5.583) {K1};
\node[etichetta, anchor=base] at (0.798,-6.952) {a};
\node[etichetta, anchor=base] at (8.256,-7.000) {b};

%--------------------------------------------------------------
% Linea di alimentazione, derivazione verso ulteriori utenze e
% pressostato 0S1
%--------------------------------------------------------------
\draw[linea] (4.530,-6.762) -- (4.530,-12.565);
\node[nodo] at (4.530,-8.970) {};
\draw[linea] (4.530,-8.970) -- (5.976,-8.970);
\draw[sottile] (5.976,-8.851) -- (6.179,-8.970) -- (5.976,-9.089) -- cycle;
\node[testo, anchor=base west] at (6.607,-8.845) {ulteriori utenze e};
\node[testo, anchor=base west] at (6.607,-9.238) {sistemi di comando};
\node[nodo] at (4.530,-10.887) {};
\draw[linea] (4.530,-10.887) -- (5.970,-10.887);
\draw[simbolo] (5.970,-10.399) rectangle (6.935,-11.363);
\draw[sottile] (5.976,-11.357) -- (6.929,-10.405);
\node[etichetta, anchor=base] at (6.327,-10.881) {P};
\draw[sottile] (6.330,-11.250) -- (6.500,-11.250) -- (6.500,-11.131) -- (6.699,-11.131);
\node[etichetta, anchor=base] at (5.321,-10.417) {0S1};

%--------------------------------------------------------------
% PLC K1 con i segnali di ingresso da 1V1a, 1V1b, 1S1 e l'uscita
% verso 0V1
%--------------------------------------------------------------
\draw[simbolo] (11.655,-7.065) rectangle (15.435,-9.232);
\draw[simbolo] (11.655,-7.613) -- (15.435,-7.613);
\draw[simbolo] (11.655,-8.685) -- (15.435,-8.685);
\node[piccolo, anchor=base] at (13.565,-7.393) {Ingressi};
\node[grande, anchor=base] at (13.565,-8.345) {PLC};
\node[piccolo, anchor=base] at (13.565,-9.036) {Uscite};
\node[etichetta, anchor=base] at (11.083,-7.405) {K1};
\draw[ingresso] (12.143,-5.932) -- (12.143,-7.065);
\draw[segnale] (12.143,-5.932) -- (11.532,-6.289);
\draw[ingresso] (13.577,-5.146) -- (13.577,-7.065);
\draw[segnale] (13.577,-5.146) -- (12.968,-5.504);
\draw[ingresso] (15.012,-4.881) -- (15.012,-7.065);
\draw[segnale] (15.012,-4.881) -- (14.527,-4.615);
\node[etichetta, anchor=base] at (10.756,-6.381) {1V1a};
\node[etichetta, anchor=base] at (12.185,-5.607) {1V1b};
\node[etichetta, anchor=base] at (13.762,-4.750) {1S1};
\draw[segnale] (13.545,-9.232) -- (13.545,-10.446) -- (13.062,-10.729);
\node[etichetta, anchor=base] at (12.446,-10.869) {0V1};

%--------------------------------------------------------------
% 0V1: valvola 3/2 con elettromagnete e ritorno a molla (dalla
% Figura 8.5), collegamento con K1
%--------------------------------------------------------------
\draw[simbolo] (3.339,-12.565) rectangle (5.254,-13.522);
\draw[simbolo] (4.291,-12.565) -- (4.291,-13.522);
% elettromagnete
\draw[simbolo] (1.419,-13.055) rectangle (3.339,-13.522);
\draw[simbolo] (2.382,-13.055) -- (2.382,-13.522);
\draw[sottile] (1.791,-13.522) -- (2.015,-13.055);
\draw[sottile] (2.382,-13.055) -- (2.801,-13.294) -- (2.382,-13.522);
% molla
\draw[sottile] (5.254,-13.294) -- (5.315,-13.055) -- (5.434,-13.522) -- (5.562,-13.055)
  -- (5.672,-13.522) -- (5.801,-13.055) -- (5.919,-13.522) -- (5.976,-13.279);
% posizione di passaggio
\draw[freccia] (3.576,-13.522) -- (3.576,-12.565);
\draw[sottile] (4.072,-13.522) -- (4.072,-13.241);
\draw[sottile] (3.982,-13.241) -- (4.144,-13.241);
% posizione di scarico
\draw[freccia] (4.539,-12.565) -- (5.015,-13.522);
\draw[sottile] (4.448,-13.241) -- (4.615,-13.241);
\draw[linea] (5.015,-13.522) -- (5.015,-14.041);
\draw[sottile] (4.896,-14.041) -- (5.134,-14.041) -- (5.015,-14.241) -- cycle;
\draw[segnale] (1.783,-13.518) -- (1.783,-14.976) -- (2.393,-14.623);
\node[etichetta, anchor=base] at (1.786,-12.405) {0V1};
\node[etichetta, anchor=base] at (2.893,-14.738) {K1};

%--------------------------------------------------------------
% Linea di alimentazione verso l'unita' di manutenzione 0Z
%--------------------------------------------------------------
\draw[linea] (4.530,-13.226) -- (4.530,-17.601) -- (8.649,-17.601);
\draw[linea] (9.613,-17.601) -- (10.369,-17.601);
\draw[linea] (11.726,-17.601) -- (12.958,-17.601);
\draw[linea] (13.911,-17.601) -- (15.702,-17.601);
% sorgente di aria compressa
\draw[sottile] (15.702,-17.601) -- (16.119,-17.363) -- (16.119,-17.845) -- cycle;
\draw[zona] (6.101,-14.958) rectangle (15.030,-19.089);
\node[etichetta, anchor=base] at (5.631,-15.333) {0Z};

%--------------------------------------------------------------
% Componenti dell'unita' di manutenzione 0Z (dalla Figura 8.5)
%--------------------------------------------------------------
% manometro
\draw[simbolo] (7.214,-16.643) circle (0.479);
\draw[sottile, -{Latex[length=6pt,width=3.6pt]}] (7.428,-16.862) -- (7.000,-16.433);
\draw[linea] (7.214,-17.122) -- (7.214,-17.601);
\node[nodo] at (7.214,-17.601) {};

% regolatore di pressione con molla regolabile, sfiato secondario
% e pilotaggio
\draw[simbolo] (8.649,-16.887) rectangle (9.602,-17.845);
\draw[doppia] (8.656,-17.595) -- (9.594,-17.595);
\draw[sottile] (9.602,-17.016) -- (9.852,-17.123) -- (9.602,-17.245);
\draw[sottile] (9.380,-16.887) -- (9.227,-16.816) -- (9.602,-16.730) -- (9.227,-16.630)
  -- (9.602,-16.545) -- (9.227,-16.445) -- (9.602,-16.352) -- (9.416,-16.302);
\draw[sottile, -{Latex[length=6pt,width=3.6pt]}] (9.844,-16.687) -- (8.763,-16.445);
\draw[pilota] (8.649,-17.595) -- (8.409,-17.863) -- (8.409,-18.263) -- (9.387,-18.263)
  -- (9.387,-17.845);

% filtro con separatore d'acqua e indicatore di intasamento
\draw[simbolo] (11.057,-16.925) -- (11.724,-17.607) -- (11.057,-18.288) -- (10.371,-17.607) -- cycle;
\draw[pilota] (11.057,-16.925) -- (11.057,-17.845);
\draw[sottile] (10.610,-17.845) -- (11.500,-17.845);
\draw[sottile] (10.738,-17.845) -- (11.057,-18.164) -- (11.381,-17.845);
\draw[linea] (11.057,-18.288) -- (11.057,-19.607);
\draw[simbolo] (10.477,-16.211) rectangle (11.657,-16.925);
\draw[sottile, -{Latex[length=6pt,width=3.6pt]}] (11.010,-16.782) -- (10.581,-16.354);
\draw[sottile] (11.305,-16.582) circle (0.186);
\draw[sottile] (11.174,-16.451) -- (11.437,-16.714);
\draw[sottile] (11.174,-16.714) -- (11.437,-16.451);

% valvola di intercettazione manuale 3/2 con scarico
\draw[simbolo] (12.958,-16.411) rectangle (13.909,-18.335);
\draw[simbolo] (12.958,-17.372) -- (13.909,-17.372);
% posizione di passaggio (bocca chiusa a destra)
\draw[doppia] (12.965,-16.658) -- (13.902,-16.658);
\draw[sottile] (13.629,-17.063) -- (13.629,-17.221);
\draw[sottile] (13.629,-17.135) -- (13.909,-17.135);
% posizione di scarico (bocca chiusa sulla linea)
\draw[doppia] (12.965,-17.613) -- (13.902,-18.083);
\draw[sottile] (13.629,-17.525) -- (13.629,-17.692);
\draw[sottile] (13.629,-17.611) -- (13.909,-17.611);
\draw[linea] (13.909,-18.086) -- (14.386,-18.086);
\draw[sottile] (14.386,-17.968) -- (14.600,-18.086) -- (14.386,-18.215) -- cycle;
% azionamento manuale con arresto
\draw[linea] (13.615,-16.411) -- (13.615,-15.458);
\draw[linea] (13.366,-16.411) -- (13.366,-16.049) -- (13.495,-15.935) -- (13.366,-15.825)
  -- (13.366,-15.458);
\draw[sottile] (13.248,-15.935) -- (13.366,-15.935);
\draw[simbolo] (13.248,-15.458) -- (13.733,-15.458);


\end{tikzpicture}
}
        \captionof{figure}{Valvola pneumatica con test elettronico per il comando di movimenti pericolosi.}
        \label{fig:08-21}
    \end{minipage}
\end{center}

\textbf{Funzioni di sicurezza}
\begin{itemize}
    \item funzione di arresto legata alla sicurezza: arresto di un movimento pericoloso e prevenzione dell'avviamento inatteso dalla posizione di riposo, realizzata mediante SSC e, in caso di guasti rilevati (rilevamento dei guasti), mediante SDE
    \item qui è rappresentata unicamente la parte pneumatica del sistema di comando, sotto forma di sottosistema. Per completare la funzione di sicurezza devono essere aggiunte, sempre sotto forma di sottosistemi, ulteriori SRP/CS (per es. dispositivi di protezione ed elementi logici elettrici).
\end{itemize}

\textbf{Descrizione delle funzionalità}
\begin{itemize}
    \item i movimenti pericolosi sono comandati da una valvola distributrice 1V1.
    \item un guasto della valvola distributrice 1V1 fra due test funzionali può comportare la perdita della funzione di sicurezza. Il guasto dipende dall'affidabilità della valvola distributrice.
    \item il test della funzione di sicurezza è forzato dal PLC K1 mediante un sistema di misura dello spostamento 1S1. Il test viene eseguito a intervalli adeguati e in occasione di una richiesta della funzione di sicurezza. Il rilevamento di un guasto di 1V1 determina la disinserzione della valvola di scarico 0V1.
    \item l'interruzione del movimento pericoloso mediante la valvola di scarico 0V1 comporta in genere una sovracorsa maggiore. La distanza dalla zona pericolosa deve essere scelta tenendo conto di tale maggiore sovracorsa.
    \item la funzione di test non deve essere compromessa da un guasto della valvola distributrice. Un guasto della funzione di test non deve provocare il guasto della valvola distributrice.
    \item se l'aria compressa intrappolata costituisce un ulteriore pericolo, sono necessarie misure aggiuntive.
\end{itemize}

\textbf{Caratteristiche progettuali}
\begin{itemize}
    \item sono rispettati i principi di sicurezza di base e i principi di sicurezza ben collaudati e sono soddisfatti i requisiti della Categoria B.
    \item 1V1 è una valvola distributrice con posizione centrale chiusa, ricoprimento sufficiente e posizione centrale centrata a molla.
    \item la posizione di commutazione orientata alla sicurezza è raggiunta con l'annullamento del segnale di comando.
    \item il test può consistere, per esempio, nella verifica della caratteristica tempo/spazio (sistema di misura dello spostamento 1S1) dei movimenti pericolosi in relazione alla posizione di commutazione della valvola distributrice, con valutazione in un PLC (K1).
    \item K1 non deve essere utilizzato per il comando elettrico di 1V1.
    \item per prevenire un guasto sistematico, la funzione sovraordinata di messa fuori energia (che in questo esempio agisce sulla valvola di scarico 0V1) viene verificata a intervalli adeguati, per es. quotidianamente.
    \item per l'impiego in applicazioni con interventi poco frequenti dell'operatore nella zona pericolosa. Ciò consente di soddisfare il requisito dell'architettura designata per la Categoria 2. Il requisito prevede che il test venga eseguito immediatamente al momento della richiesta della funzione di sicurezza e che il tempo complessivo necessario per rilevare il guasto e portare la macchina in uno stato non pericoloso - tenuto conto per esempio della sovracorsa, che dipende tra l'altro dai tempi di depressurizzazione e di commutazione delle valvole (in questo caso la depressurizzazione avviene a livello superiore attraverso la valvola 0V1) - sia inferiore al tempo necessario per raggiungere il punto pericoloso (vedere anche la EN ISO 13855 e cfr. il punto 6.2.14).
    \item il componente standard K1 è impiegato conformemente alle indicazioni del punto 6.3.10.
    \item il software (SRASW) è programmato secondo i requisiti per il PL b (declassato in ragione della diversità) e le indicazioni del punto 6.3.
\end{itemize}

\textbf{Calcolo della probabilità di guasto}
\begin{itemize}
    \item MTTF\textsubscript{D} del canale funzionale: per la valvola distributrice 1V1 si assume un valore B\textsubscript{10D} di 20.000.000 di cicli di commutazione [S]. Con 240 giorni lavorativi, 16 ore lavorative al giorno e un tempo di ciclo di 5 secondi, n\textsubscript{op} è pari a 2.764.800 cicli di commutazione all'anno e l'MTTF\textsubscript{D} è di 72,3 anni. Questo è anche il valore dell'MTTF\textsubscript{D} del canale funzionale.
    \item MTTF\textsubscript{D} del canale di test: per il sistema di misura dello spostamento 1S1 si assume un valore MTTF\textsubscript{D} di 150 anni [E]. Per il PLC K1 si assume un valore MTTFD di 50 anni [E]. Per la valvola di scarico 0V1 vale un valore B10D di 20.000.000 di cicli [S]. Con un azionamento una volta al giorno su 240 giorni lavorativi, il valore MTTFD di 0V1 è di 833.333 anni. L'MTTFD del canale di test risulta quindi di 37,5 anni.
    \item DC\textsubscript{avg}: il DC del 60\% per 1V1 si basa sul confronto della caratteristica spazio/tempo del movimento pericoloso in relazione allo stato di commutazione della valvola distributrice. Questo è anche il DC\textsubscript{avg} ("basso").
    \item misure adeguate contro i guasti per causa comune (85 punti): separazione (15), diversità (20), protezione contro le sovratensioni ecc. (15) e condizioni ambientali (25 + 10)
    \item la combinazione degli elementi di comando pneumatici soddisfa la Categoria 2 con un MTTFD elevato (72,3 anni) e un DCavg basso (60\%). Ne risulta una probabilità media di guasto pericoloso di 7,6 10\textsuperscript{-7} all'ora. Ciò soddisfa il PL d. L'aggiunta di ulteriori SRP/CS sotto forma di sottosistemi per completare la funzione di sicurezza può, in determinate circostanze, comportare un PL inferiore. L'elemento soggetto a usura 1V1 dovrebbe essere sostituito all'incirca ogni sette anni (T\textsubscript{10D}).
\end{itemize}
\clearpage

\textbf{Esempio 12}
\textit{Valvola idraulica testata (sottosistema) - Categoria 2 - PL d}

% La minipage tiene insieme disegno e didascalia (figura non fluttuante)
\begin{center}
    \begin{minipage}{\textwidth}
        \centering
        \captionsetup{hypcap=false}
        \resizebox{0.83\textwidth}{!}{%==============================================================
%  Figura 8.23 - Valvola idraulica con test elettronico per il
%  comando di movimenti pericolosi (Esempio 12)
%  Adattamento in italiano della Figura 8.23 dell'IFA Report 2/2017
%
%  I simboli dei componenti riproducono quelli dell'originale
%  (ISO 1219): le coordinate sono ricavate dal disegno di partenza
%  (1 cm = 77,4 px, origine nell'angolo in alto a sinistra della
%  cornice); a questa scala i simboli hanno le stesse dimensioni
%  della Figura 8.7, da cui sono ripresi cilindro, valvole, pompa,
%  motore, filtri e strumenti del serbatoio. Nessun cambio di corpo
%  del font: le etichette sono ingrandite con "scale" (vedere la
%  nota nella Figura 8.3).
%==============================================================
\begin{tikzpicture}[
  font=\normalsize,
  linea/.style={draw=black, line width=0.8pt},
  simbolo/.style={draw=black, line width=0.8pt},
  sottile/.style={draw=black, line width=0.6pt},
  freccia/.style={sottile, -{Latex[length=9pt,width=4pt]}},
  segnale/.style={sottile, -{Latex[length=8pt,width=4pt]}},
  ingresso/.style={linea, -{Latex[length=5.5pt,width=3.6pt]}},
  pilota/.style={draw=black, line width=0.6pt, dash pattern=on 3.6pt off 2.4pt},
  elemento/.style={draw=black, line width=0.6pt, dash pattern=on 0.25cm off 0.121cm},
  etichetta/.style={inner sep=0pt, scale=1.15},
  testo/.style={inner sep=0pt, scale=0.9},
  piccolo/.style={inner sep=0pt, scale=0.85},
  grande/.style={inner sep=0pt, scale=1.65},
  nodo/.style={circle, fill=black, inner sep=0pt, minimum size=4.8pt},
  nodogrande/.style={circle, fill=black, inner sep=0pt, minimum size=8pt}
]

%--------------------------------------------------------------
% Cornice
%--------------------------------------------------------------
\draw[draw=black, line width=1.2pt] (0.000,-0.000) rectangle (17.804,-17.513);

%--------------------------------------------------------------
% Cilindro 1A: camicia, pistone e stelo a doppia linea, blocco di
% attacco con le due bocche
%--------------------------------------------------------------
\draw[simbolo] (4.561,-2.413) rectangle (6.914,-3.699);
\draw[simbolo] (4.828,-2.413) -- (4.828,-3.699);
\draw[simbolo] (5.352,-2.413) -- (5.352,-3.699);
\draw[simbolo, fill=white] (5.352,-2.927) rectangle (7.695,-3.189);
\draw[linea] (4.624,-3.699) -- (4.624,-4.627) -- (5.471,-4.627) -- (5.471,-6.303);
\draw[linea] (6.852,-3.699) -- (6.852,-4.627) -- (5.981,-4.627) -- (5.981,-6.303);
\node[etichetta, anchor=base] at (4.018,-2.713) {1A};

%--------------------------------------------------------------
% Sensore di posizione 1S3 sull'estremita' dello stelo, con uscita
% verso K1
%--------------------------------------------------------------
\draw[linea] (7.700,-2.933) -- (7.700,-2.106);
\draw[simbolo] (7.700,-1.059) rectangle (8.747,-2.106);
\draw[sottile] (7.713,-2.096) -- (8.737,-1.079);
\draw[sottile] (8.091,-1.983) -- (8.285,-1.983) -- (8.285,-1.854) -- (8.488,-1.854);
\node[grande, anchor=base] at (8.023,-1.486) {G};
\draw[segnale] (8.747,-1.059) -- (8.747,-0.362) -- (9.415,-0.740);
\node[etichetta, anchor=base] at (9.774,-0.685) {K1};
\node[etichetta, anchor=base] at (6.983,-1.357) {1S3};

%--------------------------------------------------------------
% Movimento pericoloso
%--------------------------------------------------------------
\draw[sottile, {Latex[length=6pt,width=3.4pt]}-{Latex[length=6pt,width=3.4pt]}] (9.160,-3.075) -- (12.804,-3.075);
\draw[sottile] (11.001,-2.946) -- (11.001,-3.204);
\node[testo, anchor=base] at (11.008,-2.248) {movimento};
\node[testo, anchor=base] at (11.008,-2.674) {pericoloso};

%--------------------------------------------------------------
% 1V3: valvola 4/3 a centro chiuso (P e T collegati), azionata da
% elettromagneti a e b (comandati da K1) e centrata da molle
%--------------------------------------------------------------
\draw[simbolo] (4.147,-6.059) rectangle (7.286,-7.106);
\draw[simbolo] (5.190,-6.059) -- (5.190,-7.106);
\draw[simbolo] (6.238,-6.059) -- (6.238,-7.106);
% elettromagnete a
\draw[simbolo] (2.571,-6.587) rectangle (4.147,-7.106);
\draw[simbolo] (3.357,-6.587) -- (3.357,-7.106);
\draw[sottile] (2.847,-7.106) -- (3.096,-6.587);
\fill[black] (3.357,-6.587) -- (3.633,-6.844) -- (3.357,-7.106) -- cycle;
% elettromagnete b
\draw[simbolo] (7.286,-6.587) rectangle (8.838,-7.106);
\draw[simbolo] (8.061,-6.587) -- (8.061,-7.106);
\draw[sottile] (8.314,-6.587) -- (8.586,-7.106);
\fill[black] (8.061,-6.587) -- (7.790,-6.844) -- (8.061,-7.106) -- cycle;
% posizione a: collegamenti incrociati
\draw[freccia] (4.404,-6.059) -- (4.928,-7.106);
\draw[freccia] (4.404,-7.106) -- (4.928,-6.059);
% posizione centrale: utilizzi chiusi, P collegato a T
\draw[sottile] (5.324,-6.320) -- (5.586,-6.320);
\draw[sottile] (5.833,-6.320) -- (6.110,-6.320);
\draw[sottile] (5.457,-7.106) -- (5.457,-6.587) -- (5.967,-6.587) -- (5.967,-7.106);
% posizione b: collegamenti paralleli
\draw[freccia] (6.490,-7.106) -- (6.490,-6.059);
\draw[freccia] (7.014,-6.059) -- (7.014,-7.106);
% molle di centraggio
\draw[sottile] (3.496,-6.324) -- (3.590,-6.576) -- (3.682,-6.059) -- (3.783,-6.576) -- (3.876,-6.059) -- (3.970,-6.576) -- (4.065,-6.059) -- (4.154,-6.305);
\draw[sottile] (7.293,-6.292) -- (7.391,-6.059) -- (7.486,-6.576) -- (7.578,-6.059) -- (7.672,-6.576) -- (7.765,-6.059) -- (7.862,-6.576) -- (7.953,-6.318);
\draw[segnale] (3.101,-6.576) -- (3.101,-6.030) -- (3.640,-5.727);
\draw[segnale] (8.342,-6.576) -- (8.342,-6.030) -- (8.880,-5.727);
\node[etichetta, anchor=base] at (2.061,-6.111) {1V3};
\node[etichetta, anchor=base] at (3.663,-5.568) {K1};
\node[etichetta, anchor=base] at (8.908,-5.568) {K1};
\node[etichetta, anchor=base] at (2.216,-6.951) {a};
\node[etichetta, anchor=base] at (9.309,-6.977) {b};

%--------------------------------------------------------------
% Linea di pressione P: valvola limitatrice 1V2, valvola di non
% ritorno 1V1, pompa 1P
%--------------------------------------------------------------
\draw[linea] (5.465,-7.106) -- (5.465,-12.791);
\draw[linea] (5.465,-13.127) -- (5.465,-14.186);
\draw[linea] (5.465,-15.749) -- (5.465,-16.796);
\node[nodogrande] at (5.465,-11.292) {};
\draw[linea] (5.465,-11.292) -- (7.041,-11.292);
% 1V2: valvola limitatrice di pressione con molla regolabile
\draw[simbolo] (7.041,-10.762) rectangle (8.088,-11.805);
\draw[sottile, -{Latex[length=8pt,width=3.6pt]}] (7.041,-11.548) -- (8.088,-11.548);
\draw[sottile] (7.248,-10.138) -- (7.041,-10.180) -- (7.470,-10.295) -- (7.041,-10.395)
  -- (7.470,-10.483) -- (7.041,-10.590) -- (7.470,-10.709) -- (7.327,-10.762);
\draw[sottile, -{Latex[length=8pt,width=4pt]}] (6.791,-10.562) -- (7.977,-10.288);
\draw[pilota] (7.041,-11.288) -- (6.524,-11.798) -- (6.524,-12.333) -- (7.310,-12.333)
  -- (7.310,-11.805);
\draw[linea] (8.088,-11.292) -- (8.863,-11.292) -- (8.863,-16.796);
\node[etichetta, anchor=base] at (6.208,-11.047) {1V2};
% 1V1: valvola di non ritorno con molla
\draw[sottile] (5.465,-12.270) -- (5.736,-12.347) -- (5.184,-12.431) -- (5.736,-12.493)
  -- (5.184,-12.577) -- (5.736,-12.646) -- (5.184,-12.731) -- (5.465,-12.789);
\draw[simbolo] (5.465,-12.921) circle (0.139);
\draw[simbolo] (5.264,-12.921) -- (5.465,-13.127) -- (5.651,-12.921);
\node[etichetta, anchor=base] at (4.399,-12.687) {1V1};
% 1P: pompa a cilindrata fissa con senso di rotazione
\draw[simbolo] (5.465,-14.939) circle (0.790);
\fill[black] (5.465,-14.153) -- (5.736,-14.606) -- (5.189,-14.606) -- cycle;
\draw[sottile, -{Latex[length=8pt,width=4pt]}] (5.465,-14.939) ++(-30:1.171) arc[start angle=-30, end angle=30, radius=1.171];
\node[etichetta, anchor=base] at (7.403,-15.078) {1P};
% 1M: motore elettrico trifase e giunto
\draw[simbolo] (2.081,-14.939) circle (0.786);
\draw[sottile] (2.856,-14.816) -- (3.642,-14.816);
\draw[sottile] (2.856,-15.077) -- (3.642,-15.077);
\draw[simbolo] (3.642,-14.677) -- (3.642,-15.201);
\draw[simbolo] (3.907,-14.677) -- (3.907,-15.201);
\draw[sottile] (3.907,-14.816) -- (4.682,-14.816);
\draw[sottile] (3.907,-15.077) -- (4.682,-15.077);
\node[etichetta, anchor=base] at (2.099,-14.819) {M};
\node[etichetta, anchor=base] at (2.054,-15.297) {3\,$\sim$};
\node[etichetta, anchor=base] at (0.717,-15.090) {1M};

%--------------------------------------------------------------
% Alimentazione del motore: linea L e contatto principale di Q1
%--------------------------------------------------------------
\node[etichetta, anchor=base] at (2.112,-10.762) {L};
\draw[linea] (2.083,-11.034) -- (2.083,-12.626);
\draw[sottile] (1.890,-11.592) -- (2.277,-11.399);
\draw[sottile] (1.890,-11.705) -- (2.277,-11.512);
\draw[sottile] (1.890,-11.828) -- (2.277,-11.625);
\draw[sottile] (2.083,-12.416) arc[start angle=90, end angle=270, x radius=0.081, y radius=0.105];
\draw[sottile] (1.809,-12.555) -- (2.083,-13.007);
\draw[linea] (2.083,-13.007) -- (2.083,-14.154);
\node[etichetta, anchor=base] at (1.214,-12.739) {Q1};

%--------------------------------------------------------------
% Linea di ritorno T: filtro 1Z2 con valvola di bypass e indicatore
% di intasamento
%--------------------------------------------------------------
\draw[linea] (5.982,-7.106) -- (5.982,-9.199) -- (11.072,-9.199) -- (11.072,-13.488);
\draw[linea] (11.072,-14.961) -- (11.072,-16.796);
\draw[simbolo] (9.224,-12.961) rectangle (12.139,-15.486);
\node[nodo] at (11.057,-13.180) {};
\node[nodo] at (11.057,-15.271) {};
\draw[simbolo] (11.057,-13.490) -- (11.804,-14.228) -- (11.057,-14.966) -- (10.319,-14.228) -- cycle;
\draw[elemento] (10.319,-14.228) -- (11.804,-14.228);
% bypass con valvola di non ritorno a molla
\draw[linea] (11.057,-13.180) -- (9.763,-13.180) -- (9.763,-13.814);
\draw[simbolo] (9.566,-13.931) -- (9.763,-13.744) -- (9.941,-13.931);
\draw[simbolo] (9.763,-13.948) circle (0.134);
\draw[sottile] (9.763,-14.056) -- (10.030,-14.154) -- (9.486,-14.247) -- (10.030,-14.341)
  -- (9.486,-14.437) -- (10.030,-14.530) -- (9.486,-14.618) -- (9.763,-14.708);
\draw[linea] (9.763,-14.083) -- (9.763,-15.271) -- (11.057,-15.271);
\draw[linea] (11.072,-13.178) -- (12.649,-13.178);
\draw[simbolo] (12.653,-12.791) rectangle (13.949,-13.577);
\draw[sottile, -{Latex[length=6pt,width=3.6pt]}] (13.244,-13.415) -- (12.767,-12.940);
\draw[sottile] (13.559,-13.187) circle (0.207);
\draw[sottile] (13.412,-13.040) -- (13.706,-13.334);
\draw[sottile] (13.412,-13.334) -- (13.706,-13.040);
\node[etichetta, anchor=base] at (9.528,-12.597) {1Z2};

%--------------------------------------------------------------
% PLC K1 con i segnali di ingresso da 1V3a, 1V3b, 1S3 e l'uscita
% verso la bobina del contattore Q1 del motore
%--------------------------------------------------------------
\draw[simbolo] (13.101,-6.189) rectangle (17.222,-8.540);
\draw[simbolo] (13.101,-6.770) -- (17.222,-6.770);
\draw[simbolo] (13.101,-7.959) -- (17.222,-7.959);
\node[piccolo, anchor=base] at (15.187,-6.550) {Ingressi};
\node[grande, anchor=base] at (15.213,-7.558) {PLC};
\node[piccolo, anchor=base] at (15.207,-8.346) {Uscite};
\node[etichetta, anchor=base] at (12.578,-6.499) {K1};
\draw[ingresso] (13.621,-4.952) -- (13.621,-6.189);
\draw[segnale] (13.621,-4.952) -- (12.963,-5.332);
\draw[ingresso] (15.194,-4.096) -- (15.194,-6.189);
\draw[segnale] (15.194,-4.096) -- (14.535,-4.479);
\draw[ingresso] (16.766,-3.815) -- (16.766,-6.189);
\draw[segnale] (16.766,-3.815) -- (16.236,-3.501);
\node[etichetta, anchor=base] at (12.099,-5.388) {1V3a};
\node[etichetta, anchor=base] at (13.786,-4.561) {1V3b};
\node[etichetta, anchor=base] at (16.053,-3.282) {1S3};
\draw[linea] (15.194,-8.540) -- (15.194,-9.716);
\draw[simbolo] (14.800,-9.716) rectangle (15.581,-10.116);
\draw[linea] (15.194,-10.116) -- (15.194,-10.891);
\draw[linea] (14.432,-10.891) -- (16.860,-10.891);
\node[nodo] at (15.194,-10.891) {};
\node[nodo] at (16.027,-10.891) {};
\draw[linea] (16.027,-10.891) -- (16.027,-11.499);
\draw[linea] (15.891,-11.499) -- (16.176,-11.499);
\node[etichetta, anchor=base] at (14.289,-10.039) {Q1};

%--------------------------------------------------------------
% Serbatoio con indicatore di livello 1S1, termometro 1S2 e
% filtro di aerazione 1Z1
%--------------------------------------------------------------
\draw[simbolo] (0.749,-15.956) -- (0.749,-17.080) -- (17.532,-17.080) -- (17.532,-15.956);
\draw[simbolo] (13.333,-15.588) circle (0.526);
\draw[sottile] (13.333,-15.188) -- (13.333,-15.922);
\node[nodo, minimum size=4.5pt] at (13.333,-15.922) {};
\draw[linea] (13.333,-16.112) -- (13.333,-16.802);
\node[etichetta, anchor=base] at (13.262,-14.767) {1S1};
\draw[simbolo] (14.910,-15.588) circle (0.529);
\draw[sottile] (14.381,-15.588) -- (15.438,-15.588);
\fill[black] (14.761,-15.351) -- (15.047,-15.351) -- (14.910,-15.588) -- cycle;
\draw[linea] (14.910,-16.116) -- (14.910,-16.802);
\node[etichetta, anchor=base] at (14.877,-14.755) {1S2};
\draw[simbolo] (16.745,-14.645) -- (17.484,-15.382) -- (16.745,-16.120) -- (15.998,-15.382) -- cycle;
\draw[elemento] (15.998,-15.382) -- (17.484,-15.382);
\draw[sottile] (16.617,-14.420) -- (16.874,-14.420) -- (16.745,-14.645) -- cycle;
\draw[sottile] (16.745,-14.420) -- (16.745,-14.159);
\draw[sottile] (16.617,-14.159) -- (16.874,-14.159) -- (16.745,-13.930) -- cycle;
\draw[linea] (16.745,-16.120) -- (16.745,-16.796);
\node[etichetta, anchor=base] at (15.730,-14.070) {1Z1};

\end{tikzpicture}
}
        \captionof{figure}{Valvola idraulica con test elettronico per il comando di movimenti pericolosi.}
        \label{fig:08-23}
    \end{minipage}
\end{center}

\textbf{Funzioni di sicurezza}
\begin{itemize}
    \item funzione di arresto legata alla sicurezza: arresto di un movimento pericoloso e prevenzione dell'avviamento inatteso dalla posizione di riposo.
    \item qui è rappresentata unicamente la parte idraulica del sistema di comando, sotto forma di sottosistema. Per completare la funzione di sicurezza devono essere aggiunti, sempre sotto forma di sottosistemi, ulteriori componenti di comando legati alla sicurezza (per es. dispositivi di protezione ed elementi logici elettrici).
\end{itemize}

\textbf{Descrizione funzionale}
\begin{itemize}
    \item i movimenti pericolosi sono comandati dalla valvola distributrice 1V3.
    \item un guasto della valvola distributrice 1V3 fra due test funzionali può comportare la perdita della funzione di sicurezza. La probabilità di guasto dipende dall'affidabilità della valvola distributrice.
    \item il test della funzione di sicurezza è forzato dal PLC K1 mediante un sistema di misura dello spostamento. L'interruzione del movimento pericoloso da parte della pompa idraulica comporta in genere una sovracorsa maggiore. La distanza dalla zona pericolosa deve essere scelta tenendo conto di tale maggiore sovracorsa.
    \item la funzione di test non deve essere compromessa da un guasto della valvola distributrice. Un guasto della funzione di test non deve provocare il guasto della valvola distributrice. Il test viene eseguito a intervalli adeguati e in occasione di una richiesta della funzione di sicurezza. Il rilevamento di un guasto di 1V3 determina la disinserzione della pompa idraulica 1M/1P per mezzo del contattore Q1.
\end{itemize}

\textbf{Caratteristiche di progetto}
\begin{itemize}
    \item sono rispettati i principi di sicurezza di base e i principi di sicurezza ben collaudati e sono soddisfatti i requisiti della Categoria B.
    \item 1V3 è una valvola distributrice con posizione centrale chiusa, ricoprimento sufficiente e posizione centrale centrata a molla.
    \item la posizione di commutazione orientata alla sicurezza è raggiunta con l'annullamento del segnale di comando.
    \item il test può consistere, per esempio, nella verifica della caratteristica spazio/tempo (sistema di misura dello spostamento 1S3) dei movimenti pericolosi in relazione alla posizione di commutazione della valvola distributrice, con valutazione in un PLC (K1). K1 non deve essere utilizzato per il comando elettrico di 1V3.
    \item per prevenire un guasto sistematico, la funzione sovraordinata di messa fuori energia (che in questo esempio agisce sulla pompa idraulica) viene verificata a intervalli adeguati, per es. quotidianamente.
    \item per l'impiego in applicazioni con interventi poco frequenti dell'operatore nella zona pericolosa. Ciò consente di soddisfare il requisito dell'architettura designata per la Categoria 2. Il requisito prevede che il test venga eseguito immediatamente al momento della richiesta della funzione di sicurezza e che il tempo complessivo necessario per rilevare il guasto e portare la macchina in uno stato non pericoloso - tenuto conto, per esempio, della sovracorsa - sia inferiore al tempo necessario per raggiungere il punto pericoloso (vedere anche la EN ISO 13855 e cfr. il punto 6.2.14)
    \item il componente standard K1 è impiegato conformemente alle indicazioni del punto 6.3.10.
    \item il software (SRASW) è programmato secondo i requisiti per il PL b (declassato in ragione della diversità) e le indicazioni del punto 6.3.
\end{itemize}

\textbf{Calcolo della probabilità di guasto}
\begin{itemize}
    \item MTTF\textsubscript{D} del canale funzionale: per la valvola distributrice 1V3 si assume un MTTF\textsubscript{D} di 150 anni [M]. Questo è anche il valore dell'MTTF\textsubscript{D} del canale funzionale, che viene dapprima limitato a 100 anni.
    \item MTTF\textsubscript{D} del canale di test: per il sistema di misura dello spostamento 1S3 si assume un valore MTTF\textsubscript{D} di 91 anni [M]. Per il PLC K1 si assume un valore MTTF\textsubscript{D} di 50 anni [E]. Per il contattore Q1 vale un valore B\textsubscript{10D} di 1.300.000 cicli [S]. Con un azionamento una volta al giorno su 240 giorni lavorativi, il valore MTTF\textsubscript{D} di Q1 è di 54.166 anni. L'MTTF\textsubscript{D} del canale di test risulta quindi di 32,3 anni. L'MTTF\textsubscript{D} del canale funzionale deve pertanto essere ridotto a 64,5 anni, conformemente al modello di analisi adottato.
    \item DC\textsubscript{avg}: il DC del 60\% per 1V3 si basa sul confronto della caratteristica spazio/tempo del movimento pericoloso in relazione allo stato di commutazione della valvola distributrice. Questo è anche il DC\textsubscript{avg} ("basso").
    \item misure adeguate contro i guasti per causa comune (85 punti): separazione (15), diversità (20), protezione contro le sovratensioni ecc. (15) e condizioni ambientali (25 + 10)
    \item la combinazione degli elementi di comando soddisfa la Categoria 2 con un MTTFD elevato (75 anni) e un DC\textsubscript{avg} basso (60\%). Ne risulta una probabilità media di guasto pericoloso di 8,7 10\textsuperscript{-7} all'ora. Ciò soddisfa il PL d. L'aggiunta di ulteriori SRP/CS sotto forma di sottosistemi per completare la funzione di sicurezza può, in determinate circostanze, comportare un PL inferiore.
\end{itemize}

